\documentclass[a4paper,12pt]{report}

% --------------------------------------------------
% Codifica e lingua
% --------------------------------------------------
\usepackage[T1]{fontenc}
\usepackage[utf8]{inputenc}
\usepackage[italian]{babel}

% --------------------------------------------------
% Impaginazione e tipografia
% --------------------------------------------------
\usepackage[a4paper,margin=2.5cm]{geometry}
\usepackage{microtype}
\usepackage{setspace}
\usepackage{parskip}

% --------------------------------------------------
% Matematica
% --------------------------------------------------
\usepackage{amsmath,amssymb,mathtools}
\usepackage{amsthm}
\usepackage{bm}
\usepackage{physics}

% --------------------------------------------------
% Figure, grafici, colori
% --------------------------------------------------
\usepackage{graphicx}
\usepackage{xcolor}
\usepackage{float}
\usepackage{caption}
\usepackage{subcaption}
\usepackage{booktabs}
\usepackage{array}
\usepackage{multirow}

% --------------------------------------------------
% TikZ e PGFPlots per diagrammi, blocchi, grafici
% --------------------------------------------------
\usepackage{tikz}
\usepackage{pgfplots}
\pgfplotsset{compat=1.18}

\usetikzlibrary{
    arrows.meta,
    calc,
    positioning,
    fit,
    matrix,
    backgrounds,
    decorations.pathreplacing,
    shapes.geometric,
    shapes.symbols,
    shapes.misc,
    patterns
}

% Se vuoi anche circuiti/diagrammi più tecnici
\usepackage{circuitikz}

% --------------------------------------------------
% Link ipertestuali
% --------------------------------------------------
\usepackage[
    colorlinks=true,
    linkcolor=blue!50!black,
    urlcolor=blue!50!black,
    citecolor=blue!50!black
]{hyperref}

% --------------------------------------------------
% Comandi utili
% --------------------------------------------------
\newcommand{\R}{\mathbb{R}}
\newcommand{\C}{\mathbb{C}}

\begin{document}

% --------------------------------------------------
% Frontespizio
% --------------------------------------------------
\begin{titlepage}
    \centering
    \vspace*{2cm}

    {\Huge\bfseries Teoria del Controllo\par}
    \vspace{0.8cm}

    {\Large Professor Caiti\par}
    \vspace{0.5cm}

    {\large Università di Pisa\par}
    \vspace{0.3cm}

    {\large Ingegneria Robotica e dell'Automazione\par}

    \vfill

    {\Large\bfseries Autori: Ingegneri Robotizzati\par}
    \vspace{0.8cm}

    {\large\textbf{Materiale:}\par}
    \vspace{0.2cm}
    {\large Alessio, Francesca, Mery, Gabriele, Samuele, Lucia, Donatella\par}

    \vfill

\end{titlepage}

% --------------------------------------------------
% Disclaimer
% --------------------------------------------------
\chapter*{Disclaimer}
\addcontentsline{toc}{chapter}{Disclaimer}

Questo documento è stato realizzato con il supporto dell'intelligenza artificiale a partire dagli appunti presi in aula durante le lezioni. Il materiale è stato successivamente rielaborato, organizzato e ampliato mediante l'uso di libri, documenti di supporto e strumenti di IA, con l'obiettivo di renderlo più chiaro, completo e utile allo studio.

Il testo ha esclusivamente finalità didattiche e di supporto personale allo studio. Non sostituisce in alcun modo le lezioni, il materiale ufficiale del corso, né i testi consigliati dal docente.

Questo documento non è assolutamente pensato per la vendita, la distribuzione commerciale o qualunque altra forma di utilizzo a scopo di lucro.

\tableofcontents

\chapter{Lezione 1 --- Sistemi in retroazione}

\section{Schema generale del sistema in retroazione}

Consideriamo un sistema lineare tempo-invariante descritto nel dominio di Laplace
da un regolatore \(R(s)\) e da un impianto \(G(s)\), con retroazione unitaria negativa.

Indicando con:
\begin{itemize}
    \item \(y^\star(t)\) il riferimento,
    \item \(e(t)\) l'errore,
    \item \(u(t)\) il segnale di controllo,
    \item \(y(t)\) l'uscita,
\end{itemize}
si ha nel dominio di Laplace
\[
E(s)=Y^\star(s)-Y(s),
\qquad
U(s)=R(s)E(s),
\qquad
Y(s)=G(s)U(s).
\]

La funzione d'anello, o funzione di catena diretta, è
\[
L(s)=R(s)G(s).
\]

La funzione di trasferimento in ciclo chiuso dal riferimento all'uscita è
\[
H(s)=\frac{Y(s)}{Y^\star(s)}=\frac{L(s)}{1+L(s)}=\frac{R(s)G(s)}{1+R(s)G(s)}.
\]

Nel seguito si assume che le eventuali dinamiche non visibili dall'esterno
(non osservabili o non raggiungibili) siano stabili, in modo da poter
ricondurre l'analisi della stabilità del sistema alla stabilità del ciclo chiuso.

\section{Obiettivo del progetto del regolatore}

Dato l'impianto \(G(s)\), il problema del controllo consiste nel progettare \(R(s)\)
in modo che il sistema in retroazione soddisfi alcune specifiche.

Le richieste principali si dividono tipicamente in tre gruppi:
\begin{enumerate}
    \item stabilità;
    \item prestazioni statiche, cioè a regime;
    \item prestazioni dinamiche, cioè nel transitorio.
\end{enumerate}

\section{Stabilità}

Una prima richiesta fondamentale è la stabilità BIBO
(\emph{Bounded Input, Bounded Output}):
a ogni ingresso limitato deve corrispondere un'uscita limitata.

In formule:
\[
y^\star \in L_\infty \quad \Longrightarrow \quad y \in L_\infty.
\]

La stabilità BIBO è la proprietà minima richiesta al sistema in ciclo chiuso.
In molti problemi si richiede anche la stabilità asintotica, cioè che i transitori
si estinguano nel tempo.

\section{Prestazioni statiche}

Le prestazioni statiche riguardano l'errore a regime, cioè la differenza tra
riferimento e uscita quando il transitorio si è esaurito.

Definiamo
\[
e_\infty=\lim_{t\to+\infty}\bigl(y^\star(t)-y(t)\bigr),
\]
quando tale limite esiste.

L'errore a regime dipende dal comportamento di \(L(s)\) alle basse frequenze
e, in particolare, dal numero di poli nell'origine presenti nella funzione d'anello.

\subsection{Tipo del sistema}

Si dice che il sistema è di tipo \(\nu\) se la funzione d'anello \(L(s)\)
ha \(\nu\) poli nell'origine, cioè se si può scrivere localmente per \(s\to 0\) come
\[
L(s)\sim \frac{K}{s^\nu},
\qquad \nu=0,1,2,\dots
\]
con \(K>0\).

Il numero di integratori nella catena diretta è quindi ciò che determina
la capacità del sistema di inseguire segnali di riferimento ``lenti''.

\subsection{Ingressi tipici}

I riferimenti canonici sono:
\[
\text{gradino: } r(t)=1(t),
\qquad
\text{rampa: } r(t)=t\,1(t),
\qquad
\text{parabola: } r(t)=\frac{t^2}{2}\,1(t),
\]
le cui trasformate di Laplace sono rispettivamente
\[
R(s)=\frac{1}{s},
\qquad
R(s)=\frac{1}{s^2},
\qquad
R(s)=\frac{1}{s^3}.
\]

\subsection{Errore a regime per gradino, rampa e parabola}

Per sistemi in retroazione unitaria, le regole classiche sono:

\[
\begin{array}{c|ccc}
 & \nu=0 & \nu=1 & \nu=2 \\ \hline
\text{gradino} & \dfrac{1}{1+K_p} & 0 & 0 \\[1.2ex]
\text{rampa} & +\infty & \dfrac{1}{K_v} & 0 \\[1.2ex]
\text{parabola} & +\infty & +\infty & \dfrac{1}{K_a}
\end{array}
\]
dove:
\[
K_p=\lim_{s\to 0}L(s),\qquad
K_v=\lim_{s\to 0}sL(s),\qquad
K_a=\lim_{s\to 0}s^2L(s).
\]

Questa tabella mostra che:
\begin{itemize}
    \item un alto guadagno a bassa frequenza riduce l'errore statico;
    \item la presenza di integratori è essenziale per annullare l'errore
    su riferimenti di tipo rampa o parabola.
\end{itemize}

\subsection{Interpretazione in frequenza}

Le prestazioni statiche si leggono dal comportamento di \(L(j\omega)\) per \(\omega\to 0\).

Se
\[
|L(j\omega)|\gg 1 \qquad \text{per } \omega\to 0,
\]
allora
\[
H(j\omega)=\frac{L(j\omega)}{1+L(j\omega)}\approx 1,
\]
cioè il sistema in ciclo chiuso riproduce bene il riferimento alle basse frequenze.

Questa è la ragione per cui si desidera che la funzione d'anello abbia
modulo elevato in bassa frequenza.

\section{Risposta armonica}

Se il sistema è lineare e asintoticamente stabile, a un ingresso sinusoidale
\[
u(t)=A\sin(\omega t)
\]
corrisponde, a regime, un'uscita sinusoidale alla stessa pulsazione:
\[
y(t)=A\,|H(j\omega)|\,\sin\bigl(\omega t+\varphi(\omega)\bigr),
\]
dove
\[
\varphi(\omega)=\arg H(j\omega).
\]

In forma complessa:
\[
H(j\omega)\in\mathbb{C},
\qquad
H(j\omega)=|H(j\omega)|\,e^{j\varphi(\omega)}.
\]

Il modulo \(|H(j\omega)|\) dice quanto il sistema amplifica o attenua
una sinusoide a quella frequenza; la fase \(\varphi(\omega)\) dice quanto la ritarda o anticipa.

\section{Prestazioni dinamiche}

Le prestazioni dinamiche riguardano il transitorio, cioè il modo in cui il sistema
raggiunge il regime.

\subsection{Sistema del primo ordine}

Un modello elementare è
\[
H(s)=\frac{1}{1+s\tau}.
\]

La risposta al gradino unitario è
\[
y(t)=1-e^{-t/\tau}.
\]

Il parametro \(\tau\) è la costante di tempo:
più \(\tau\) è piccola, più il sistema è rapido.

\subsection{Sistema del secondo ordine}

Un modello standard del secondo ordine è
\[
H(s)=\frac{\omega_n^2}{s^2+2\zeta\omega_n s+\omega_n^2}
=
\frac{1}{1+2\zeta \dfrac{s}{\omega_n}+\dfrac{s^2}{\omega_n^2}}.
\]

I parametri principali sono:
\begin{itemize}
    \item \(\omega_n\): pulsazione naturale;
    \item \(\zeta\): coefficiente di smorzamento.
\end{itemize}

Se \(0<\zeta<1\), la risposta al gradino è oscillante smorzata.
All'aumentare di \(\zeta\), le oscillazioni si riducono.

\subsection{Margine di fase}

Nel progetto in frequenza è molto importante il margine di fase \(M_\varphi\).

Sia \(\omega_c\) la pulsazione di attraversamento in ampiezza, definita da
\[
|L(j\omega_c)|=1
\qquad
\text{(oppure 0 dB)}.
\]

Il margine di fase è
\[
M_\varphi = \pi + \arg L(j\omega_c)
\]
oppure, in gradi,
\[
M_\varphi = 180^\circ + \arg L(j\omega_c).
\]

Interpretazione:
\begin{itemize}
    \item se il margine di fase è troppo piccolo, il sistema è poco smorzato
    o addirittura instabile;
    \item un margine di fase intorno a \(45^\circ\)-\(70^\circ\) è spesso
    associato a un comportamento soddisfacente.
\end{itemize}

In modo euristico, per sistemi dominati dal secondo ordine, si usa spesso la regola:
\[
\zeta \approx \frac{M_\varphi}{100}
\qquad
\text{(con \(M_\varphi\) espresso in gradi, stima grossolana).}
\]

\section{Bassa e alta frequenza}

Il progetto del controllore deve tenere conto sia del comportamento a bassa frequenza
sia di quello ad alta frequenza.

\subsection{Bassa frequenza}

Alle basse frequenze si vuole:
\[
|L(j\omega)| \gg 1,
\]
in modo da avere
\[
H(j\omega)\approx 1.
\]

Questo garantisce buone prestazioni statiche e un buon inseguimento del riferimento.

\subsection{Alta frequenza}

Alle alte frequenze, se
\[
|L(j\omega)|\ll 1,
\]
allora
\[
H(j\omega)\approx L(j\omega).
\]

In genere è desiderabile attenuare le alte frequenze, perché lì si collocano
rumore di misura, componenti non modellate e dinamiche indesiderate.

\section{Disturbi, sensibilità e sensibilità complementare}

Introduciamo ora due disturbi tipici:
\begin{itemize}
    \item \(m\): disturbo di misura, cioè rumore sul segnale misurato;
    \item \(d\): disturbo di carico, applicato in uscita all'impianto.
\end{itemize}

Assumiamo il modello standard:
\[
e(t)=y^\star(t)-\bigl(y(t)+m(t)\bigr),
\]
e che il disturbo \(d(t)\) si sommi all'uscita dell'impianto.

Nel dominio di Laplace si ottiene
\[
Y(s)=T(s)\,Y^\star(s)-T(s)\,M(s)+S(s)\,D(s),
\]
dove:
\[
T(s)=\frac{R(s)G(s)}{1+R(s)G(s)}
\]
è la \emph{sensibilità complementare}, mentre
\[
S(s)=\frac{1}{1+R(s)G(s)}
\]
è la \emph{sensibilità}.

Le due funzioni soddisfano la relazione fondamentale
\[
S(s)+T(s)=1.
\]

\subsection{Significato fisico}

\begin{itemize}
    \item \(T(s)\) governa l'effetto del riferimento e del rumore di misura sull'uscita.
    \item \(S(s)\) governa l'effetto dei disturbi di carico e l'errore residuo.
\end{itemize}

Di conseguenza:
\begin{itemize}
    \item a bassa frequenza si desidera \(S(j\omega)\approx 0\), quindi \(T(j\omega)\approx 1\),
    per rigettare i disturbi lenti e seguire bene il riferimento;
    \item ad alta frequenza si desidera \(T(j\omega)\approx 0\), quindi \(S(j\omega)\approx 1\),
    per non amplificare il rumore di misura.
\end{itemize}

Questa è la classica tensione di progetto del controllo in frequenza:
non si può rendere contemporaneamente piccolo sia \(S\) sia \(T\) alla stessa frequenza,
perché la loro somma è sempre uguale a \(1\).

\section{Filtri ideali}

Gli andamenti qualitativi richiamati nella lezione sono quelli dei filtri ideali.

\subsection{Passa-basso}

Un filtro passa-basso lascia passare le basse frequenze e attenua le alte:
\[
|H(j\omega)| \approx
\begin{cases}
1, & \omega \ll \omega_c,\\
0, & \omega \gg \omega_c.
\end{cases}
\]

\subsection{Passa-alto}

Un filtro passa-alto attenua le basse frequenze e lascia passare le alte:
\[
|H(j\omega)| \approx
\begin{cases}
0, & \omega \ll \omega_c,\\
1, & \omega \gg \omega_c.
\end{cases}
\]

\subsection{Passa-banda}

Un filtro passa-banda lascia passare un intervallo di frequenze compreso tra
\(\omega_1\) e \(\omega_2\), attenuando quelle molto basse e quelle molto alte.

\section{Sintesi della lezione}

La progettazione del regolatore \(R(s)\) si basa quindi su pochi principi chiave:
\begin{enumerate}
    \item rendere stabile il sistema in ciclo chiuso;
    \item garantire prestazioni statiche adeguate attraverso alto guadagno a bassa frequenza
    e opportuno numero di integratori;
    \item garantire prestazioni dinamiche adeguate controllando banda e margine di fase;
    \item attenuare i disturbi mediante le funzioni di sensibilità
    \[
    S(s)=\frac{1}{1+R(s)G(s)},\qquad
    T(s)=\frac{R(s)G(s)}{1+R(s)G(s)}.
    \]
\end{enumerate}
\chapter{Lezione 2 --- Spazi di funzioni, serie di Fourier, trasformate e distribuzioni}

\section{Spazio di funzioni periodiche}

Consideriamo uno spazio di funzioni periodiche di periodo \(T\).
L'idea è lavorare in uno spazio funzionale su cui sia possibile introdurre
una struttura geometrica analoga a quella degli spazi vettoriali.

Indichiamo con
\[
L_T^2
\]
lo spazio delle funzioni periodiche di periodo \(T\) che siano quadrato-sommabili
su un periodo, cioè tali che
\[
\int_{t_0}^{t_0+T} |f(t)|^2\,dt < +\infty.
\]

Poiché \(f\) è periodica, il valore di questo integrale non dipende dalla scelta di \(t_0\).

\section{Richiamo su integrabilità}

Negli appunti si distinguono due nozioni di integrale.

\subsection{Integrale di Riemann}

Nell'integrale di Riemann si approssima l'area sotto il grafico della funzione
discretizzando il dominio e costruendo somme di rettangoli.

\subsection{Integrale di Lebesgue}

Nell'integrale di Lebesgue la costruzione è più generale:
invece di discretizzare il dominio, si discretizza il codominio,
cioè si ragiona sui livelli assunti dalla funzione.

L'integrale di Lebesgue consente di trattare classi più ampie di funzioni
e, sotto ipotesi opportune, permette di scambiare limiti e integrali.

Nel seguito assumiamo che l'integrabilità sia nel senso di Lebesgue.

\section{Perché funzioni a quadrato sommabile}

Vogliamo introdurre nello spazio delle funzioni una struttura geometrica,
in particolare prodotto scalare e norma.

Per questo consideriamo funzioni
\[
f:\mathbb{R}\to\mathbb{R}
\quad \text{oppure} \quad
f:\mathbb{R}\to\mathbb{C}
\]
tali che
\[
\int_{t_0}^{t_0+T} |f(t)|^2\,dt < +\infty.
\]

Questa ipotesi garantisce che la norma della funzione sia finita.

\section{Prodotto scalare in \(L_T^2\)}

Se \(f_1,f_2\in L_T^2\), definiamo il prodotto scalare come
\[
\langle f_1,f_2\rangle
=
\frac{1}{T}\int_{t_0}^{t_0+T} f_1(t)\,\overline{f_2(t)}\,dt.
\]

Se le funzioni sono reali, la coniugazione si può omettere:
\[
\langle f_1,f_2\rangle
=
\frac{1}{T}\int_{t_0}^{t_0+T} f_1(t)\,f_2(t)\,dt.
\]

Il fattore \(1/T\) serve a normalizzare il prodotto scalare rispetto alla lunghezza del periodo.

\section{Norma}

Una volta definito il prodotto scalare, la norma associata è
\[
\|f\|=\sqrt{\langle f,f\rangle}.
\]

Questo chiarisce perché si richieda che le funzioni appartengano a \(L_T^2\):
senza quadrato-sommabilità la norma potrebbe non esistere oppure essere infinita.

\section{Interpretazione geometrica}

Con prodotto scalare e norma, lo spazio delle funzioni acquista una struttura
simile a quella di uno spazio vettoriale euclideo.

Di conseguenza si può parlare di:
\begin{itemize}
    \item basi;
    \item ortogonalità;
    \item ortonormalità;
    \item proiezioni.
\end{itemize}

L'idea fondamentale di Fourier è che una funzione periodica possa essere espressa
come combinazione lineare di sinusoidi a frequenze diverse.

\section{Formula di Eulero}

Ricordiamo la formula di Eulero:
\[
e^{j\omega t}=\cos(\omega t)+j\sin(\omega t).
\]

Questa scrittura permette di trattare seni e coseni con un'unica notazione
tramite esponenziali complessi.

\section{Serie di Fourier complessa}

Sia \(f\in L_T^2\). Allora, in forma complessa, la sua serie di Fourier è
\[
f(t)=\sum_{n=-\infty}^{+\infty} C_n e^{j n\omega_0 t},
\qquad
\omega_0=\frac{2\pi}{T}.
\]

Qui:
\begin{itemize}
    \item \(T\) è il periodo della funzione;
    \item \(\omega_0\) è la pulsazione fondamentale;
    \item la frequenza fondamentale in Hertz è
    \[
    \nu_0=\frac{1}{T}=\frac{\omega_0}{2\pi}.
    \]
\end{itemize}

La formula dice che ogni funzione del nostro spazio si rappresenta come combinazione lineare
di armoniche della fondamentale.

\section{Che funzioni base compaiono}

Le funzioni base sono
\[
e^{j n\omega_0 t},
\qquad n\in\mathbb{Z}.
\]

Equivalentemente si può parlare di seni e coseni alle pulsazioni
\[
n\omega_0.
\]

Queste sono precisamente le sinusoidi che hanno periodo \(T\) oppure un suo sottomultiplo:
\[
\frac{T}{|n|}, \qquad n\neq 0.
\]

La famiglia delle funzioni base è infinita ma numerabile,
perché è indicizzata dagli interi.

\section{Ortogonalità della base}

Nel prodotto scalare definito sopra, le funzioni esponenziali complesse risultano ortogonali:
\[
\left\langle e^{j m\omega_0 t}, e^{j n\omega_0 t}\right\rangle
=
\frac{1}{T}\int_{t_0}^{t_0+T}
e^{j m\omega_0 t}\,\overline{e^{j n\omega_0 t}}\,dt.
\]

Poiché
\[
\overline{e^{j n\omega_0 t}}=e^{-j n\omega_0 t},
\]
si ottiene
\[
\left\langle e^{j m\omega_0 t}, e^{j n\omega_0 t}\right\rangle
=
\frac{1}{T}\int_{t_0}^{t_0+T} e^{j(m-n)\omega_0 t}\,dt
=
\begin{cases}
1, & m=n,\\
0, & m\neq n.
\end{cases}
\]

Quindi la famiglia
\[
\left\{e^{j n\omega_0 t}\right\}_{n\in\mathbb{Z}}
\]
è ortonormale.

\section{Coefficienti della serie di Fourier}

I coefficienti \(C_n\) si ottengono proiettando \(f\) sulla funzione base corrispondente:
\[
C_n=\left\langle f(t), e^{j n\omega_0 t}\right\rangle.
\]

Usando la definizione di prodotto scalare si ricava
\[
C_n
=
\frac{1}{T}\int_{t_0}^{t_0+T} f(t)e^{-j n\omega_0 t}\,dt.
\]

Scegliendo per comodità \(t_0=0\), si ha
\[
C_n
=
\frac{1}{T}\int_0^T f(t)e^{-j n\omega_0 t}\,dt.
\]

\section{Significato dei coefficienti \(C_n\)}

In generale \(C_n\) è un numero complesso.

Possiamo quindi scrivere
\[
C_n = |C_n| e^{j\varphi_n}.
\]

Il modulo \(|C_n|\) misura il peso dell'armonica \(n\omega_0\),
mentre la fase \(\varphi_n\) ne misura lo sfasamento.

Se una funzione coincide con una singola armonica, allora un solo coefficiente è non nullo;
negli altri casi il segnale è una combinazione lineare di più armoniche.

\section{Interpretazione per i sistemi lineari}

Se un segnale periodico entra in un sistema lineare, possiamo scomporlo in armoniche.
Per linearità, la risposta complessiva si ottiene sommando le risposte del sistema
alle singole componenti sinusoidali.

Questa è la base concettuale dell'analisi in frequenza dei sistemi lineari.

\section{Dal caso periodico al caso non periodico}

Se il periodo \(T\) cresce, la pulsazione fondamentale
\[
\omega_0=\frac{2\pi}{T}
\]
diventa sempre più piccola.

Nel limite \(T\to+\infty\):
\begin{itemize}
    \item le frequenze discrete \(n\omega_0\) si addensano;
    \item la sommatoria sulle armoniche tende a un integrale;
    \item si passa da una descrizione discreta in frequenza a una descrizione continua.
\end{itemize}

Questo passaggio porta dalla serie di Fourier alla trasformata di Fourier.

\section{Trasformata di Fourier}

Nel caso non periodico si scrive formalmente
\[
f(t)=\frac{1}{2\pi}\int_{-\infty}^{+\infty} F(\omega)e^{j\omega t}\,d\omega,
\]
dove
\[
\mathcal{F}\{f(t)\}=F(\omega)
=
\int_{-\infty}^{+\infty} f(t)e^{-j\omega t}\,dt.
\]

Qui \(\omega\) è una variabile continua e \(F(\omega)\) è, in generale, una funzione complessa.

Possiamo scrivere
\[
F(\omega)=|F(\omega)|e^{j\arg F(\omega)}.
\]

Il modulo \(|F(\omega)|\) descrive il contenuto in frequenza del segnale;
la fase \(\arg F(\omega)\) descrive lo sfasamento delle varie componenti.

\section{Segnali stazionari e significato della Fourier}

La trasformata di Fourier è lo strumento naturale per analizzare segnali stazionari
oppure segnali definiti su tutto l'asse reale.

Nel caso di segnali periodici o stazionari, il contenuto in frequenza non dipende
dall'istante da cui si comincia a osservare il segnale.

Per questo la Fourier è particolarmente adatta a descrivere il regime
e il contenuto armonico del segnale.

\section{Trasformata di Laplace}

La trasformata di Laplace unilaterale si definisce come
\[
\mathcal{L}\{f(t)\}=F(s)
=
\int_{0^-}^{+\infty} f(t)e^{-st}\,dt,
\]
dove
\[
s\in\mathbb{C},
\qquad
s=\sigma+j\omega.
\]

La variabile complessa \(s\) permette di pesare il segnale non solo con sinusoidi,
ma con esponenziali del tipo \(e^{st}\).

\section{Significato della Laplace}

La trasformata di Laplace è particolarmente utile per:
\begin{itemize}
    \item segnali causali;
    \item fenomeni transitori;
    \item studio dei sistemi dinamici;
    \item trattamento delle condizioni iniziali;
    \item descrizione di segnali che non ammettono Fourier classica.
\end{itemize}

Per questo, dal punto di vista fisico e ingegneristico:
\begin{itemize}
    \item Fourier è soprattutto uno strumento per il contenuto in frequenza
    e per i segnali stazionari;
    \item Laplace è soprattutto uno strumento per sistemi causali, transitori
    e risposte forzate.
\end{itemize}

\section{Rapporto tra Fourier e Laplace}

Formalmente, la trasformata di Fourier si può vedere come la trasformata di Laplace
valutata sull'asse immaginario:
\[
s=j\omega.
\]

Tuttavia questo passaggio è lecito solo quando l'asse immaginario appartiene
alla regione di convergenza della trasformata di Laplace.

Quindi Fourier non è semplicemente ``Laplace con \(s=j\omega\)''
in modo automatico: il legame è corretto, ma richiede condizioni di esistenza.

\section{Distribuzioni}

Per trattare in modo rigoroso oggetti come impulso, gradino e loro derivate,
si introducono le distribuzioni.

Una distribuzione non è, in generale, una funzione regolare nel senso classico,
ma un oggetto definito attraverso il suo effetto integrale su funzioni test.

\section{Delta di Dirac}

L'esempio fondamentale è la delta di Dirac \(\delta(t)\), definita dalla proprietà
\[
\int_{-\infty}^{+\infty} f(t)\delta(t)\,dt = f(0).
\]

Questa relazione dice che la delta estrae il valore della funzione nel punto \(t=0\).

Intuitivamente:
\begin{itemize}
    \item la delta è concentrata in \(t=0\);
    \item il suo effetto è nullo ovunque tranne che in quell'istante;
    \item ha senso solo sotto segno di integrale.
\end{itemize}

\section{Altre distribuzioni utili}

Anche il gradino, la rampa e le derivate della delta
si interpretano correttamente nel quadro delle distribuzioni.

Questo è importante perché trasformate come Laplace e Fourier,
essendo operazioni integrali, possono essere estese anche a questi oggetti generalizzati.

\section{Sintesi della lezione}

I punti chiave della lezione sono:
\begin{enumerate}
    \item introduzione dello spazio \(L_T^2\) delle funzioni periodiche quadrato-sommabili;
    \item definizione di prodotto scalare e norma;
    \item interpretazione geometrica dello spazio delle funzioni;
    \item base ortonormale di esponenziali complessi;
    \item serie di Fourier e coefficienti \(C_n\);
    \item passaggio dal caso periodico al caso non periodico;
    \item trasformata di Fourier;
    \item trasformata di Laplace;
    \item distribuzioni e delta di Dirac.
\end{enumerate}
\chapter{Lezione 3 --- Distribuzioni, regolatore digitale, campionamento e teorema di Shannon}

\section{Richiami dalla lezione precedente}

Riprendiamo le tre rappresentazioni introdotte nella lezione precedente:
\begin{itemize}
    \item \textbf{Serie di Fourier} per segnali periodici di periodo \(T\):
    \[
    f(t)=\sum_{n=-\infty}^{+\infty} C_n e^{j n\omega_0 t},
    \qquad
    \omega_0=\frac{2\pi}{T},
    \]
    con coefficienti
    \[
    C_n=\frac{1}{T}\int_{-T/2}^{T/2} f(t)e^{-j n\omega_0 t}\,dt.
    \]

    \item \textbf{Trasformata di Fourier}:
    \[
    \mathcal{F}\{f(t)\}=F(\omega)=\int_{-\infty}^{+\infty} f(t)e^{-j\omega t}\,dt.
    \]

    \item \textbf{Trasformata di Laplace unilaterale}:
    \[
    \mathcal{L}\{f(t)\}=F(s)=\int_{0^-}^{+\infty} f(t)e^{-st}\,dt.
    \]
\end{itemize}

Questi tre strumenti verranno ora usati per introdurre alcune distribuzioni fondamentali
e per arrivare al problema del campionamento.

\section{Distribuzioni}

Una distribuzione non viene definita come una funzione ``ordinaria'', ma attraverso
le sue proprietà integrali quando viene moltiplicata per una funzione regolare.

In modo intuitivo, se \(g\) è una distribuzione e \(f\) è una funzione regolare,
l'oggetto significativo è
\[
\int_{-\infty}^{+\infty} f(t)\,g(t)\,dt.
\]

Questo numero reale o complesso rappresenta l'azione della distribuzione \(g\)
sulla funzione test \(f\).

\subsection{Distribuzioni fondamentali richiamate a lezione}

Le distribuzioni citate negli appunti sono:
\[
\delta(t),\qquad \delta'(t),\qquad \delta''(t),
\]
e inoltre il gradino unitario
\[
1(t)\equiv H(t),
\]
detto anche \textbf{funzione gradino} o \textbf{funzione di Heaviside}.

\section{Delta di Dirac}

La distribuzione fondamentale è la \textbf{delta di Dirac}, definita dalla proprietà
\[
\int_{-\infty}^{+\infty} f(t)\,\delta(t)\,dt = f(0).
\]

Poiché la delta è concentrata nell'origine, si può anche scrivere
\[
\int_{0^-}^{+\infty} f(t)\,\delta(t)\,dt = f(0).
\]

Dal punto di vista intuitivo:
\begin{itemize}
    \item \(\delta(t)\) è nulla per \(t\neq 0\);
    \item la sua area totale vale \(1\);
    \item il suo effetto ha senso solo sotto segno di integrale.
\end{itemize}

\section{Derivata della delta di Dirac}

Negli appunti compare la derivata della delta. Con la convenzione standard delle distribuzioni,
vale
\[
\int_{-\infty}^{+\infty} f(t)\,\delta'(t)\,dt = -f'(0).
\]

Più in generale, per la derivata \(n\)-esima della delta,
\[
\int_{-\infty}^{+\infty} f(t)\,\delta^{(n)}(t)\,dt
=
(-1)^n f^{(n)}(0).
\]

Questa relazione si ottiene per integrazione per parti in senso distribuzionale.

\section{Gradino di Heaviside}

Il gradino unitario si definisce come
\[
H(t)=1(t)=
\begin{cases}
0, & t<0,\\
1, & t>0.
\end{cases}
\]

Dal punto di vista distribuzionale, il gradino è legato alla delta tramite
\[
H(t)=\int_{-\infty}^{t}\delta(\tau)\,d\tau,
\]
e quindi
\[
\frac{d}{dt}H(t)=\delta(t).
\]

Questo è il collegamento fondamentale tra impulso e gradino.

\section{Tabella elementare di trasformate}

Raccogliamo i segnali richiamati negli appunti e le relative trasformate.

\subsection*{Trasformate di Laplace}

\[
\mathcal{L}\{\delta(t)\}=1,
\]
\[
\mathcal{L}\{H(t)\}=\frac{1}{s},
\]
\[
\mathcal{L}\{tH(t)\}=\frac{1}{s^2},
\]
\[
\mathcal{L}\{e^{-\sigma t}H(t)\}=\frac{1}{s+\sigma},
\qquad \Re(s)>-\sigma.
\]

\subsection*{Trasformate di Fourier}

Nel contesto ingegneristico degli appunti si usano spesso le forme sintetiche:
\[
\mathcal{F}\{\delta(t)\}=1,
\]
\[
\mathcal{F}\{H(t)\}\sim \frac{1}{j\omega},
\]
\[
\mathcal{F}\{tH(t)\}\sim -\frac{1}{\omega^2},
\]
\[
\mathcal{F}\{e^{-\sigma t}H(t)\}=\frac{1}{\sigma+j\omega},
\qquad \sigma>0.
\]

\medskip

\noindent
\textbf{Osservazione.}
In senso pienamente distribuzionale, le trasformate di Fourier di \(H(t)\) e \(tH(t)\)
contengono anche termini proporzionali a \(\delta(\omega)\) e alle sue derivate.
Negli appunti il punto centrale è però la forma operativa usata nei calcoli.

\section{Dal regolatore analogico al regolatore digitale}

Nelle applicazioni reali oggi non si usano quasi più regolatori interamente analogici,
ma regolatori \textbf{digitali}.

Nel caso continuo lo schema classico è:
\[
y^\star \longrightarrow \text{sommatore} \longrightarrow e \longrightarrow R(s) \longrightarrow u \longrightarrow G(s) \longrightarrow y,
\]
con retroazione negativa dell'uscita.

Nel caso digitale, invece, tra errore e impianto compaiono:
\begin{itemize}
    \item un \textbf{ADC} (convertitore analogico-digitale),
    \item un \textbf{microprocessore / PLC} che implementa il regolatore,
    \item un \textbf{DAC} (convertitore digitale-analogico).
\end{itemize}

Quindi il segnale analogico \(e(t)\) viene campionato e trasformato in una sequenza
\[
e[k]=e(kT),
\]
il controllore elabora
\[
u[k],
\]
e il DAC ricostruisce un segnale analogico da applicare all'impianto.

\subsection{Osservazione sul modello dell'impianto}

Negli appunti è sottolineato che il modello dell'impianto deve includere anche
le funzioni di trasferimento dei sensori, se queste sono rilevanti per la dinamica complessiva.

\section{Modello a tempo discreto}

Quando si campiona con periodo \(T\), lo stato viene osservato agli istanti
\[
t=kT, \qquad k\in\mathbb{Z}.
\]

Il modello discreto del sistema assume la forma
\[
x(k+1)=Ax(k)+Bu(k),
\]
dove
\[
x(k)=x(kT).
\]

Questa è la rappresentazione discreta che si usa per il progetto del regolatore digitale.

\section{Operazione di campionamento}

Dato un segnale continuo \(f(t)\), il campionamento consiste nel prelevarne il valore
agli istanti \(t=kT\), ottenendo la sequenza
\[
f[k]=f(kT).
\]

Per descrivere il campionamento nel continuo si introduce il \textbf{pettine di Dirac}
\[
\operatorname{III}_T(t)=\sum_{k=-\infty}^{+\infty}\delta(t-kT).
\]

Il segnale campionato si può allora modellare come
\[
f_s(t)=f(t)\operatorname{III}_T(t)
=
f(t)\sum_{k=-\infty}^{+\infty}\delta(t-kT).
\]

Usando la proprietà della delta si ottiene anche
\[
f_s(t)=\sum_{k=-\infty}^{+\infty} f(kT)\,\delta(t-kT).
\]

Questa scrittura mette in evidenza che il campionamento produce un treno di impulsi
pesati con i valori del segnale agli istanti di campionamento.

\section{Trasformata di Fourier del pettine di Dirac}

Una proprietà fondamentale è che la trasformata di Fourier del pettine di Dirac
è ancora un pettine di Dirac.

Se definiamo
\[
\omega_s=\frac{2\pi}{T},
\]
allora
\[
\mathcal{F}\left\{\operatorname{III}_T(t)\right\}
=
\frac{2\pi}{T}\sum_{n=-\infty}^{+\infty}\delta(\omega-n\omega_s).
\]

Quindi:
\begin{itemize}
    \item nel dominio del tempo il pettine ha passo \(T\);
    \item nel dominio della frequenza il pettine ha passo \(\omega_s=\dfrac{2\pi}{T}\).
\end{itemize}

Se invece si usa la frequenza ordinaria \(\nu=\omega/(2\pi)\), allora la periodicità in frequenza diventa
\[
f_s=\frac{1}{T}.
\]

\section{Digressione sulla convoluzione}

Per capire l'effetto del campionamento in frequenza è utile richiamare la convoluzione.

\subsection{Definizione}

Per due segnali \(f_1(t)\) e \(f_2(t)\), la convoluzione si definisce come
\[
g(t)=f_1(t)\ast f_2(t)=\int_{-\infty}^{+\infty} f_1(\tau)\,f_2(t-\tau)\,d\tau.
\]

Nel caso di segnali causali si può scrivere anche
\[
g(t)=\int_{0}^{t} f_1(\tau)\,f_2(t-\tau)\,d\tau.
\]

La convoluzione è commutativa:
\[
f_1(t)\ast f_2(t)=f_2(t)\ast f_1(t).
\]

\subsection{Uso nei sistemi dinamici}

Se \(g(t)\) è la risposta impulsiva di un sistema lineare tempo-invariante e \(u(t)\)
è l'ingresso, allora l'uscita è
\[
y(t)=u(t)\ast g(t).
\]

Se l'ingresso è la delta di Dirac,
\[
u(t)=\delta(t),
\]
allora
\[
y(t)=g(t).
\]

Questo giustifica il nome di \textbf{risposta impulsiva}.

\subsection{Convoluzione nel dominio di Laplace}

Passando nel dominio di Laplace, la convoluzione diventa prodotto:
\[
\mathcal{L}\{g(t)\}=G(s), \qquad \mathcal{L}\{u(t)\}=U(s),
\]
e quindi
\[
Y(s)=G(s)\,U(s).
\]

\section{Effetto del campionamento nel dominio della frequenza}

Poiché
\[
f_s(t)=f(t)\operatorname{III}_T(t),
\]
nel dominio della frequenza otteniamo una convoluzione:
\[
F_s(\omega)=\frac{1}{2\pi}F(\omega)\ast
\mathcal{F}\{\operatorname{III}_T(t)\}.
\]

Sostituendo la trasformata del pettine:
\[
F_s(\omega)
=
\frac{1}{2\pi}
F(\omega)\ast
\left(
\frac{2\pi}{T}\sum_{n=-\infty}^{+\infty}\delta(\omega-n\omega_s)
\right),
\]
da cui segue
\[
F_s(\omega)
=
\frac{1}{T}\sum_{n=-\infty}^{+\infty} F(\omega-n\omega_s).
\]

Questa è la formula fondamentale del campionamento:
\begin{itemize}
    \item lo spettro del segnale campionato è ottenuto replicando periodicamente
    lo spettro originale;
    \item le repliche sono centrate in multipli della frequenza di campionamento.
\end{itemize}

Nel dominio della frequenza ordinaria \(\nu\), la formula si scrive analogamente come
\[
F_s(\nu)=\frac{1}{T}\sum_{n=-\infty}^{+\infty}F\!\left(\nu-\frac{n}{T}\right).
\]

Quindi lo spettro del segnale campionato è \textbf{periodico in frequenza}
con periodo
\[
\frac{1}{T}.
\]

\section{Interpretazione grafica}

L'effetto geometrico è il seguente:
\begin{itemize}
    \item si parte da uno spettro \(F(\nu)\) centrato attorno a \(0\);
    \item il campionamento genera copie traslate di \(F(\nu)\);
    \item le copie si ripetono attorno a
    \[
    \nu=0,\ \pm \frac{1}{T},\ \pm \frac{2}{T},\dots
    \]
\end{itemize}

Se tali copie non si sovrappongono, il segnale originale può essere ricostruito.
Se invece si sovrappongono, compare aliasing.

\section{Teorema del campionamento}

Supponiamo che \(f(t)\) sia un segnale a banda limitata, cioè che il suo spettro sia nullo
al di fuori di un intervallo
\[
|\nu|\leq \bar{\nu}.
\]

Se la frequenza di campionamento soddisfa
\[
\frac{1}{T}=f_s > 2\bar{\nu},
\]
allora le repliche spettrali non si sovrappongono e il segnale può essere ricostruito
perfettamente con un filtro passa-basso ideale.

La quantità
\[
2\bar{\nu}
\]
è detta \textbf{frequenza di Nyquist} minima richiesta per evitare aliasing.

\subsection{Caso limite}

Se
\[
f_s = 2\bar{\nu},
\]
si è nel \textbf{caso critico} o \textbf{limite di Nyquist}:
le repliche si toccano esattamente ai bordi.

Nella pratica si preferisce lavorare con
\[
f_s > 2\bar{\nu},
\]
in modo da mantenere un margine di sicurezza.

\subsection{Aliasing}

Se invece
\[
f_s < 2\bar{\nu},
\]
le repliche dello spettro si sovrappongono e componenti a frequenze diverse
diventano indistinguibili dopo il campionamento.

Questo fenomeno prende il nome di \textbf{aliasing}.

\section{Significato fisico del teorema}

Il teorema del campionamento dice che un segnale continuo a banda limitata
può essere rappresentato da una sequenza discreta di campioni, purché i campioni
siano abbastanza fitti nel tempo.

In altre parole:
\begin{itemize}
    \item più il segnale contiene frequenze alte, più bisogna campionare velocemente;
    \item se si campiona troppo lentamente, si perde informazione in modo irreversibile.
\end{itemize}

\section{Sintesi della lezione}

I punti chiave della lezione sono:
\begin{enumerate}
    \item richiamo a serie di Fourier, trasformata di Fourier e trasformata di Laplace;
    \item introduzione delle distribuzioni fondamentali:
    \[
    \delta(t),\ \delta'(t),\ \delta''(t),\ H(t);
    \]
    \item relazione tra gradino e delta:
    \[
    \frac{d}{dt}H(t)=\delta(t);
    \]
    \item passaggio dal regolatore analogico al regolatore digitale;
    \item introduzione del modello discreto
    \[
    x(k+1)=Ax(k)+Bu(k);
    \]
    \item modellazione del campionamento tramite pettine di Dirac;
    \item periodicizzazione dello spettro nel dominio della frequenza;
    \item teorema del campionamento e condizione
    \[
    f_s>\,2\bar{\nu};
    \]
    \item comparsa dell'aliasing quando il campionamento è insufficiente.
\end{enumerate}
\chapter{Lezione 4 --- Teorema del campionamento, banda pratica, rumore, filtro anti-alias e modello di sampling in Laplace}

\section{Richiamo: campionamento nel tempo e periodicizzazione in frequenza}

Ripartiamo dal modello già introdotto nella lezione precedente.
Un segnale continuo \(f(t)\) viene campionato moltiplicandolo per un pettine di Dirac
di periodo \(T\):
\[
\operatorname{III}_T(t)=\sum_{k=-\infty}^{+\infty}\delta(t-kT).
\]

Il segnale campionato è quindi
\[
\hat f(t)=f(t)\,\operatorname{III}_T(t)
=
f(t)\sum_{k=-\infty}^{+\infty}\delta(t-kT)
=
\sum_{k=-\infty}^{+\infty} f(kT)\,\delta(t-kT).
\]

Dal punto di vista del tempo, questa operazione sostituisce il segnale continuo
con un treno di impulsi pesati con i valori dei campioni.

Nel dominio della frequenza, la trasformata di Fourier del segnale campionato
risulta periodica e ripete lo spettro originale con periodo
\[
\omega_s=\frac{2\pi}{T},
\]
se si lavora nel dominio delle pulsazioni, oppure con periodo
\[
f_s=\frac{1}{T},
\]
se si lavora nella frequenza ordinaria in Hertz.

Più precisamente, se \(F(\omega)\) è la trasformata di Fourier del segnale continuo,
allora
\[
\hat F(\omega)=\frac{1}{T}\sum_{n=-\infty}^{+\infty} F(\omega-n\omega_s).
\]

Quindi il campionamento:
\begin{itemize}
    \item discretizza il segnale nel tempo;
    \item replica periodicamente il suo spettro in frequenza.
\end{itemize}

\section{Ricostruzione ideale del segnale continuo}

Poiché tra dominio del tempo e dominio della frequenza esiste una relazione biunivoca
tramite trasformata e antitrasformata di Fourier, per ricostruire idealmente \(f(t)\)
a partire dal segnale campionato si procede così:
\begin{enumerate}
    \item si considera lo spettro periodico \(\hat F(\omega)\) del segnale campionato;
    \item si isola una sola copia dello spettro originale;
    \item si applica l'antitrasformata di Fourier di quella sola porzione spettrale.
\end{enumerate}

Questo procedimento è possibile solo se le repliche spettrali non si sovrappongono.

\section{Teorema del campionamento}

Supponiamo che il segnale sia a banda limitata, cioè che esista una pulsazione \(\bar\omega\)
tale che
\[
F(\omega)=0
\qquad
\text{per } |\omega|>\bar\omega.
\]

Allora il segnale è ricostruibile dai suoi campioni se la frequenza di campionamento
è sufficientemente alta da impedire la sovrapposizione delle copie periodiche dello spettro.

La condizione fondamentale è
\[
\omega_s > 2\bar\omega.
\]

Poiché
\[
\omega_s=\frac{2\pi}{T},
\]
questa condizione è equivalente a
\[
\frac{2\pi}{T}>2\bar\omega
\qquad \Longleftrightarrow \qquad
T<\frac{\pi}{\bar\omega}.
\]

Nel dominio delle frequenze ordinarie, la stessa condizione si scrive
\[
f_s>\,2\bar f.
\]

\subsection{Interpretazione}

La distanza tra le repliche dello spettro è fissata da \(1/T\) oppure da \(\omega_s\):
\begin{itemize}
    \item se \(T\) è piccolo, le repliche in frequenza sono lontane;
    \item se \(T\) è grande, le repliche in frequenza si avvicinano.
\end{itemize}

Per campionare segnali con frequenze elevate bisogna quindi ridurre il periodo di campionamento,
cioè aumentare la frequenza di campionamento.

\section{Dualismo tempo--frequenza}

Negli appunti viene richiamato un principio qualitativo molto importante:
c'è un dualismo tra localizzazione temporale e localizzazione in frequenza.

Nel caso del campionamento:
\begin{itemize}
    \item campioni molto distanziati nel tempo (cioè \(T\) grande) producono repliche
    molto vicine in frequenza;
    \item campioni molto fitti nel tempo (cioè \(T\) piccolo) producono repliche
    molto lontane in frequenza.
\end{itemize}

Questo spiega geometricamente perché la ricostruzione peggiora quando si campiona troppo lentamente.

\section{Segnali a banda limitata: caso ideale e caso pratico}

Dal punto di vista matematico, un segnale \emph{strettamente} a banda limitata
ha trasformata di Fourier nulla oltre una certa frequenza di taglio.

Tuttavia un segnale davvero a banda limitata è un modello ideale:
in molti casi reali il segnale non si esaurisce perfettamente in frequenza,
ma presenta code spettrali che decadono progressivamente.

Per questo, in pratica, si introduce una \textbf{banda passante effettiva}
o \textbf{banda pratica}, definita tramite una soglia convenzionale.

\section{Definizione pratica di banda: il criterio dei \(-3\) dB}

Negli appunti compare il richiamo al valore \(-3\) dB.

Se \(F(\omega)\) ha un massimo a bassa frequenza, si può definire la banda passante
come l'intervallo di frequenze per cui il modulo resta sopra una soglia prefissata.
Una scelta molto comune è la soglia a \(-3\) dB rispetto al massimo.

In questo modo si definisce una pulsazione di banda \(\bar\omega\) come il punto
in cui il modulo è sceso di \(3\) dB rispetto al valore di riferimento.

\subsection{Osservazione importante}

La soglia dei \(-3\) dB non è una legge assoluta:
\begin{itemize}
    \item è una convenzione pratica molto diffusa;
    \item in altri casi si può usare una soglia diversa, ad esempio \(-6\) dB;
    \item la scelta va fatta in base al sistema e alla precisione desiderata.
\end{itemize}

Quindi, per i segnali reali, la nozione di banda è spesso una nozione \emph{operativa}
più che rigorosamente matematica.

\section{Rumore di misura}

Quando si misura un segnale fisico con un sensore, la misura non è mai esatta.
Si può modellare il segnale misurato come
\[
f_m(t)=f(t)+m(t),
\]
dove:
\begin{itemize}
    \item \(f(t)\) è il segnale ``vero'';
    \item \(m(t)\) è il rumore di misura.
\end{itemize}

Il rumore dipende da:
\begin{itemize}
    \item sensibilità del sensore;
    \item elettronica di misura;
    \item ambiente;
    \item processo di acquisizione.
\end{itemize}

Per questo, nel problema del campionamento, non si ha quasi mai a che fare
con un segnale puro, ma con un segnale contaminato dal rumore.

\section{Rumore bianco}

Negli appunti compare il richiamo al \textbf{processo bianco}.

\subsection{Autocorrelazione}

Un processo di rumore bianco ideale ha autocorrelazione del tipo
\[
R_m(\tau)=\mathbb{E}\{m(t)m(t+\tau)\}=K\,\delta(\tau),
\]
dove \(K\) è una costante positiva.

Questo significa che campioni presi a istanti diversi sono incorrelati.

\subsection{Spettro del rumore bianco}

La trasformata di Fourier della delta è costante, quindi la densità spettrale
del rumore bianco è piatta:
\[
S_m(\omega)=\mathcal{F}\{R_m(\tau)\}=K.
\]

Per questo si dice che un rumore bianco ha energia distribuita su tutte le frequenze.

\subsection{Rumore bianco e rumore gaussiano}

Le due proprietà non coincidono:
\begin{itemize}
    \item \textbf{bianco} significa spettro piatto / campioni incorrelati;
    \item \textbf{gaussiano} si riferisce invece alla distribuzione di probabilità.
\end{itemize}

Un rumore può essere gaussiano ma non bianco, oppure bianco ma non gaussiano.

\section{Perché il rumore crea problemi nel campionamento}

Se il segnale misurato è
\[
f_m(t)=f(t)+m(t),
\]
allora il suo spettro contiene:
\begin{itemize}
    \item la parte utile \(F(\omega)\);
    \item la parte di rumore \(M(\omega)\).
\end{itemize}

Quando si campiona, anche lo spettro del rumore viene replicato periodicamente.
Poiché il rumore bianco è presente su tutte le frequenze, le repliche possono ripiegarsi
sulle basse frequenze e contaminare la banda utile.

Questo fenomeno è una delle ragioni principali per cui, \emph{prima del campionamento},
si introduce un filtro anti-alias.

\section{Filtro anti-alias}

Prima di campionare un segnale reale è quindi necessario inserire un
\textbf{filtro anti-alias}, cioè un filtro analogico posto a monte del campionatore.

Lo schema concettuale è:
\[
f(t)\ \longrightarrow\ \text{filtro anti-alias}\ \longrightarrow\ \text{campionatore}\ \longrightarrow\ \hat f(kT).
\]

\subsection{Funzione del filtro anti-alias}

Il filtro anti-alias deve:
\begin{itemize}
    \item attenuare le componenti ad alta frequenza;
    \item limitare il contenuto spettrale che, dopo il campionamento,
    si ripiegherebbe nella banda utile;
    \item rendere applicabile in modo pratico il teorema del campionamento.
\end{itemize}

Idealmente sarebbe un passa-basso rettangolare, ma nella realtà si usano filtri
con pendenza finita.

\section{Caratteristiche del filtro anti-alias}

Il filtro anti-alias è tipicamente un filtro passa-basso analogico.

Un esempio semplice è un filtro RC del primo ordine, con pendenza asintotica
\[
-20\ \text{dB/decade}.
\]

Filtri di ordine più elevato permettono pendenze maggiori, ma introducono
in generale maggiore complessità e maggiore sfasamento.

\subsection{Compromesso attenuazione--sfasamento}

Un punto essenziale della lezione è che il filtro anti-alias non può essere scelto
guardando solo l'attenuazione.

Infatti:
\begin{itemize}
    \item una pendenza più rapida attenua meglio le alte frequenze;
    \item però una pendenza più rapida introduce più sfasamento;
    \item lo sfasamento può degradare la dinamica del sistema e, in alcuni casi,
    compromettere anche i margini di stabilità.
\end{itemize}

La scelta del filtro è quindi sempre un compromesso tra:
\begin{itemize}
    \item quanto rumore / aliasing si vuole sopprimere;
    \item quanto ritardo di fase si può tollerare.
\end{itemize}

\section{Scelta pratica di frequenza di taglio e frequenza di campionamento}

Nel caso ideale il teorema di campionamento richiederebbe soltanto
\[
\omega_s>2\bar\omega.
\]

Nella pratica però, poiché il filtro anti-alias non è ideale, si preferisce lavorare
con un margine di sicurezza:
\begin{itemize}
    \item si sceglie la frequenza di campionamento ben più alta del doppio della banda utile;
    \item si colloca la frequenza di taglio del filtro anti-alias in modo coerente
    con la banda del segnale e con la pendenza disponibile.
\end{itemize}

Negli appunti è richiamata una regola pratica secondo cui conviene scegliere
\[
\omega_c \gg 2\bar\omega,
\]
oppure, in forma equivalente,
\[
f_s \text{ sufficientemente maggiore di } 2\bar f.
\]

Una regola empirica spesso usata in pratica è quella di scegliere la frequenza di campionamento
diverse volte più grande della banda utile, ad esempio
\[
f_s \ge 5\cdot (2\bar f)
\quad \text{oppure} \quad
f_s \ge 10\cdot (2\bar f),
\]
a seconda delle prestazioni richieste e della pendenza del filtro disponibile.

\subsection{Osservazione}

Se si sbaglia la scelta della frequenza di campionamento oppure della frequenza di taglio
del filtro anti-alias, l'informazione persa per aliasing non può più essere recuperata a valle.

\section{Sensori digitali e sensori analogici}

Molti sensori reali sono ancora analogici, e in tal caso:
\begin{itemize}
    \item il filtraggio anti-alias;
    \item il campionamento;
\end{itemize}
vengono effettuati nel sistema di acquisizione.

Se invece il sensore è digitale, spesso:
\begin{itemize}
    \item il costruttore incorpora già una fase di filtraggio;
    \item il periodo di campionamento è prefissato oppure regolabile in modo limitato.
\end{itemize}

In ogni caso, dal punto di vista progettuale, rimane centrale capire quali siano
le caratteristiche in frequenza del filtro anti-alias effettivo.

\section{Passaggio dal modello in Fourier al modello in Laplace}

Finora il campionamento è stato descritto soprattutto nel dominio della Fourier,
cioè guardando la periodicizzazione dello spettro.

Tuttavia, per l'analisi e il progetto dei sistemi di controllo, è più comodo lavorare
nel dominio di Laplace, perché:
\begin{itemize}
    \item i modelli dei sistemi dinamici continui sono dati in \(s\);
    \item il regolatore e l'impianto vengono normalmente descritti tramite
    funzioni di trasferimento;
    \item il campionamento interessa solo una parte dell'anello di controllo.
\end{itemize}

Per questo si introduce un modello di sampling compatibile con la trasformata di Laplace.

\section{Modello di sampling nel dominio di Laplace}

Per un segnale causale campionato agli istanti \(kT\), si considera
\[
\hat f(t)=\sum_{k=0}^{+\infty} f(kT)\,\delta(t-kT).
\]

Calcoliamo la sua trasformata di Laplace:
\[
\mathcal{L}\{\hat f(t)\}
=
\int_{0^-}^{+\infty}
\left(
\sum_{k=0}^{+\infty} f(kT)\,\delta(t-kT)
\right)e^{-st}\,dt.
\]

Portando fuori la sommatoria,
\[
\mathcal{L}\{\hat f(t)\}
=
\sum_{k=0}^{+\infty} f(kT)\,\mathcal{L}\{\delta(t-kT)\}.
\]

Poiché
\[
\mathcal{L}\{\delta(t-kT)\}=e^{-skT},
\]
si ottiene
\[
\mathcal{L}\{\hat f(t)\}
=
\sum_{k=0}^{+\infty} f(kT)e^{-skT}.
\]

Questa non è una funzione razionale in \(s\), perché contiene l'esponenziale \(e^{-sT}\).

\section{Introduzione della trasformata \(z\)}

Per trattare comodamente il campionamento si introduce la variabile
\[
z=e^{sT}.
\]

Di conseguenza
\[
e^{-sT}=z^{-1}.
\]

Sostituendo nella formula precedente,
\[
\mathcal{L}\{\hat f(t)\}
=
\sum_{k=0}^{+\infty} f(kT)\,z^{-k}.
\]

Questa espressione suggerisce la definizione della \textbf{trasformata \(z\)} del segnale campionato:
\[
\mathcal{Z}\{f[k]\}
=
\sum_{k=0}^{+\infty} f(kT)\,z^{-k}.
\]

Quindi la trasformata \(z\) nasce come strumento naturale per descrivere
segnali e sistemi campionati.

\section{Significato della trasformata \(z\)}

La variabile \(z\) incorpora in sé l'informazione del campionamento,
perché è legata al periodo \(T\) attraverso la relazione
\[
z=e^{sT}.
\]

La trasformata \(z\):
\begin{itemize}
    \item svolge per i sistemi discreti un ruolo analogo a quello della Laplace
    per i sistemi continui;
    \item permette di passare da equazioni alle differenze a funzioni di trasferimento discrete;
    \item è lo strumento naturale per lo studio dei regolatori digitali.
\end{itemize}

\section{Sintesi della lezione}

I punti fondamentali della lezione sono:
\begin{enumerate}
    \item richiamo del teorema del campionamento e della condizione
    \[
    \omega_s>2\bar\omega
    \qquad\text{oppure}\qquad
    T<\frac{\pi}{\bar\omega};
    \]
    \item interpretazione del dualismo tempo--frequenza;
    \item distinzione tra banda ideale e banda pratica;
    \item uso del criterio a \(-3\) dB per definire una banda passante operativa;
    \item introduzione del rumore di misura e del rumore bianco;
    \item necessità del filtro anti-alias prima del campionamento;
    \item compromesso tra attenuazione delle alte frequenze e sfasamento introdotto dal filtro;
    \item passaggio dal modello di sampling in Fourier al modello di sampling in Laplace;
    \item nascita della trasformata \(z\) dalla relazione
    \[
    z=e^{sT}.
    \]
\end{enumerate}
\chapter{Lezione 5 --- Trasformata z, esempi elementari e corrispondenza tra piano \(s\) e piano \(z\)}

\section{Richiamo dalla lezione precedente}

Nella lezione precedente abbiamo concluso che, se un segnale continuo causale \(f(t)\)
viene campionato con periodo \(T\), allora il segnale campionato in forma impulsiva può essere scritto come
\[
\hat f(t)=\sum_{k=0}^{+\infty} f(kT)\,\delta(t-kT).
\]

Applicando la trasformata di Laplace si ottiene
\[
\mathcal{L}\{\hat f(t)\}
=
\int_{0^-}^{+\infty}
\left(
\sum_{k=0}^{+\infty} f(kT)\,\delta(t-kT)
\right)e^{-st}\,dt.
\]

Scambiando somma e integrale,
\[
\mathcal{L}\{\hat f(t)\}
=
\sum_{k=0}^{+\infty} f(kT)\int_{0^-}^{+\infty}\delta(t-kT)e^{-st}\,dt
=
\sum_{k=0}^{+\infty} f(kT)e^{-skT}.
\]

Questa espressione è una successione di termini numerici pesati da \(e^{-skT}\),
e suggerisce naturalmente l'introduzione di una nuova variabile complessa.

\section{Nascita della trasformata \(z\)}

Poiché l'esponenziale \(e^{-skT}\) non è particolarmente comodo da manipolare,
si introduce la variabile
\[
z=e^{sT},
\qquad\text{per cui}\qquad
e^{-sT}=z^{-1}.
\]

La trasformata di Laplace del segnale campionato diventa allora
\[
\mathcal{L}\{\hat f(t)\}
=
\sum_{k=0}^{+\infty} f(kT)z^{-k}.
\]

Questa espressione motiva la definizione della \textbf{trasformata \(z\)}.

\section{Definizione della trasformata \(z\)}

\subsection{Trasformata \(z\) monolatera}

Per una sequenza discreta
\[
x[k], \qquad k=0,1,2,\dots
\]
si definisce la trasformata \(z\) monolatera come
\[
X(z)=\mathcal{Z}\{x[k]\}=\sum_{k=0}^{+\infty} x[k]\,z^{-k}.
\]

Nel nostro contesto,
\[
x[k]=f(kT),
\]
quindi
\[
\mathcal{Z}\{f(kT)\}=\sum_{k=0}^{+\infty} f(kT)\,z^{-k}.
\]

La trasformata \(z\) monolatera è l'analogo discreto della trasformata di Laplace unilaterale:
parte da \(k=0\) e arriva a \(+\infty\), quindi è adatta a segnali causali e transitori.

\subsection{Trasformata \(z\) bilatera}

Per sequenze definite su tutti gli interi, si introduce invece la trasformata \(z\) bilatera:
\[
X(z)=\mathcal{Z}_b\{x[k]\}=\sum_{k=-\infty}^{+\infty} x[k]\,z^{-k}.
\]

Questa forma è l'analogo discreto della trasformata di Fourier/Laplace bilatera
ed è utile per segnali stazionari o definiti su tutto l'asse discreto.

\subsection{Osservazione}

La sequenza \(x[k]\) è discreta, ma \(X(z)\) è una funzione della variabile complessa
\[
z\in\mathbb{C}.
\]
Quindi la trasformata \(z\) non è un oggetto discreto:
\[
X(z)
\]
è una funzione continua della variabile complessa \(z\), definita sul proprio dominio di convergenza.

\section{Interpretazione di \(z^{-1}\) e \(z\)}

Negli appunti compare una lettura operativa molto utile:
\begin{itemize}
    \item \(z^{-1}\) si interpreta come \textbf{operatore di ritardo unitario};
    \item \(z\) si interpreta formalmente come \textbf{operatore di anticipo unitario}.
\end{itemize}

Questa interpretazione è fondamentale nei sistemi discreti e nei diagrammi a blocchi.

\section{Sequenze equispaziate}

La trasformata \(z\) vale per sequenze discrete i cui campioni siano equispaziati nel tempo,
cioè prelevati con periodo costante \(T\).

Dire che i campioni sono equispaziati significa che tra due campioni consecutivi
passa sempre lo stesso intervallo di tempo:
\[
t_k = kT.
\]

Questo è il caso naturale che si incontra nel campionamento uniforme.

\section{Proprietà della trasformata \(z\)}

La trasformata \(z\) monolatera gode di proprietà del tutto analoghe a quelle della Laplace.

\subsection{Linearità}

Per \(\alpha,\beta\in\mathbb{R}\) oppure \(\mathbb{C}\),
\[
\mathcal{Z}\{\alpha x[k]+\beta y[k]\}
=
\alpha X(z)+\beta Y(z).
\]

\subsection{Convoluzione discreta}

Se \(x[k]\) e \(y[k]\) sono due sequenze causali, la loro convoluzione discreta è
\[
(x\ast y)[n]
=
\sum_{k=0}^{n} x[k]\,y[n-k].
\]

La trasformata \(z\) trasforma la convoluzione in prodotto:
\[
\mathcal{Z}\{x\ast y\}=X(z)Y(z).
\]

Questa proprietà è perfettamente analoga a quella già vista per Laplace nel caso continuo.

\section{Primo esempio: trasformata \(z\) del gradino}

Consideriamo il gradino continuo
\[
f(t)=1(t)=H(t),
\]
e adottiamo la convenzione coerente con gli appunti secondo cui il segnale campionato vale
\[
f(kT)=1,
\qquad k=0,1,2,\dots
\]

Allora
\[
X(z)=\sum_{k=0}^{+\infty} 1\cdot z^{-k}
=
\sum_{k=0}^{+\infty}(z^{-1})^k.
\]

Questa è una serie geometrica, che converge se
\[
|z^{-1}|<1
\qquad\Longleftrightarrow\qquad
|z|>1.
\]

Quindi
\[
X(z)=\frac{1}{1-z^{-1}}
=
\frac{z}{z-1},
\qquad
\text{ROC: } |z|>1.
\]

\subsection{Osservazioni}

\begin{itemize}
    \item La trasformata è una funzione razionale.
    \item C'è un polo in \(z=1\).
    \item C'è uno zero in \(z=0\), se si guarda la forma \(\dfrac{z}{z-1}\).
\end{itemize}

Questo mostra che, anche nel dominio \(z\), possiamo continuare a ragionare in termini di poli,
zeri e funzioni razionali.

\section{Regione di convergenza}

Come nella trasformata di Laplace, anche la trasformata \(z\) non è caratterizzata solo
dalla sua espressione algebrica, ma anche dalla \textbf{regione di convergenza} (ROC).

Nel primo esempio abbiamo ottenuto
\[
X(z)=\frac{z}{z-1},
\qquad \text{ROC: } |z|>1.
\]

Questo significa che la funzione è definita e utile nella regione esterna al cerchio di raggio \(1\),
escluso il polo \(z=1\).

\subsection{Osservazione sul dominio}

Negli appunti compare l'idea, mutuata dall'analisi matematica, di studiare i punti di bordo
del dominio guardando i limiti dai punti in cui la funzione è effettivamente definita.

Questa osservazione serve a ricordare che:
\begin{itemize}
    \item una funzione razionale può essere prolungata per continuità in un punto
    solo se il limite esiste ed è finito;
    \item in presenza di un polo il prolungamento non è possibile;
    \item quindi il punto resta escluso dal dominio.
\end{itemize}

\section{Secondo esempio: trasformata \(z\) di un esponenziale campionato}

Consideriamo il segnale continuo causale
\[
f(t)=e^{-at}1(t),
\qquad a>0.
\]

Campionandolo agli istanti \(t=kT\), otteniamo
\[
f(kT)=e^{-akT},
\qquad k=0,1,2,\dots
\]

La trasformata \(z\) è allora
\[
X(z)=\sum_{k=0}^{+\infty} e^{-akT}z^{-k}
=
\sum_{k=0}^{+\infty}\left(e^{-aT}z^{-1}\right)^k.
\]

La serie converge se
\[
|e^{-aT}z^{-1}|<1
\qquad\Longleftrightarrow\qquad
|z|>e^{-aT}.
\]

Quindi
\[
X(z)=\frac{1}{1-e^{-aT}z^{-1}}
=
\frac{z}{z-e^{-aT}},
\qquad
\text{ROC: } |z|>e^{-aT}.
\]

\subsection{Interpretazione}

La trasformata ha:
\begin{itemize}
    \item un polo in
    \[
    z=e^{-aT},
    \]
    \item uno zero in \(z=0\), nella forma razionale in \(z\).
\end{itemize}

Poiché \(a>0\), si ha
\[
0<e^{-aT}<1,
\]
quindi il polo cade all'interno del cerchio unitario.

\section{Funzioni razionali nel dominio \(z\)}

Un punto sottolineato negli appunti è che la trasformata \(z\) dei segnali elementari
porta ancora a funzioni razionali.

Questo è importante perché permette di mantenere, anche nel discreto,
molti degli strumenti concettuali usati nel continuo:
\begin{itemize}
    \item poli;
    \item zeri;
    \item fattorizzazioni;
    \item studio della stabilità;
    \item diagrammi a blocchi.
\end{itemize}

\subsection{Osservazione tecnica}

Nel dominio \(z\) occorre sempre essere attenti alla convenzione usata:
una funzione può essere scritta come polinomio in \(z\) oppure come polinomio in \(z^{-1}\).
Per questo alcune considerazioni su grado del numeratore e del denominatore
vanno interpretate con cautela, specificando sempre quale forma si sta usando.

\section{Corrispondenza tra piano \(s\) e piano \(z\)}

La relazione fondamentale tra piano continuo e piano discreto è
\[
z=e^{sT}.
\]

Se
\[
s=\sigma+j\omega,
\]
allora
\[
z=e^{(\sigma+j\omega)T}=e^{\sigma T}e^{j\omega T}.
\]

Questa formula contiene tutta la corrispondenza geometrica tra i due piani.

\subsection{Caso \(\omega=0\)}

Se \(\omega=0\), allora
\[
s=\sigma
\qquad\Longrightarrow\qquad
z=e^{\sigma T}.
\]

Quindi:
\begin{itemize}
    \item se \(\sigma=0\), allora \(z=1\);
    \item se \(\sigma>0\), allora \(z=e^{\sigma T}>1\), cioè si va sull'asse reale positivo
    all'esterno del cerchio unitario;
    \item se \(\sigma<0\), allora \(0<z=e^{\sigma T}<1\), cioè si va sull'asse reale positivo
    all'interno del cerchio unitario;
    \item se \(\sigma\to -\infty\), allora \(z\to 0\).
\end{itemize}

\subsection{Caso \(\sigma=0\)}

Se \(\sigma=0\), allora
\[
s=j\omega
\qquad\Longrightarrow\qquad
z=e^{j\omega T}=\cos(\omega T)+j\sin(\omega T).
\]

Quindi l'asse immaginario del piano \(s\) viene mappato sul cerchio unitario nel piano \(z\):
\[
|z|=1.
\]

\subsection{Osservazione importante}

Questa mappa non è biunivoca su tutto l'asse immaginario:
infatti
\[
e^{j\omega T}=e^{j(\omega+2k\pi/T)T},
\qquad k\in\mathbb{Z}.
\]

Cioè lo stesso punto del piano \(z\) corrisponde a infiniti valori di \(\omega\)
che differiscono di multipli di \(2\pi/T\).

\section{Caso generale}

Nel caso generale
\[
z=e^{\sigma T}e^{j\omega T}.
\]

Da qui si leggono due proprietà geometriche fondamentali:

\begin{itemize}
    \item fissato \(\sigma\), il modulo è
    \[
    |z|=e^{\sigma T},
    \]
    quindi le rette verticali \(\sigma=\text{costante}\) nel piano \(s\)
    si trasformano in circonferenze concentriche nel piano \(z\);

    \item fissato \(\omega\), l'argomento è
    \[
    \arg(z)=\omega T \pmod{2\pi},
    \]
    quindi le rette orizzontali \(\omega=\text{costante}\) nel piano \(s\)
    si trasformano in semirette uscenti dall'origine nel piano \(z\).
\end{itemize}

\section{Conseguenze fondamentali della mappa \(z=e^{sT}\)}

Dalla formula precedente segue immediatamente che:

\begin{itemize}
    \item l'asse immaginario del piano \(s\) diventa il cerchio unitario del piano \(z\);
    \item il semipiano sinistro
    \[
    \Re(s)<0
    \]
    viene mappato all'interno del cerchio unitario
    \[
    |z|<1;
    \]
    \item il semipiano destro
    \[
    \Re(s)>0
    \]
    viene mappato all'esterno del cerchio unitario
    \[
    |z|>1;
    \]
    \item l'origine del piano \(s\), cioè \(s=0\), viene mappata nel punto
    \[
    z=1.
    \]
\end{itemize}

\section{Fascia primaria}

Per rendere biunivoca la corrispondenza tra \(s\) e \(z\), bisogna limitarsi a una fascia
del piano \(s\) di ampiezza \(2\pi/T\) rispetto alla variabile \(\omega\).

Una scelta standard è la \textbf{fascia primaria}
\[
-\frac{\pi}{T}<\omega\le \frac{\pi}{T}.
\]

Se ci si limita a questa fascia, la mappa
\[
z=e^{sT}
\]
diventa univoca.

\subsection{Interpretazione}

La massima pulsazione contenuta nella fascia primaria è
\[
\omega_{\max}=\frac{\pi}{T}.
\]

Se definiamo la pulsazione di campionamento
\[
\omega_c=\omega_s=\frac{2\pi}{T},
\]
allora
\[
\omega_{\max}=\frac{\omega_c}{2}.
\]

Questa è un'altra forma del teorema del campionamento:
la fascia primaria contiene le frequenze non aliasate.

\section{Legame con il teorema di Shannon}

Poiché la fascia primaria è
\[
-\frac{\pi}{T}<\omega\le \frac{\pi}{T},
\]
per evitare aliasing bisogna fare in modo che il contenuto in frequenza del segnale
ricada interamente dentro questa fascia.

Equivalentemente:
\[
\omega_b < \frac{\pi}{T}
\qquad \Longleftrightarrow \qquad
\omega_s=\frac{2\pi}{T}>2\omega_b.
\]

Questa è esattamente la formulazione del teorema di Shannon nel dominio delle pulsazioni.

\section{Mappatura di poli e singolarità}

Un altro punto importante richiamato negli appunti è che le singolarità della trasformata
nel dominio continuo vengono mappate nei corrispondenti punti del piano \(z\).

Se un sistema continuo ha un polo in
\[
s=p,
\]
allora il corrispondente punto nel dominio discreto è
\[
z=e^{pT}.
\]

Questo significa che:
\begin{itemize}
    \item poli continui nel semipiano sinistro vanno all'interno del cerchio unitario;
    \item poli continui sull'asse immaginario vanno sul cerchio unitario;
    \item poli continui nel semipiano destro vanno all'esterno del cerchio unitario.
\end{itemize}

Questa corrispondenza è alla base del criterio di stabilità per i sistemi discreti.

\section{Sintesi della lezione}

I punti principali della lezione sono:
\begin{enumerate}
    \item derivazione della trasformata \(z\) dalla trasformata di Laplace del segnale campionato;
    \item definizione di trasformata \(z\) monolatera e bilatera;
    \item interpretazione di \(z^{-1}\) come operatore di ritardo;
    \item proprietà di linearità e convoluzione;
    \item calcolo della trasformata \(z\) del gradino:
    \[
    \mathcal{Z}\{1\}=\frac{z}{z-1},\qquad |z|>1;
    \]
    \item calcolo della trasformata \(z\) dell'esponenziale campionato:
    \[
    \mathcal{Z}\{e^{-akT}\}=\frac{z}{z-e^{-aT}},\qquad |z|>e^{-aT};
    \]
    \item corrispondenza tra piano \(s\) e piano \(z\) tramite
    \[
    z=e^{sT};
    \]
    \item asse immaginario di \(s\) mappato sul cerchio unitario;
    \item semipiano sinistro mappato all'interno del cerchio unitario;
    \item introduzione della fascia primaria
    \[
    -\frac{\pi}{T}<\omega\le \frac{\pi}{T},
    \]
    che rappresenta la porzione non aliasata del piano \(s\).
\end{enumerate}
\chapter{Lezione 6 --- Trasformata z: propriet\`a, tabelle fondamentali e funzione di trasferimento dei sistemi discreti}

\section{Richiamo: trasformata z della successione campionata}

Riprendiamo il risultato della lezione precedente. Se un segnale continuo causale \(f(t)\)
viene campionato con periodo \(T\), allora il segnale campionato in forma impulsiva si pu\`o scrivere come
\[
\hat f(t)=\sum_{k=0}^{+\infty} f(kT)\,\delta(t-kT).
\]

Applicando la trasformata di Laplace si ottiene
\[
\mathcal{L}\{\hat f(t)\}
=
\int_{0^-}^{+\infty}
\left(
\sum_{k=0}^{+\infty} f(kT)\,\delta(t-kT)
\right)e^{-st}\,dt
=
\sum_{k=0}^{+\infty} f(kT)e^{-skT}.
\]

Introducendo la variabile
\[
z=e^{sT},
\qquad\text{quindi}\qquad
e^{-sT}=z^{-1},
\]
si arriva alla definizione di trasformata z monolatera:
\[
X(z)=\mathcal{Z}\{x(k)\}
=
\sum_{k=0}^{+\infty} x(k)\,z^{-k}.
\]

Nel nostro caso \(x(k)=f(kT)\), quindi pi\`u rigorosamente si dovrebbe scrivere
\[
\mathcal{Z}\{f(kT)\}
=
\sum_{k=0}^{+\infty} f(kT)\,z^{-k}.
\]

\subsection*{Osservazione importante}

Nella trasformata z la dipendenza esplicita da \(T\) non compare pi\`u nella formula,
ma \`e \emph{nascosta} nella variabile \(z=e^{sT}\).
Per questo due successioni formalmente uguali possono corrispondere a segnali continui
campionati con periodi diversi.

\section{Trasformata z monolatera e bilatera}

Per segnali causali o comunque definiti per \(k=0,1,2,\dots\), si usa la trasformata z \textbf{monolatera}:
\[
X(z)=\sum_{k=0}^{+\infty} x(k)\,z^{-k}.
\]

Per segnali definiti su tutti gli interi, si pu\`o introdurre la trasformata z \textbf{bilatera}:
\[
X(z)=\sum_{k=-\infty}^{+\infty} x(k)\,z^{-k}.
\]

Nel seguito lavoriamo prevalentemente con la forma monolatera, perch\'e stiamo studiando
sistemi causali e sequenze con transitorio.

\section{Propriet\`a della trasformata z}

La trasformata z monolatera eredita molte propriet\`a della trasformata di Laplace.

\subsection{Linearit\`a}

Per \(\alpha,\beta\in\mathbb{R}\) oppure \(\mathbb{C}\),
\[
\mathcal{Z}\{\alpha x(k)+\beta y(k)\}
=
\alpha X(z)+\beta Y(z).
\]

\subsection{Convoluzione discreta}

Se definiamo la convoluzione discreta di due successioni causali come
\[
(x\ast y)(n)=\sum_{k=0}^{n} x(k)\,y(n-k),
\]
allora vale
\[
\mathcal{Z}\{x\ast y\}=X(z)Y(z).
\]

\subsection{Traslazione temporale}

Supponiamo di avere la successione
\[
x(k)\ \longleftrightarrow\ X(z)
\]
e consideriamo la successione anticipata di un passo
\[
\tilde x(k)=x(k+1).
\]

Allora
\[
\mathcal{Z}\{x(k+1)\}=zX(z)-zx(0).
\]

Questa propriet\`a \`e l'analogo discreto della propriet\`a della derivata
nella trasformata di Laplace.

\subsection{Teorema del valore iniziale}

Se la trasformata z esiste, allora
\[
x(0)=\lim_{z\to\infty} X(z).
\]

\subsection{Teorema del valore finale}

Se esiste il limite \(\lim_{k\to+\infty}x(k)\) e se sono soddisfatte le usuali ipotesi
sui poli di \((z-1)X(z)\), allora
\[
\lim_{k\to+\infty}x(k)
=
\lim_{z\to 1}(z-1)X(z).
\]

\section{Coppie fondamentali tra trasformata z, segnale campionato e trasformata di Laplace}

Raccogliamo alcune coppie molto usate negli esercizi.
Nella tabella seguente:
\begin{itemize}
    \item la prima colonna contiene la trasformata z della sequenza campionata;
    \item la seconda colonna il corrispondente segnale continuo causale;
    \item la terza colonna la relativa trasformata di Laplace.
\end{itemize}

\[
\renewcommand{\arraystretch}{1.5}
\begin{array}{|c|c|c|}
\hline
X(z) & x(t) & X(s) \\
\hline
1 & \delta(t) & 1 \\
\hline
\dfrac{z}{z-1} & 1(t) & \dfrac{1}{s} \\
\hline
\dfrac{Tz}{(z-1)^2} & t\,1(t) & \dfrac{1}{s^2} \\
\hline
\dfrac{T^2 z(z+1)}{2(z-1)^3} & \dfrac{1}{2}t^2\,1(t) & \dfrac{1}{s^3} \\
\hline
\dfrac{z}{z-e^{-aT}} & e^{-at}1(t) & \dfrac{1}{s+a} \\
\hline
\dfrac{Tze^{-aT}}{(z-e^{-aT})^2} & t\,e^{-at}1(t) & \dfrac{1}{(s+a)^2} \\
\hline
\dfrac{z\sin(\omega T)}{z^2-2\cos(\omega T)z+1}
& \sin(\omega t)\,1(t)
& \dfrac{\omega}{s^2+\omega^2}
\\
\hline
\dfrac{z\bigl(z-\cos(\omega T)\bigr)}{z^2-2\cos(\omega T)z+1}
& \cos(\omega t)\,1(t)
& \dfrac{s}{s^2+\omega^2}
\\
\hline
\dfrac{ze^{-aT}\sin(\omega T)}
{z^2-2e^{-aT}\cos(\omega T)z+e^{-2aT}}
& e^{-at}\sin(\omega t)\,1(t)
& \dfrac{\omega}{(s+a)^2+\omega^2}
\\
\hline
\end{array}
\]

\subsection*{Osservazione sui poli complessi coniugati}

Per il segnale
\[
e^{-at}\sin(\omega t)\,1(t),
\]
i poli della trasformata z sono
\[
z_{1,2}=e^{-aT}\bigl(\cos(\omega T)\pm j\sin(\omega T)\bigr)
=
e^{-aT}e^{\pm j\omega T}.
\]

Quindi:
\begin{itemize}
    \item il modulo dei poli \`e \(e^{-aT}\);
    \item l'angolo dei poli \`e \(\pm \omega T\).
\end{itemize}

\section{Sistema in forma di stato discreta}

Consideriamo ora un sistema lineare tempo-invariante nel discreto, in forma di stato:
\[
\begin{cases}
x(k+1)=Ax(k)+Bu(k),\\[4pt]
y(k)=Cx(k),
\end{cases}
\]
con condizione iniziale \(x(0)\).

Applicando la trasformata z monolatera all'equazione di stato, e usando la propriet\`a
di traslazione, si ottiene
\[
zX(z)-zx(0)=AX(z)+BU(z).
\]

Dall'equazione di uscita segue invece
\[
Y(z)=CX(z).
\]

Riordinando l'equazione di stato:
\[
(zI-A)X(z)=BU(z)+zx(0).
\]

Quindi
\[
X(z)=(zI-A)^{-1}BU(z)+(zI-A)^{-1}zx(0).
\]

Sostituendo nell'equazione di uscita:
\[
Y(z)=C(zI-A)^{-1}BU(z)+C(zI-A)^{-1}zx(0).
\]

\subsection{Interpretazione dei due termini}

I due termini hanno un significato fisico ben preciso:
\[
Y(z)=
\underbrace{C(zI-A)^{-1}BU(z)}_{\text{risposta forzata}}
+
\underbrace{C(zI-A)^{-1}zx(0)}_{\text{risposta libera}}.
\]

Il primo dipende dall'ingresso, il secondo dalla condizione iniziale.

\section{Matrice di trasferimento}

In assenza di condizioni iniziali, cio\`e ponendo \(x(0)=0\), resta
\[
Y(z)=C(zI-A)^{-1}BU(z).
\]

Da qui si definisce la \textbf{matrice di trasferimento}
\[
G(z)=C(zI-A)^{-1}B.
\]

Se il sistema ha anche un termine diretto \(D\), la formula generale diventa
\[
G(z)=C(zI-A)^{-1}B+D.
\]

Nel caso trattato negli appunti si assume \(D=0\).

\section{Calcolo di \((zI-A)^{-1}\) mediante aggiunta e determinante}

Per calcolare esplicitamente la matrice di trasferimento, pu\`o essere utile richiamare
la formula generale dell'inversa di una matrice invertibile.

Se \(M\in\mathbb{R}^{n\times n}\) \`e invertibile, allora
\[
M^{-1}=\frac{1}{\det(M)}\operatorname{Adj}(M),
\]
dove \(\operatorname{Adj}(M)\) \`e la matrice aggiunta di \(M\).

L'elemento \((i,j)\) dell'aggiunta si ottiene a partire dal complemento algebrico:
\[
c_{ij}=(-1)^{i+j}\det(M^{\ast}_{ji}),
\]
dove \(M^{\ast}_{ji}\) \`e il minore ottenuto eliminando la riga \(j\) e la colonna \(i\).

Applicando questa formula a \(M=zI-A\), si ottiene
\[
(zI-A)^{-1}
=
\frac{1}{\det(zI-A)}\operatorname{Adj}(zI-A).
\]

Quindi la matrice di trasferimento pu\`o essere scritta come
\[
G(z)=C(zI-A)^{-1}B
=
\frac{1}{\det(zI-A)}\,C\,\operatorname{Adj}(zI-A)\,B.
\]

\section{Caso SISO}

Nel caso \textbf{SISO} (Single Input Single Output),
\[
u(k)\in\mathbb{R},\qquad y(k)\in\mathbb{R},\qquad x(k)\in\mathbb{R}^n,
\]
e si pu\`o scrivere
\[
C=c^T,\qquad c\in\mathbb{R}^n,
\qquad
B=b,\qquad b\in\mathbb{R}^n.
\]

La funzione di trasferimento diventa quindi
\[
G(z)=c^T(zI-A)^{-1}b.
\]

Poich\'e compare l'inversa di una matrice polinomiale in \(z\),
la funzione di trasferimento \`e una \textbf{funzione razionale}:
\[
G(z)=
\frac{b_m z^m+b_{m-1}z^{m-1}+\dots+b_0}
{z^n+a_{n-1}z^{n-1}+\dots+a_1 z+a_0}.
\]

\subsection*{Osservazione}

Nel caso SISO, il denominatore coincide con il polinomio caratteristico
\[
\det(zI-A),
\]
mentre il numeratore dipende da \(A,B,C\).

\section{Equazione alle differenze associata alla funzione di trasferimento}

Sotto ipotesi di condizioni iniziali nulle, dalla relazione
\[
Y(z)=G(z)U(z)
\]
si ha
\[
\bigl(z^n+a_{n-1}z^{n-1}+\dots+a_0\bigr)Y(z)
=
\bigl(b_m z^m+b_{m-1}z^{m-1}+\dots+b_0\bigr)U(z).
\]

Formalmente, questa relazione porta a una equazione alle differenze.
Negli appunti compare una scrittura del tipo
\[
y(k+n)+a_{n-1}y(k+n-1)+\dots+a_0y(k)
=
b_m u(k+m)+b_{m-1}u(k+m-1)+\dots+b_0u(k).
\]

\subsection*{Forma causale pi\`u usata}

Per implementare il sistema come filtro digitale causale, \`e per\`o pi\`u comodo
scrivere la funzione di trasferimento in potenze di \(z^{-1}\):
\[
G(z)=
\frac{\beta_0+\beta_1 z^{-1}+\dots+\beta_m z^{-m}}
{1+\alpha_1 z^{-1}+\dots+\alpha_n z^{-n}}.
\]

In questo caso si ottiene direttamente la forma causale
\[
y(k)+\alpha_1 y(k-1)+\dots+\alpha_n y(k-n)
=
\beta_0 u(k)+\beta_1 u(k-1)+\dots+\beta_m u(k-m).
\]

Questa \`e l'equazione che definisce l'algoritmo di calcolo dell'uscita
quando l'ingresso attraversa un filtro con funzione di trasferimento \(G(z)\).

\section{Esempi di aggiunta e inversa di matrice}

Negli appunti compaiono due esempi espliciti di calcolo dell'aggiunta.

\subsection{Primo esempio: matrice \(2\times 2\)}

Sia
\[
M=
\begin{bmatrix}
3 & 6\\
7 & 9
\end{bmatrix}.
\]

Allora
\[
\det(M)=3\cdot 9-6\cdot 7=-15,
\]
e
\[
\operatorname{Adj}(M)=
\begin{bmatrix}
9 & -6\\
-7 & 3
\end{bmatrix}.
\]

Quindi
\[
M^{-1}
=
\frac{1}{-15}
\begin{bmatrix}
9 & -6\\
-7 & 3
\end{bmatrix}.
\]

\subsection{Secondo esempio: matrice \(3\times 3\)}

Sia
\[
M=
\begin{bmatrix}
5 & 2 & 0\\
6 & 6 & 1\\
0 & 3 & 2
\end{bmatrix}.
\]

Il determinante vale
\[
\det(M)=21,
\]
mentre l'aggiunta \`e
\[
\operatorname{Adj}(M)=
\begin{bmatrix}
9 & -4 & 2\\
-12 & 10 & -5\\
18 & -15 & 18
\end{bmatrix}.
\]

Quindi
\[
M^{-1}
=
\frac{1}{21}
\begin{bmatrix}
9 & -4 & 2\\
-12 & 10 & -5\\
18 & -15 & 18
\end{bmatrix}.
\]

\section{Sintesi della lezione}

I punti chiave della lezione sono:
\begin{enumerate}
    \item la trasformata z di una successione causale si definisce come
    \[
    X(z)=\sum_{k=0}^{+\infty}x(k)z^{-k};
    \]
    \item la dipendenza dal periodo di campionamento \(T\) scompare formalmente,
    ma resta implicita nella relazione \(z=e^{sT}\);
    \item la trasformata z monolatera ha propriet\`a di linearit\`a, convoluzione,
    traslazione nel tempo, valore iniziale e valore finale;
    \item esistono tavole di corrispondenza tra segnali continui causali,
    trasformate di Laplace e trasformate z delle rispettive sequenze campionate;
    \item per un sistema discreto in forma di stato
    \[
    \begin{cases}
    x(k+1)=Ax(k)+Bu(k),\\
    y(k)=Cx(k),
    \end{cases}
    \]
    si ottiene
    \[
    Y(z)=C(zI-A)^{-1}BU(z)+C(zI-A)^{-1}zx(0);
    \]
    \item la matrice di trasferimento \`e
    \[
    G(z)=C(zI-A)^{-1}B;
    \]
    \item nel caso SISO, \(G(z)\) \`e una funzione razionale;
    \item dalla funzione di trasferimento si ricava l'equazione alle differenze
    che implementa il filtro discreto;
    \item il calcolo esplicito di \(G(z)\) pu\`o essere fatto usando
    \[
    (zI-A)^{-1}=\frac{1}{\det(zI-A)}\operatorname{Adj}(zI-A).
    \]
\end{enumerate}
\chapter{Lezione 7 --- Anti-trasformata z, fratti semplici e funzione di trasferimento discreta}

\section{Richiamo: trasformata z e sistemi discreti}

Riprendiamo la definizione di trasformata z monolatera di una successione
\[
x(k), \qquad k=0,1,2,\dots
\]
data da
\[
X(z)=\mathcal{Z}\{x(k)\}=\sum_{k=0}^{+\infty} x(k)\,z^{-k}.
\]

Nel caso di un segnale continuo \(f(t)\) campionato con periodo \(T\), si scrive pi\`u rigorosamente
\[
\mathcal{Z}\{f(kT)\}=\sum_{k=0}^{+\infty} f(kT)\,z^{-k}.
\]

Negli appunti si sottolinea che la trasformata z nasce come riscrittura della trasformata di Laplace
del segnale campionato, con il cambio di variabile
\[
z=e^{sT}.
\]

\section{Richiamo: matrice e funzione di trasferimento}

Per un sistema discreto in forma di stato
\[
\begin{cases}
x(k+1)=Ax(k)+Bu(k),\\[4pt]
y(k)=Cx(k),
\end{cases}
\]
la matrice di trasferimento \`e
\[
G(z)=C(zI-A)^{-1}B.
\]

Nel caso \textbf{SISO}, si ha
\[
u(k)\in\mathbb{R}, \qquad y(k)\in\mathbb{R}, \qquad x(k)\in\mathbb{R}^n,
\]
e quindi
\[
G(z)=c^T(zI-A)^{-1}b=\frac{N(z)}{D(z)},
\]
cio\`e una funzione razionale in \(z\).

In condizioni iniziali nulle vale
\[
Y(z)=G(z)U(z).
\]

Il problema affrontato in questa lezione \`e allora il seguente:
dato \(Y(z)\), oppure pi\`u in generale data una funzione razionale \(G(z)\),
\emph{come si torna alla successione nel tempo discreto}?

\section{Questione di unicit\`a}

Negli appunti compare un'osservazione importante sul confronto tra Laplace e zeta.

Nel continuo, due funzioni che differiscono solo su un insieme di misura nulla
possono avere la stessa trasformata di Laplace, perch\'e la trasformata si basa su un integrale.

Nel discreto, invece, la successione \`e definita punto per punto:
se due successioni sono diverse, differiscono necessariamente su almeno un campione,
e quindi in condizioni standard danno trasformate z diverse.

Per questo, in ambito discreto, si ragiona di fatto con una corrispondenza biunivoca
tra successione e trasformata z, purch\'e si tenga conto anche della regione di convergenza.

\section{Metodi per l'anti-trasformata z}

Negli appunti vengono richiamati due metodi principali per ricavare la successione \(x(k)\)
a partire da \(X(z)\), quando \(X(z)\) \`e razionale:

\begin{enumerate}
    \item \textbf{fratti semplici};
    \item \textbf{lunga divisione}.
\end{enumerate}

Il primo \`e il metodo principale quando il denominatore si pu\`o fattorizzare bene
e si vogliono ricavare formule chiuse.
Il secondo \`e utile quando si vuole sviluppare \(X(z)\) come serie in potenze di \(z^{-1}\)
e leggere direttamente i coefficienti della successione.

\section{Richiamo sui fratti semplici nel dominio di Laplace}

Negli appunti si fa prima un parallelo con la trasformata di Laplace.

Se
\[
G(s)=\frac{N(s)}{D(s)}
\]
e il denominatore si fattorizza come
\[
D(s)=\prod_{i=1}^{n}(s-\alpha_i),
\]
con radici semplici, allora si pu\`o scrivere
\[
G(s)=\sum_{i=1}^{n}\frac{A_i}{s-\alpha_i}.
\]

Da qui segue immediatamente
\[
g(t)=\sum_{i=1}^{n} A_i e^{\alpha_i t}.
\]

Se invece una radice \(\alpha_i\) ha molteplicit\`a maggiore di \(1\), compaiono termini del tipo
\[
\frac{A_{i,1}}{s-\alpha_i}+\frac{A_{i,2}}{(s-\alpha_i)^2}+\dots+\frac{A_{i,\mu_i}}{(s-\alpha_i)^{\mu_i}}.
\]

I coefficienti si possono determinare:
\begin{itemize}
    \item per confronto dei coefficienti;
    \item con il metodo dei residui, nel caso di poli semplici:
    \[
    A_i=\lim_{s\to \alpha_i}(s-\alpha_i)G(s).
    \]
\end{itemize}

\section{Applicazione del metodo dei fratti semplici nel dominio \(z\)}

L'idea \`e del tutto analoga nel dominio z.

Se
\[
G(z)=\frac{N(z)}{D(z)},
\]
si vorrebbe decomporre \(G(z)\) in termini elementari di cui si conosce gi\`a
l'anti-trasformata z.

\subsection*{Osservazione pratica}

Negli appunti si osserva che spesso, scrivendo \(G(z)\) in funzione di \(z\),
il numeratore e il denominatore hanno lo stesso grado.
In questi casi, per applicare comodamente i fratti semplici,
si lavora su
\[
\frac{G(z)}{z}=\frac{N(z)}{z\,D(z)},
\]
cio\`e si aggiunge un polo nell'origine.

Questo rende pi\`u naturale la decomposizione in termini elementari
ed \`e coerente con il fatto che molte tavole di anti-trasformata z
sono espresse in forme del tipo
\[
\frac{z}{z-a}, \qquad
\frac{z\sin\theta}{z^2-2\cos\theta\,z+1},
\qquad
\frac{z(z-\cos\theta)}{z^2-2\cos\theta\,z+1}.
\]

\section{Tavola di riferimento per l'anti-trasformata z}

Le forme elementari da tenere a mente sono:

\[
\mathcal{Z}\{\delta(k)\}=1,
\]
\[
\mathcal{Z}\{a^k 1(k)\}=\frac{z}{z-a},
\qquad |z|>|a|,
\]
\[
\mathcal{Z}\{\sin(k\theta)\,1(k)\}
=
\frac{z\sin\theta}{z^2-2\cos\theta\,z+1},
\]
\[
\mathcal{Z}\{\cos(k\theta)\,1(k)\}
=
\frac{z(z-\cos\theta)}{z^2-2\cos\theta\,z+1}.
\]

Queste forme sono quelle che permettono di riconoscere rapidamente
i contributi elementari nella decomposizione in fratti semplici.

\section{Esempio di decomposizione nel dominio \(z\)}

Negli appunti compare l'esempio
\[
G_1(z)=\frac{z+1}{\left(z-\frac13\right)\left(z^2-\frac14 z+1\right)}.
\]

I poli sono:
\[
z_1=\frac13,
\qquad
z_{2,3}=\frac18 \pm j\,\frac{3\sqrt7}{8}.
\]

Poich\'e
\[
z^2-\frac14 z+1
=
z^2-2\cos\theta\,z+1,
\]
si ricava
\[
2\cos\theta=\frac14
\qquad\Longrightarrow\qquad
\cos\theta=\frac18,
\qquad
\theta=\arccos\!\left(\frac18\right).
\]

Per applicare il metodo, si considera
\[
\frac{G_1(z)}{z}
=
\frac{z+1}{z\left(z-\frac13\right)\left(z^2-\frac14 z+1\right)}.
\]

Una forma di decomposizione coerente con le tavole note \`e
\[
\frac{G_1(z)}{z}
=
\frac{A}{z}
+
\frac{B}{z-\frac13}
+
C\,\frac{z\sin\theta}{z^2-\frac14 z+1}
+
D\,\frac{z(z-\cos\theta)}{z^2-\frac14 z+1}.
\]

I coefficienti \(A,B,C,D\) si determinano per uguaglianza tra i due membri,
oppure combinando confronto dei coefficienti e uso delle forme note.

Una volta trovati i coefficienti, l'anti-trasformata si legge termine a termine come combinazione di:
\begin{itemize}
    \item un termine impulsivo;
    \item un'esponenziale discreta \(\left(\frac13\right)^k\);
    \item un termine sinusoidale \(\sin(k\theta)\);
    \item un termine cosinusoidale \(\cos(k\theta)\).
\end{itemize}

\subsection*{Osservazione}

Negli appunti \`e evidenziato che i vari metodi possono anche essere \emph{combinati}:
per esempio si possono isolare alcuni termini con i fratti semplici
e ricavare altri sviluppando in serie.

\section{Metodo della lunga divisione}

Il secondo metodo consiste nello sviluppare \(X(z)\) in serie di potenze di \(z^{-1}\):
\[
X(z)=x(0)+x(1)z^{-1}+x(2)z^{-2}+\dots
\]

Se si riesce a scrivere la funzione razionale in questa forma,
allora i coefficienti della serie sono direttamente i valori della successione.

Questo metodo \`e particolarmente utile:
\begin{itemize}
    \item quando si vogliono calcolare i primi campioni;
    \item quando la fattorizzazione del denominatore \`e scomoda;
    \item quando interessa una realizzazione ricorsiva.
\end{itemize}

\section{Collegamento con la funzione di trasferimento discreta}

Nel caso SISO, la funzione di trasferimento
\[
G(z)=\frac{N(z)}{D(z)}
\]
descrive il legame tra ingresso e uscita:
\[
Y(z)=G(z)U(z).
\]

Per tornare al dominio discreto si possono seguire due strade:
\begin{enumerate}
    \item anti-trasformare direttamente \(Y(z)\) con fratti semplici o lunga divisione;
    \item riscrivere \(G(z)\) in forma polinomiale in \(z^{-1}\) e ottenere
    l'equazione alle differenze.
\end{enumerate}

La seconda strada \`e quella pi\`u utile quando si vuole implementare il sistema,
mentre la prima \`e quella pi\`u utile quando si vuole trovare una formula esplicita
per la risposta.

\section{Osservazione sul continuo e sul discreto}

Negli appunti compare un'ultima osservazione concettuale:
nel continuo esistono formalmente funzioni diverse che differiscono solo su insiemi
di misura nulla e che quindi hanno la stessa trasformata integrale.

Nel discreto questo problema, dal punto di vista pratico, non si presenta:
i segnali reali sono successioni di valori campionati, quindi l'informazione
\`e contenuta esattamente nei campioni.

Per questo, nell'uso ingegneristico della trasformata z,
si ragiona come se ci fosse una corrispondenza biunivoca tra sequenza e trasformata.

\section{Sintesi della lezione}

I punti principali della lezione sono:
\begin{enumerate}
    \item richiamo della matrice di trasferimento discreta
    \[
    G(z)=C(zI-A)^{-1}B;
    \]
    \item nel caso SISO, \(G(z)\) \`e una funzione razionale
    \[
    G(z)=\frac{N(z)}{D(z)};
    \]
    \item il problema dell'anti-trasformata z consiste nel tornare da \(X(z)\)
    alla successione \(x(k)\);
    \item i due metodi principali sono:
    \begin{itemize}
        \item fratti semplici;
        \item lunga divisione;
    \end{itemize}
    \item il metodo dei fratti semplici nel dominio \(z\) ricalca quello della Laplace;
    \item spesso conviene lavorare su
    \[
    \frac{G(z)}{z}
    \]
    per ottenere una forma pi\`u comoda da decomporre;
    \item i termini elementari dell'anti-trasformata sono esponenziali discreti,
    seni e coseni discreti, oltre agli impulsi;
    \item l'esempio svolto mostra come trattare un polo reale e una coppia di poli
    complessi coniugati nel piano \(z\).
\end{enumerate}
\chapter{Lezione 8 --- Anti-trasformata z, implementazione dei filtri digitali, trasformata bilineare di Tustin e pre-warping}

\section{Richiamo: tavola essenziale per l'anti-trasformata z}

Negli appunti viene ripresa una tabella di coppie utili tra segnale continuo causale,
successione campionata e trasformata z.

Ponendo
\[
\varphi=e^{-aT},
\]
si hanno le seguenti corrispondenze fondamentali:

\[
\mathcal{Z}\{1\}=\frac{z}{z-1},
\]
\[
\mathcal{Z}\{kT\}=\frac{Tz}{(z-1)^2},
\]
\[
\mathcal{Z}\left\{\frac{1}{2}k^2T^2\right\}
=
\frac{T^2}{2}\frac{z(z+1)}{(z-1)^3},
\]
\[
\mathcal{Z}\{\varphi^k\}=\frac{z}{z-\varphi},
\]
\[
\mathcal{Z}\{kT\,\varphi^k\}
=
\frac{T\varphi z}{(z-\varphi)^2},
\]
\[
\mathcal{Z}\{1-\varphi^k\}
=
\frac{z(1-\varphi)}{(z-1)(z-\varphi)},
\]
\[
\mathcal{Z}\{\sin(\omega kT)\}
=
\frac{z\sin(\omega T)}{z^2-2\cos(\omega T)\,z+1},
\]
\[
\mathcal{Z}\{\cos(\omega kT)\}
=
\frac{z\bigl(z-\cos(\omega T)\bigr)}{z^2-2\cos(\omega T)\,z+1},
\]
\[
\mathcal{Z}\{\varphi^k\sin(\omega kT)\}
=
\frac{z\varphi\sin(\omega T)}{z^2-2\varphi\cos(\omega T)\,z+\varphi^2},
\]
\[
\mathcal{Z}\{\varphi^k\cos(\omega kT)\}
=
\frac{z\bigl(z-\varphi\cos(\omega T)\bigr)}{z^2-2\varphi\cos(\omega T)\,z+\varphi^2}.
\]

Questa tabella \`e il riferimento operativo per tornare da una funzione razionale \(G(z)\)
alla corrispondente successione nel tempo discreto.

\section{Anti-trasformata z con il metodo dei fratti semplici}

Negli appunti si osserva che, diversamente da quanto accade in Laplace,
nelle forme tabulate della trasformata z compare quasi sempre un fattore \(z\) al numeratore.
Per questo, data una funzione razionale
\[
G(z)=\frac{N(z)}{D(z)},
\]
per applicare il metodo dei fratti semplici conviene spesso considerare
\[
\frac{G(z)}{z}.
\]

L'idea operativa \`e:
\[
G(z)
\quad\longrightarrow\quad
\frac{G(z)}{z}
\quad\longrightarrow\quad
\text{espansione in fratti semplici}
\quad\longrightarrow\quad
\text{calcolo dei residui}
\quad\longrightarrow\quad
\text{moltiplico per } z
\quad\longrightarrow\quad
\text{anti-trasformo via tabella}.
\]

In questo modo si ottiene una soluzione analitica della successione.

\subsection*{Osservazione}

Il procedimento \`e del tutto analogo a quello usato nel dominio di Laplace:
si fattorizza il denominatore, si espande in termini elementari e poi si riconosce
l'anti-trasformata di ciascun termine.

\section{Metodo non analitico: lunga divisione}

Negli appunti viene richiamato anche un secondo metodo, pi\`u algoritmico:
la \textbf{lunga divisione} tra numeratore e denominatore.

Se
\[
G(z)=\frac{N(z)}{D(z)},
\]
si sviluppa la funzione come serie in potenze di \(z^{-1}\):
\[
G(z)=g(0)+g(1)z^{-1}+g(2)z^{-2}+\dots
\]

I coefficienti letti nella serie sono proprio i valori della successione.

\subsection*{Interpretazione}

Ogni moltiplicazione per \(z^{-1}\) corrisponde a un ritardo di un campione,
mentre ogni moltiplicazione per \(z\) corrisponde formalmente a un anticipo.

Il metodo della lunga divisione \`e particolarmente utile quando:
\begin{itemize}
    \item si vogliono i primi campioni della risposta;
    \item la fattorizzazione del denominatore \`e scomoda;
    \item interessa un approccio numerico pi\`u che una formula chiusa.
\end{itemize}

\section{Dalla funzione di trasferimento all'equazione alle differenze}

Consideriamo una funzione di trasferimento discreta del tipo
\[
G(z)=\frac{n(z)}{d(z)}
=
\frac{b_m z^m+\dots+b_0}{z^n+a_{n-1}z^{n-1}+\dots+a_0}.
\]

Assumendo la relazione ingresso--uscita
\[
Y(z)=G(z)U(z),
\]
moltiplicando per il denominatore si ottiene
\[
d(z)Y(z)=n(z)U(z).
\]

Anti-trasformando, si ricava una equazione alle differenze del tipo
\[
y(k)+a_{n-1}y(k-1)+\dots+a_0 y(k-n)
=
b_m u(k)+b_{m-1}u(k-1)+\dots+b_0 u(k-m).
\]

Quindi si pu\`o risolvere esplicitamente rispetto a \(y(k)\):
\[
y(k)=
-a_{n-1}y(k-1)-\dots-a_0y(k-n)
+b_m u(k)+\dots+b_0u(k-m).
\]

Questa \`e l'equazione che permette di implementare il filtro nel tempo discreto.

\section{Implementazione digitale del filtro}

Per implementare l'equazione alle differenze su una macchina digitale serve:
\begin{itemize}
    \item un \textbf{buffer di memoria}, per memorizzare i campioni passati di ingresso e uscita;
    \item una \textbf{ALU} (\emph{Arithmetic Logic Unit}), che esegue somme e moltiplicazioni;
    \item una logica di aggiornamento dei campioni nel buffer.
\end{itemize}

Se l'ordine del filtro \`e \(n\), in una realizzazione tipica bisogna conservare:
\begin{itemize}
    \item \(n\) campioni passati dell'uscita;
    \item \(n+1\) campioni dell'ingresso.
\end{itemize}

Quindi il buffer totale contiene, a meno di convenzioni di indicizzazione,
circa \(2n+1\) campioni.

\section{Filtri IIR e filtri FIR}

\subsection{Filtri IIR}

Un filtro descritto da una equazione alle differenze con dipendenza
anche dai valori passati dell'uscita \`e un filtro \textbf{ricorsivo}.

Questi filtri si chiamano
\[
\text{IIR} \quad \text{(\emph{Infinite Impulse Response})}.
\]

Il nome deriva dal fatto che, in generale, la risposta all'impulso discreto
non si annulla in un numero finito di passi.

\subsection{Filtri FIR}

Se invece la funzione di trasferimento ha tutti i poli nell'origine, cio\`e pu\`o essere scritta come
\[
G(z)=\frac{m(z)}{z^n},
\]
allora il filtro non \`e ricorsivo e l'uscita dipende solo da valori presenti e passati dell'ingresso:
\[
y(k)=b_m u(k)+b_{m-1}u(k-1)+\dots+b_0 u(k-m).
\]

Questi filtri si chiamano
\[
\text{FIR} \quad \text{(\emph{Finite Impulse Response})}.
\]

Se l'ingresso \(u(k)\) \`e un impulso discreto, la risposta di un filtro FIR va a zero
in un numero finito di passi.

\section{Specifiche di un filtro passa-basso}

Negli appunti si introduce il problema progettuale:
si vogliono assegnare le specifiche di un filtro passa-basso nel continuo
e poi costruire un filtro digitale equivalente.

Un modello analogico elementare \`e
\[
G(s)=\frac{\omega_b}{s+\omega_b},
\]
dove \(\omega_b\) \`e la pulsazione di banda.

Si tratta di un filtro del primo ordine, con pendenza asintotica
\[
-20\ \text{dB/decade}
\]
oltre la frequenza di taglio.

\section{Problema: come passare da \(s\) a \(z\)?}

Se si vuole lavorare non pi\`u su segnali continui ma su segnali campionati,
bisogna sostituire il filtro analogico con un filtro discreto \(G(z)\).

La relazione esatta tra \(s\) e \(z\),
\[
z=e^{sT},
\]
\`e per\`o trascendente, quindi non porta in generale a una funzione razionale semplice in \(z\).

Per ottenere una funzione di trasferimento digitale razionale
si introducono trasformazioni approssimate da \(s\) a \(z\).

\section{Idea generale: approssimazioni razionali}

Negli appunti si richiama che la trasformazione non pu\`o essere esatta,
ma si pu\`o cercare una buona approssimazione.

Una prima idea \`e quella di usare lo sviluppo di Taylor troncato.
Una scelta pi\`u efficace \`e usare un approssimante di Pad\'e,
cio\`e una funzione razionale che approssima una funzione trascendente.

Qui si sceglie un approssimante di Pad\'e del primo ordine per \(e^{sT}\).

\section{Approssimazione di Pad\'e del primo ordine}

Scriviamo
\[
e^{sT}= \frac{e^{\frac{sT}{2}}}{e^{-\frac{sT}{2}}}.
\]

Approssimando numeratore e denominatore al primo ordine di Taylor:
\[
e^{\frac{sT}{2}} \approx 1+\frac{sT}{2},
\qquad
e^{-\frac{sT}{2}} \approx 1-\frac{sT}{2}.
\]

Si ottiene quindi
\[
e^{sT}\approx \frac{1+\frac{sT}{2}}{1-\frac{sT}{2}}.
\]

Poich\'e \(z=e^{sT}\), segue l'approssimazione
\[
z\approx \frac{1+\frac{sT}{2}}{1-\frac{sT}{2}}.
\]

Risolvendo rispetto a \(s\):
\[
\left(1-\frac{sT}{2}\right)z = 1+\frac{sT}{2},
\]
\[
z-1=\frac{sT}{2}(z+1),
\]
\[
s=\frac{2}{T}\frac{z-1}{z+1}.
\]

Questa \`e la \textbf{trasformata di Tustin}, detta anche \textbf{trasformata bilineare}.

\section{Trasformata di Tustin o bilineare}

La sostituzione fondamentale \`e
\[
s=\frac{2}{T}\frac{z-1}{z+1}.
\]

Essa permette di trasformare una funzione di trasferimento continua \(G(s)\)
in una funzione di trasferimento discreta \(G(z)\) razionale.

\subsection*{Osservazione}

La trasformazione \`e semplice e molto usata in pratica,
ma resta comunque un'approssimazione del legame esatto tra \(s\) e \(z\).

\section{Esempio: discretizzazione di un passa-basso del primo ordine}

Consideriamo il filtro analogico
\[
G(s)=\frac{\omega_b}{s+\omega_b}.
\]

Sostituendo \(s=\frac{2}{T}\frac{z-1}{z+1}\), si ottiene
\[
G(z)=\frac{\omega_b}{\frac{2}{T}\frac{z-1}{z+1}+\omega_b}.
\]

Moltiplicando numeratore e denominatore per \(z+1\):
\[
G(z)=
\frac{\omega_b(z+1)}
{\frac{2}{T}(z-1)+\omega_b(z+1)}.
\]

Moltiplicando ancora per \(T\):
\[
G(z)=
\frac{T\omega_b(z+1)}
{2(z-1)+T\omega_b(z+1)}.
\]

Sviluppando il denominatore:
\[
G(z)=
\frac{T\omega_b(z+1)}
{(2+T\omega_b)z-2+T\omega_b}.
\]

Questa \`e una funzione razionale in \(z\), quindi pu\`o essere implementata
come filtro digitale tramite le tecniche gi\`a viste.

\section{Dove vanno i punti del piano \(s\) con la trasformata bilineare?}

La lezione si chiede anche quale sia l'effetto geometrico della trasformazione di Tustin
sui punti del piano \(s\).

Prendendo un punto sull'asse immaginario,
\[
s=j\omega,
\]
si ha
\[
z=\frac{1+j\omega T/2}{1-j\omega T/2}.
\]

Calcoliamo il modulo:
\[
|z|^2
=
\frac{\left(1+j\omega T/2\right)\left(1-j\omega T/2\right)}
{\left(1-j\omega T/2\right)\left(1+j\omega T/2\right)}
=1.
\]

Quindi
\[
|z|=1.
\]

Ne segue che l'asse immaginario del piano \(s\) viene mappato sul cerchio unitario del piano \(z\).

Inoltre:
\begin{itemize}
    \item il semipiano sinistro di \(s\) viene mappato all'interno del cerchio unitario;
    \item il semipiano destro di \(s\) viene mappato all'esterno del cerchio unitario.
\end{itemize}

Questa propriet\`a \`e molto importante perch\'e preserva il criterio di stabilit\`a:
un sistema stabile nel continuo viene trasformato in un sistema stabile nel discreto.

\section{Distorsione in frequenza: frequency warping}

La trasformata di Tustin non preserva linearmente la frequenza.
Per capire cosa succede alla risposta armonica,
si valuta il filtro discreto sul cerchio unitario:
\[
z=e^{j\omega_d T}.
\]

Allora
\[
G_d(e^{j\omega_d T})
=
G_a\!\left(\frac{2}{T}\frac{e^{j\omega_d T}-1}{e^{j\omega_d T}+1}\right).
\]

Semplificando il rapporto:
\[
\frac{e^{j\omega_d T}-1}{e^{j\omega_d T}+1}
=
\frac{e^{j\omega_d T/2}\left(e^{j\omega_d T/2}-e^{-j\omega_d T/2}\right)}
{e^{j\omega_d T/2}\left(e^{j\omega_d T/2}+e^{-j\omega_d T/2}\right)}
=
\frac{2j\sin(\omega_d T/2)}{2\cos(\omega_d T/2)}
=
j\tan\!\left(\frac{\omega_d T}{2}\right).
\]

Quindi
\[
G_d(e^{j\omega_d T})
=
G_a\!\left(
j\,\frac{2}{T}\tan\!\left(\frac{\omega_d T}{2}\right)
\right).
\]

Questo mostra che le pulsazioni analogiche e discrete sono legate da
\[
\omega_a=\frac{2}{T}\tan\!\left(\frac{\omega_d T}{2}\right),
\]
oppure inversamente
\[
\omega_d=\frac{2}{T}\arctan\!\left(\frac{\omega_a T}{2}\right).
\]

\subsection*{Interpretazione}

La risposta armonica viene distorta poco alle basse frequenze
e molto di pi\`u alle alte frequenze.

Per questo si parla di \textbf{warping} della frequenza:
la scala delle pulsazioni nel continuo non viene trasportata linearmente
nella scala discreta.

\section{Pre-warping}

Spesso interessa che il filtro discreto riproduca in modo esatto
il filtro analogico almeno a una specifica pulsazione \(\omega_0\),
ad esempio proprio alla frequenza di taglio.

Per ottenere questo, si modifica la trasformazione bilineare introducendo una costante \(k\):
\[
s=k\,\frac{z-1}{z+1}.
\]

Si sceglie \(k\) imponendo che la corrispondenza sia esatta alla pulsazione desiderata \(\omega_0\):
\[
\omega_0 = k\,\tan\!\left(\frac{\omega_0 T}{2}\right),
\]
da cui
\[
k=\frac{\omega_0}{\tan\!\left(\frac{\omega_0 T}{2}\right)}.
\]

Questa scelta prende il nome di \textbf{pre-warping}.

\subsection*{Significato}

Con il pre-warping si forza la trasformazione in modo che
la risposta armonica del filtro discreto coincida con quella del filtro analogico
alla frequenza \(\omega_0\) di progetto.

In altre parole:
\begin{itemize}
    \item senza pre-warping, la Tustin \`e semplice ma distorce la scala in frequenza;
    \item con pre-warping, si corregge localmente la distorsione nel punto di maggiore interesse.
\end{itemize}

\section{Sintesi della lezione}

I punti chiave della lezione sono:
\begin{enumerate}
    \item richiamo dei metodi per anti-trasformare una funzione razionale in \(z\):
    \begin{itemize}
        \item fratti semplici;
        \item lunga divisione;
    \end{itemize}
    \item da \(Y(z)=G(z)U(z)\) si ricava l'equazione alle differenze che implementa il filtro;
    \item i filtri ricorsivi sono filtri IIR, mentre quelli non ricorsivi con poli nell'origine
    sono filtri FIR;
    \item per discretizzare un filtro analogico si cerca una trasformazione razionale da \(s\) a \(z\);
    \item usando l'approssimante di Pad\'e del primo ordine si ottiene la trasformata di Tustin:
    \[
    s=\frac{2}{T}\frac{z-1}{z+1};
    \]
    \item applicando Tustin a
    \[
    G(s)=\frac{\omega_b}{s+\omega_b}
    \]
    si ottiene un filtro discreto razionale \(G(z)\);
    \item la trasformata bilineare manda l'asse immaginario di \(s\) sul cerchio unitario di \(z\);
    \item la corrispondenza tra frequenze \`e
    \[
    \omega_a=\frac{2}{T}\tan\!\left(\frac{\omega_d T}{2}\right),
    \]
    quindi la scala in frequenza \`e deformata;
    \item con il pre-warping si usa
    \[
    s=k\frac{z-1}{z+1},
    \qquad
    k=\frac{\omega_0}{\tan(\omega_0 T/2)},
    \]
    per garantire corrispondenza esatta alla frequenza \(\omega_0\).
\end{enumerate}
\chapter{Lezione 9 --- Pre-warping, trasformata di Tustin e spettro discreto tramite FFT}

\section{Dal regolatore continuo \(R(s)\) al regolatore discreto \(\tilde R(z)\)}

L'obiettivo della lezione è capire come passare da un regolatore o da un filtro
progettato nel continuo a un corrispondente oggetto discreto.

Nel continuo si lavora con
\[
R(s),
\]
mentre nel discreto si vuole ottenere una funzione di trasferimento
\[
\tilde R(z).
\]

Il legame esatto tra le due descrizioni è dato dalla relazione
\[
z=e^{sT},
\qquad\text{equivalentemente}\qquad
s=\frac{1}{T}\ln z.
\]

Tuttavia questa relazione è trascendente e non è comoda da usare direttamente
per ottenere una funzione razionale in \(z\).
Per questo si usa una sostituzione approssimata, ma razionale, cioè la
\textbf{trasformata di Tustin}.

\section{Trasformata di Tustin}

La trasformata di Tustin, o trasformata bilineare, si scrive come
\[
s=\frac{2}{T}\frac{z-1}{z+1}.
\]

In forma inversa, si può anche scrivere
\[
z=\frac{\frac{2}{T}+s}{\frac{2}{T}-s}.
\]

Questa trasformazione ha due proprietà molto importanti:
\begin{itemize}
    \item manda il semipiano sinistro del piano \(s\) all'interno del cerchio unitario;
    \item manda l'asse immaginario del piano \(s\) sul cerchio unitario del piano \(z\).
\end{itemize}

Quindi:
\begin{itemize}
    \item poli stabili nel continuo vengono mappati in poli stabili nel discreto;
    \item la stabilità viene preservata.
\end{itemize}

\section{Risposta armonica e problema della distorsione}

Se si discretizza un filtro analogico con Tustin, la risposta armonica del filtro discreto
non coincide esattamente con quella del filtro continuo: la scala delle frequenze viene distorta.

Per un filtro passa-basso, in genere interessa preservare bene il comportamento
in una frequenza significativa, tipicamente la frequenza di taglio.

Ad esempio, se nel continuo si ha un filtro del primo ordine,
alla pulsazione di taglio \(\omega_b=1/\tau\) si ha tipicamente:
\begin{itemize}
    \item modulo pari a \(-3\) dB rispetto al valore a bassa frequenza;
    \item fase pari a \(-45^\circ\).
\end{itemize}

L'idea è allora imporre che queste caratteristiche siano preservate,
non ovunque, ma almeno in una pulsazione di interesse.

\section{Pre-warping}

Per correggere la distorsione della scala in frequenza si introduce una costante \(k\)
nella trasformata bilineare:
\[
s=k\frac{z-1}{z+1}.
\]

Si sceglie \(k\) in modo che la corrispondenza tra filtro analogico e filtro discreto
sia esatta in una pulsazione \(\omega_0\) prefissata.

Il valore corretto di \(k\) è
\[
k=\frac{\omega_0}{\tan\!\left(\frac{\omega_0 T}{2}\right)}.
\]

Questa scelta prende il nome di \textbf{pre-warping} alla pulsazione \(\omega_0\).

\subsection*{Condizione di progetto}

Con questa scelta si forza l'uguaglianza
\[
R(j\omega_0)=\tilde R(e^{j\omega_0 T})
\]
nel senso della risposta armonica: il filtro discreto riproduce esattamente
il comportamento del filtro continuo alla frequenza \(\omega_0\).

\subsection*{Caso limite a bassa frequenza}

Se \(\omega_0\to 0\), si ha
\[
\tan\!\left(\frac{\omega_0 T}{2}\right)\sim \frac{\omega_0 T}{2},
\]
e quindi
\[
k=\frac{\omega_0}{\tan(\omega_0 T/2)}
\longrightarrow
\frac{\omega_0}{\omega_0 T/2}
=
\frac{2}{T}.
\]

Quindi la trasformata di Tustin standard è il caso limite del pre-warping
quando si vuole privilegiare la bassa frequenza.

\section{Più frequenze di interesse}

Negli appunti compare un'osservazione importante:
se ci sono più frequenze da preservare, non è possibile renderne esatta la corrispondenza
per tutte contemporaneamente con una sola trasformazione di Tustin pre-warpata.

In pratica:
\begin{itemize}
    \item si sceglie una frequenza prioritaria;
    \item oppure si lavora per tentativi e compromessi;
    \item spesso si privilegia la frequenza di taglio di un passa-basso.
\end{itemize}

Questo è coerente con il fatto che la trasformazione resta un'approssimazione,
non una equivalenza esatta tra continuo e discreto.

\section{Campionamento di una sequenza lunga e calcolo dello spettro}

Nella seconda parte della lezione si passa al problema pratico:
se si campiona una lunga sequenza di dati e si vuole ricavare lo spettro,
si calcola la trasformata di Fourier discreta della sequenza campionata.

Se il segnale continuo \(x(t)\) viene campionato con periodo \(T\), si ottiene la successione
\[
\tilde x(kT),\qquad k=0,1,\dots,N-1.
\]

A partire da questi \(N\) campioni si costruisce il vettore dei campioni nel tempo
e si calcola la trasformata discreta di Fourier:
\[
\tilde X(h)=\sum_{k=0}^{N-1}\tilde x(kT)\,e^{-j\frac{2\pi}{N}hk},
\qquad
h=0,1,\dots,N-1.
\]

Il risultato \(\tilde X(h)\) è un vettore di numeri complessi che rappresenta modulo e fase
dello spettro campionato.

\section{FFT}

Il calcolo diretto della DFT richiederebbe un numero di operazioni dell'ordine di
\[
N^2.
\]

Per rendere il calcolo efficiente si usa la \textbf{FFT} (\emph{Fast Fourier Transform}),
cioè un algoritmo che calcola la DFT con costo molto minore, tipicamente dell'ordine di
\[
N\log N.
\]

Per questo motivo, quando si devono elaborare sequenze lunghe,
si preferisce usare la FFT invece della formula diretta della DFT.

\subsection*{Osservazione pratica}

È molto conveniente scegliere un numero di campioni \(N\) che sia una potenza di \(2\),
perché molti algoritmi FFT lavorano in modo particolarmente efficiente in questo caso.

\section{Asse delle frequenze restituito dalla FFT}

Un punto essenziale della lezione è capire che i valori \(\tilde X(h)\)
non sono associati a frequenze arbitrarie, ma a frequenze equispaziate.

La pulsazione corrispondente al bin \(h\) è
\[
\omega_h = h\,\frac{2\pi}{NT}.
\]

Quindi:
\[
\Delta\omega = \frac{2\pi}{NT}
\]
è il passo tra due campioni consecutivi in frequenza.

In particolare:
\[
h=0 \quad \Longrightarrow \quad \omega_0=0,
\]
\[
h=\frac{N}{2} \quad \Longrightarrow \quad \omega_{N/2}=\frac{\pi}{T}.
\]

Quindi il bin \(h=N/2\) corrisponde alla pulsazione di Nyquist.

\subsection*{Osservazione sull'intervallo utile}

Per segnali reali, lo spettro contiene informazione ridondante nella seconda metà,
per simmetria complessa coniugata.
Per questo, in molti casi pratici, l'intervallo utile è solo
\[
0\le \omega \le \frac{\pi}{T},
\]
cioè da \(0\) fino alla metà della pulsazione di campionamento.

In questo intervallo si hanno, a seconda della convenzione di indicizzazione,
circa \(N/2\) oppure \(N/2+1\) campioni utili.

\section{Interpretazione MATLAB}

Negli appunti si osserva che ambienti come MATLAB, quando si usa \texttt{fft},
restituiscono il vettore dei bin senza costruire automaticamente la scala fisica delle frequenze.

Di fatto, se non si imposta esplicitamente il periodo di campionamento \(T\),
la rappresentazione sembra fatta come se \(T=1\).

Quindi, per riportare correttamente l'asse delle pulsazioni,
bisogna costruire a mano il vettore
\[
\omega_h = h\,\frac{2\pi}{NT}.
\]

\section{Scala lineare e scala logaritmica}

Negli appunti viene poi discusso il problema della rappresentazione grafica dello spettro.

I campioni in frequenza prodotti dalla FFT sono \textbf{equispaziati in scala lineare},
cioè il passo tra una frequenza e la successiva è costante:
\[
\Delta\omega = \frac{2\pi}{NT}.
\]

Se però lo spettro viene riportato su asse logaritmico, i punti non appaiono più equidistanti.

In particolare:
\begin{itemize}
    \item alle basse frequenze i punti si distinguono bene;
    \item alle alte frequenze i punti si affollano e si distinguono sempre meno.
\end{itemize}

Questo è un fatto puramente geometrico dovuto al cambio di scala,
non a una diversa densità dei campioni della FFT.

\section{Limite inferiore della griglia spettrale}

La minima pulsazione positiva rappresentabile con \(N\) campioni e periodo \(T\) è
\[
\omega_{\min}=\frac{2\pi}{NT}.
\]

Questa è la frequenza del primo bin non nullo.
Per questo, se si vuole vedere bene il comportamento a bassissima frequenza,
bisogna aumentare \(N\), cioè osservare il segnale per un tempo più lungo.

\section{Normalizzazione dell'ampiezza}

Negli appunti compare un'ultima osservazione importante:
nel passaggio dal dominio del tempo alla trasformata discreta di Fourier
possono comparire discrepanze nei valori delle ampiezze se non si usa la corretta normalizzazione.

I coefficienti di scala dipendono dalla convenzione scelta:
\begin{itemize}
    \item definizione di trasformata continua;
    \item definizione di DFT/FFT;
    \item spettro monolatero o bilatero;
    \item ampiezza di picco oppure valore efficace (RMS).
\end{itemize}

Per questo, in applicazioni pratiche, possono comparire fattori del tipo:
\[
\frac{1}{N},\qquad
T,\qquad
2,\qquad
\frac{1}{2\pi},\qquad
\sqrt{2},
\]
o combinazioni di essi, a seconda di ciò che si vuole rappresentare.

\subsection*{Interpretazione}

Se si vuole ottenere dallo spettro FFT un'ampiezza fisicamente coerente con il segnale originale,
bisogna sempre specificare:
\begin{itemize}
    \item quale trasformata si sta usando;
    \item come si sta costruendo l'asse delle frequenze;
    \item quale normalizzazione si sta adottando.
\end{itemize}

\section{Modulo in dB}

Quando si vuole rappresentare il modulo dello spettro in forma logaritmica,
si usa tipicamente
\[
20\log_{10}|\tilde X(h)|.
\]

In questo modo si ottiene il diagramma del modulo in dB,
analogo a quello usato nei diagrammi di Bode.

\section{Sintesi della lezione}

I punti principali della lezione sono:
\begin{enumerate}
    \item il passaggio \(R(s)\mapsto \tilde R(z)\) si fa in pratica tramite la trasformata di Tustin;
    \item la trasformata di Tustin standard è
    \[
    s=\frac{2}{T}\frac{z-1}{z+1};
    \]
    \item con il pre-warping si usa
    \[
    s=k\frac{z-1}{z+1},
    \qquad
    k=\frac{\omega_0}{\tan(\omega_0 T/2)},
    \]
    per preservare la risposta armonica alla frequenza \(\omega_0\);
    \item al limite \(\omega_0\to 0\) si ritrova
    \[
    k\to \frac{2}{T};
    \]
    \item una sequenza di \(N\) campioni può essere trasformata nel dominio della frequenza
    tramite DFT o, più efficientemente, tramite FFT;
    \item il costo computazionale della FFT è molto più basso di quello della DFT diretta;
    \item i bin spettrali sono equispaziati in scala lineare con passo
    \[
    \Delta\omega=\frac{2\pi}{NT};
    \]
    \item la pulsazione del bin \(h\) è
    \[
    \omega_h=h\frac{2\pi}{NT};
    \]
    \item per segnali reali, l'intervallo utile è
    \[
    0\le \omega \le \frac{\pi}{T};
    \]
    \item su asse logaritmico i punti FFT risultano sempre più fitti alle alte frequenze;
    \item per rappresentare correttamente le ampiezze occorre usare la normalizzazione coerente
    con la convenzione adottata.
\end{enumerate}
\chapter{Lezione 10 --- STFT, regolatore digitale, convertitore D/A e tenuta di ordine zero}

\section{Richiamo: comportamento in frequenza della tenuta di ordine zero}

Nella parte iniziale della lezione si riprende la funzione di trasferimento della
tenuta di ordine zero, che era stata introdotta come modello del convertitore
digitale-analogico.

Se la risposta impulsiva della tenuta è un rettangolo di altezza \(1\) e durata \(T\),
allora
\[
h(t)=1(t)-1(t-T),
\]
e la sua trasformata di Laplace è
\[
F(s)=\mathcal{L}\{h(t)\}
=
\frac{1-e^{-sT}}{s}.
\]

Valutando la risposta armonica:
\[
F(j\omega)=\frac{1-e^{-j\omega T}}{j\omega}.
\]

Moltiplicando e dividendo opportunamente, si ottiene la forma standard
\[
F(j\omega)
=
e^{-j\omega T/2}\,\frac{2j\sin(\omega T/2)}{j\omega}
=
T\,e^{-j\omega T/2}\,\frac{\sin(\omega T/2)}{\omega T/2}.
\]

Questa espressione mette in evidenza due fattori distinti:
\begin{itemize}
    \item un termine di ritardo puro
    \[
    e^{-j\omega T/2},
    \]
    \item un termine reale di tipo sinc
    \[
    \frac{\sin(\omega T/2)}{\omega T/2}.
    \]
\end{itemize}

\section{Modulo della tenuta di ordine zero}

Dal risultato precedente segue
\[
|F(j\omega)|=
T\left|\frac{\sin(\omega T/2)}{\omega T/2}\right|.
\]

Quindi:
\begin{itemize}
    \item per \(\omega\to 0\), il modulo tende a \(T\);
    \item ci sono zeri quando
    \[
    \sin(\omega T/2)=0
    \quad\Longrightarrow\quad
    \omega=\frac{2k\pi}{T},\qquad k\in\mathbb{Z}\setminus\{0\};
    \]
    \item il modulo ha quindi l'andamento tipico di una sinc.
\end{itemize}

Se si considera la versione normalizzata dividendo per \(T\), si ottiene
\[
\left|\frac{F(j\omega)}{T}\right|
=
\left|\frac{\sin(\omega T/2)}{\omega T/2}\right|.
\]

Alla pulsazione di Nyquist
\[
\omega_N=\frac{\pi}{T},
\]
si ha
\[
\frac{\sin(\omega_N T/2)}{\omega_N T/2}
=
\frac{\sin(\pi/2)}{\pi/2}
=
\frac{2}{\pi}.
\]

Quindi il modulo normalizzato alla Nyquist vale
\[
20\log_{10}\!\left(\frac{2}{\pi}\right)\approx -3.92\ \mathrm{dB}.
\]

Negli appunti questa attenuazione viene approssimata qualitativamente come circa \(-3\) dB;
il valore esatto è più vicino a \(-4\) dB.

\section{Fase della tenuta di ordine zero}

La fase di \(F(j\omega)\) si può scrivere come
\[
\angle F(j\omega)
=
-\omega T/2
+
\angle\!\left(\frac{\sin(\omega T/2)}{\omega T/2}\right).
\]

Poiché il termine sinc è reale:
\begin{itemize}
    \item quando
    \[
    \frac{\sin(\omega T/2)}{\omega T/2}>0,
    \]
    il suo argomento è \(0\);
    \item quando
    \[
    \frac{\sin(\omega T/2)}{\omega T/2}<0,
    \]
    il suo argomento è \(\pi\) oppure \(-\pi\), a seconda della convenzione scelta.
\end{itemize}

Per questo la fase complessiva:
\begin{itemize}
    \item decresce linearmente come \(-\omega T/2\);
    \item subisce salti di \(\pi\) in corrispondenza degli zeri del sinc.
\end{itemize}

Se si rappresenta la fase principale, il diagramma ha il classico andamento a rampe
con riavvolgimenti periodici.

\section{Approssimazione a bassa frequenza}

Per frequenze sufficientemente piccole rispetto alla Nyquist, il termine sinc vale circa \(1\):
\[
\frac{\sin(\omega T/2)}{\omega T/2}\approx 1.
\]

Quindi
\[
F(j\omega)\approx T e^{-j\omega T/2}.
\]

Nel dominio di Laplace, questo corrisponde a
\[
F(s)\approx T e^{-sT/2}.
\]

Quindi, a bassa frequenza, la tenuta di ordine zero si comporta come:
\begin{itemize}
    \item un guadagno statico \(T\);
    \item seguito da un ritardo puro pari a \(T/2\).
\end{itemize}

Questo fatto è molto importante perché mostra che il DAC non è neutro:
introduce attenuazione alle alte frequenze e soprattutto una fase aggiuntiva
che può influenzare i margini di stabilità del sistema.

\section{Short Time Fourier Transform}

Nella prima parte della lezione viene anche introdotta la \textbf{Short Time Fourier Transform} (STFT).

L'idea è la seguente:
invece di calcolare una sola trasformata di Fourier su tutto il segnale,
si prende una finestra temporale di lunghezza finita, si calcola la trasformata
su quella finestra, poi si trasla la finestra nel tempo e si ripete il procedimento.

In questo modo si ottiene una rappresentazione bidimensionale:
\begin{itemize}
    \item asse del tempo;
    \item asse della frequenza;
    \item intensità o colore proporzionali al modulo spettrale.
\end{itemize}

\subsection{Interpretazione}

La STFT permette di vedere:
\begin{itemize}
    \item quali frequenze sono presenti;
    \item in quali intervalli temporali compaiono;
    \item come i transitori modificano il contenuto spettrale nel tempo.
\end{itemize}

Per esempio:
\begin{itemize}
    \item un transitorio veloce tende a generare contenuto ad alte frequenze;
    \item un fenomeno più lento si manifesta soprattutto alle basse frequenze.
\end{itemize}

\subsection{Compromesso tempo--frequenza}

La STFT introduce un compromesso fondamentale:
\begin{itemize}
    \item se la finestra è corta, si ha buona risoluzione temporale ma scarsa risoluzione in frequenza;
    \item se la finestra è lunga, si ha buona risoluzione in frequenza ma scarsa risoluzione temporale.
\end{itemize}

Questa è un'altra manifestazione del classico compromesso tempo--frequenza.

\section{Richiamo alle wavelet}

Negli appunti si accenna anche al fatto che esistono trasformate wavelet,
che consentono rappresentazioni tempo-frequenza con risoluzione variabile.

L'idea qualitativa è che le wavelet siano particolarmente adatte
a rappresentare fenomeni localizzati e transitori,
migliorando in molti casi il compromesso rispetto alla STFT classica.

\section{Problema di controllo digitale}

La seconda parte della lezione torna sul problema del controllo digitale.

Lo schema considerato è del tipo:
\[
y^\star \longrightarrow \text{sommatore} \longrightarrow e
\longrightarrow \text{campionatore}
\longrightarrow R(z)
\longrightarrow \text{D/A}
\longrightarrow G(s)
\longrightarrow y.
\]

Il regolatore è quindi implementato in digitale, mentre l'impianto resta continuo.

\subsection{Significato di \(G(s)\)}

Negli appunti si sottolinea che il blocco \(G(s)\) rappresenta in realtà non solo l'impianto,
ma l'intera parte continua dell'anello:
\begin{itemize}
    \item impianto fisico;
    \item attuazione;
    \item amplificazione;
    \item eventualmente anche la dinamica del sensore,
    se la si ingloba nel modello standard a retroazione unitaria.
\end{itemize}

Quindi, in pratica,
\[
G(s)=\text{impianto}+\text{attuazione}+\text{amplificazione}+\text{sensore}.
\]

\section{Tempo di calcolo del regolatore}

Se il regolatore digitale \(R(z)\) è implementato su microprocessore,
esso produce una sequenza discreta \(u(k)\) a partire dai campioni dell'errore.

In linea di principio si potrebbe pensare che, campionato l'ingresso \(e(k)\),
l'uscita \(u(k)\) venga prodotta nello stesso intervallo \([kT,(k+1)T)\).
Tuttavia il tempo di calcolo non è veramente nullo.

Negli appunti si osserva che il ritardo tra campionamento dell'ingresso
e disponibilità dell'uscita:
\begin{itemize}
    \item è inferiore a \(T\);
    \item ma non è noto con precisione;
    \item e soprattutto non è costante.
\end{itemize}

Questo è indesiderabile dal punto di vista del modello.

\subsection{Scelta pratica}

Per evitare questa incertezza, si adotta una convenzione temporale più ordinata:
\begin{itemize}
    \item il regolatore usa i campioni disponibili al passo \(k\);
    \item l'uscita effettiva viene resa coerente con il passo successivo;
    \item in questo modo ingresso e uscita del regolatore sono sempre separati da un intervallo esattamente pari a \(T\).
\end{itemize}

Questa scelta rende il modello più regolare e adatto al progetto.

\section{Due strategie di progetto del controllo digitale}

Negli appunti vengono distinte due possibili strategie.

\subsection{Prima strategia}

La più completa, ma anche più complicata, consiste nel:
\begin{itemize}
    \item discretizzare l'impianto;
    \item lavorare interamente nel dominio discreto;
    \item progettare il regolatore direttamente sul modello discreto.
\end{itemize}

\subsection{Seconda strategia}

Quella su cui il corso si concentra maggiormente è:
\begin{itemize}
    \item progettare il regolatore nel continuo, ottenendo \(R(s)\);
    \item discretizzare il solo regolatore, ad esempio con Tustin;
    \item mantenere l'impianto in forma continua.
\end{itemize}

In questo approccio serve però un buon modello del convertitore D/A,
cioè proprio della tenuta di ordine zero.

\section{Convertitore D/A come sample and hold}

Il convertitore digitale-analogico viene modellato come \textbf{sample and hold},
più precisamente \textbf{zero-order hold}.

In ingresso riceve una sequenza di campioni
\[
u(k),
\]
e in uscita produce un segnale continuo a tratti costante:
\[
u_h(t)=u(k), \qquad kT \le t < (k+1)T.
\]

Quindi il valore del campione viene mantenuto costante per tutta la durata di un periodo \(T\).

\subsection{Interpretazione}

Il segnale in uscita dal D/A non è una successione di impulsi, ma una funzione continua a tratti.
Dal punto di vista fisico, è questo segnale che entra nell'impianto \(G(s)\).

\section{Risposta impulsiva della tenuta di ordine zero}

Per determinare la funzione di trasferimento della tenuta di ordine zero,
si studia la risposta a un impulso discreto in ingresso.

Se in ingresso si applica un solo campione unitario al tempo \(k=0\),
l'uscita analogica della tenuta è un rettangolo:
\[
h(t)=
\begin{cases}
1, & 0\le t < T,\\
0, & t\ge T.
\end{cases}
\]

Equivalentemente:
\[
h(t)=1(t)-1(t-T).
\]

La trasformata di Laplace della risposta impulsiva è allora
\[
F(s)=\mathcal{L}\{h(t)\}
=
\mathcal{L}\{1(t)-1(t-T)\}
=
\frac{1}{s}-\frac{e^{-sT}}{s}
=
\frac{1-e^{-sT}}{s}.
\]

Questa è la funzione di trasferimento della tenuta di ordine zero.

\section{Risposta armonica della tenuta}

Per capire l'effetto del DAC sul comportamento complessivo del sistema,
bisogna studiare la risposta armonica della sua funzione di trasferimento.

Sostituendo \(s=j\omega\), si ottiene
\[
F(j\omega)=\frac{1-e^{-j\omega T}}{j\omega}.
\]

Come già visto, questa si può riscrivere nella forma
\[
F(j\omega)=T e^{-j\omega T/2}\frac{\sin(\omega T/2)}{\omega T/2}.
\]

\subsection{Modulo}

Il modulo vale
\[
|F(j\omega)|=T\left|\frac{\sin(\omega T/2)}{\omega T/2}\right|.
\]

Quindi:
\begin{itemize}
    \item alle basse frequenze il modulo è circa costante e pari a \(T\);
    \item all'aumentare della frequenza compare l'attenuazione di tipo sinc;
    \item il modulo si annulla ai multipli di
    \[
    \omega_c=\frac{2\pi}{T}.
    \]
\end{itemize}

\subsection{Fase}

La fase è data da
\[
\angle F(j\omega)
=
-\omega T/2
+
\angle\!\left(\frac{\sin(\omega T/2)}{\omega T/2}\right).
\]

Nel primo intervallo di frequenze, dove il sinc è positivo,
la fase è approssimativamente
\[
\angle F(j\omega)\approx -\omega T/2.
\]

Quindi, a bassa frequenza, la tenuta si comporta quasi come un ritardo puro di \(T/2\).

\section{Effetto sul diagramma di Bode}

Il DAC con tenuta di ordine zero influenza il diagramma di Bode complessivo dell'impianto:
\begin{itemize}
    \item introduce attenuazione alle alte frequenze;
    \item introduce sfasamento addizionale;
    \item può quindi peggiorare i margini di fase e il comportamento dinamico.
\end{itemize}

Per questo, nel progetto di un regolatore digitale, il blocco di tenuta
non può essere trascurato se si vuole una descrizione accurata del sistema reale.

\section{Sintesi della lezione}

I punti principali della lezione sono:
\begin{enumerate}
    \item la STFT fornisce una rappresentazione tempo--frequenza del segnale;
    \item con la STFT si vede quali frequenze sono presenti e in quali intervalli temporali;
    \item esiste un compromesso inevitabile tra risoluzione temporale e risoluzione in frequenza;
    \item nel problema di controllo digitale il regolatore \(R(z)\) è implementato in forma discreta, mentre il blocco \(G(s)\) resta continuo;
    \item il segnale prodotto dal regolatore digitale deve essere convertito in analogico tramite un convertitore D/A;
    \item il D/A è modellato come tenuta di ordine zero;
    \item la sua funzione di trasferimento è
    \[
    F(s)=\frac{1-e^{-sT}}{s};
    \]
    \item la risposta armonica della tenuta è
    \[
    F(j\omega)=T e^{-j\omega T/2}\frac{\sin(\omega T/2)}{\omega T/2};
    \]
    \item a bassa frequenza il DAC si comporta approssimativamente come un guadagno \(T\) seguito da un ritardo puro \(T/2\);
    \item il blocco di tenuta introduce attenuazione e sfasamento, e quindi modifica il Bode complessivo del sistema.
\end{enumerate}
\chapter{Lezione 11 --- Progetto del controllo digitale via loop shaping continuo, filtro anti-alias, sample-and-hold e discretizzazione del regolatore}

\section{Obiettivo della lezione}

In questa lezione si imposta una procedura pratica di progetto per un sistema di controllo digitale.

L'idea di fondo è la seguente:
\begin{itemize}
    \item l'impianto fisico resta un sistema continuo;
    \item il regolatore viene implementato in digitale;
    \item per il progetto, per\`o, si sfruttano gli strumenti pi\`u ricchi del controllo continuo;
    \item solo alla fine il regolatore continuo viene discretizzato.
\end{itemize}

Questa scelta \`e molto naturale per sistemi SISO tempo-invarianti, perch\'e il progetto nel continuo,
specialmente tramite diagrammi di Bode e loop shaping, \`e pi\`u intuitivo e pi\`u ricco
dal punto di vista degli strumenti disponibili.

\section{Schema del sistema di controllo digitale}

Lo schema fisico considerato \`e il seguente:
\[
y^\star \longrightarrow \text{sommatore} \longrightarrow \text{campionatore} \longrightarrow R(z)
\longrightarrow \text{D/A} \longrightarrow \text{amplificatore + attuatore}
\longrightarrow G(s) \longrightarrow y.
\]

Prima del campionatore bisogna inoltre prevedere un filtro anti-alias analogico.

Quindi, dal punto di vista fisico e implementativo, i blocchi aggiuntivi rispetto al solo impianto sono:
\begin{itemize}
    \item \textbf{\(Q_1(s)\)}: filtro anti-alias;
    \item \textbf{\(Q_2(s)\)}: convertitore D/A, modellato come tenuta di ordine zero;
    \item \textbf{\(Q_3(s)\)}: amplificatore + attuatore.
\end{itemize}

A seconda della convenzione usata, nel modello globale pu\`o essere inclusa anche la dinamica del sensore.

\section{Modello continuo equivalente usato per il progetto}

Per il progetto si inglobano tutte le dinamiche continue nel blocco
\[
\tilde G(s)=Q_1(s)\,Q_2(s)\,Q_3(s)\,G(s).
\]

Si progetta quindi un regolatore continuo \(C(s)\) come se il sistema fosse interamente continuo,
in modo tale che la funzione di trasferimento in ciclo chiuso
\[
H(s)=\frac{C(s)\tilde G(s)}{1+C(s)\tilde G(s)}
\]
soddisfi le specifiche assegnate.

Una volta determinato \(C(s)\), lo si discretizza per ottenere \(R(z)\),
ad esempio con la trasformata di Tustin, con o senza pre-warping.

\section{Perch\'e si progetta \(C(s)\) e non direttamente \(R(z)\)}

Il motivo \`e pratico:
nel continuo si hanno a disposizione strumenti di progetto molto efficaci:
\begin{itemize}
    \item specifiche statiche;
    \item specifiche in transitorio;
    \item margine di fase;
    \item banda passante;
    \item reiezione dei disturbi;
    \item loop shaping sulla funzione d'anello.
\end{itemize}

Per questo si sceglie di:
\begin{enumerate}
    \item progettare \(C(s)\) nel continuo;
    \item poi trasformarlo in \(R(z)\).
\end{enumerate}

\section{Specifiche del problema}

L'impostazione fatta a lezione considera specifiche tipiche del controllo:
\begin{itemize}
    \item prestazioni statiche;
    \item prestazioni di transitorio;
    \item reiezione dei disturbi e del rumore di misura.
\end{itemize}

Nell'esempio svolto, il sistema continuo da controllare \`e
\[
G(s)=\frac{10}{s+3}.
\]

Le specifiche richieste per il sistema in ciclo chiuso sono del tipo:
\begin{itemize}
    \item errore nullo a regime in risposta al gradino;
    \item tempo di assestamento al \(95\%\) inferiore a \(2\) s;
    \item assenza di sovraelongazione nella risposta al gradino;
    \item reiezione del rumore di misura di almeno \(40\) dB per frequenze superiori a \(60\) rad/s.
\end{itemize}

\section{Forma della funzione d'anello}

Nel progetto per loop shaping si lavora sulla funzione di anello aperto
\[
L(s)=C(s)G(s)
\]
oppure, pi\`u correttamente, sul sistema globale continuo
\[
L(s)=C(s)\tilde G(s).
\]

La funzione in ciclo chiuso \`e
\[
H(s)=\frac{L(s)}{1+L(s)}.
\]

\section{Specifica di errore a regime nullo}

Per avere errore nullo al gradino occorre che la funzione d'anello abbia almeno un polo nell'origine.

Quindi il regolatore deve contenere un integratore:
\[
C(s)\sim \frac{1}{s}.
\]

Questa \`e la prima scelta strutturale da fare nel progetto.

\section{Specifica di assenza di sovraelongazione}

Per evitare sovraelongazione, si vuole un sistema in ciclo chiuso con comportamento fortemente smorzato,
idealmente con polo dominante reale oppure con coppia dominante molto smorzata.

Negli appunti viene richiamata la regola empirica gi\`a vista nelle lezioni precedenti:
\[
\zeta \approx \frac{M_\varphi}{100},
\]
con \(M_\varphi\) espresso in gradi.

Per avere una risposta senza overshoot, si richiede quindi un margine di fase elevato, ad esempio
\[
M_\varphi > 75^\circ.
\]

\section{Specifica sul tempo di assestamento}

Se la risposta chiusa pu\`o essere approssimata con quella di un primo ordine
\[
H(s)=\frac{1}{1+s\theta},
\]
allora la risposta al gradino \`e
\[
y(t)=1-e^{-t/\theta}.
\]

Imponendo il tempo di assestamento al \(95\%\),
\[
1-e^{-t_a/\theta}=0.95,
\]
si ricava
\[
e^{-t_a/\theta}=0.05,
\qquad
-\frac{t_a}{\theta}=\ln(0.05),
\qquad
\theta=\frac{t_a}{-\ln(0.05)}.
\]

Con la specifica
\[
t_{a,95\%}<2\ \text{s},
\]
segue
\[
\theta<\frac{2}{-\ln(0.05)} \approx 0.668.
\]

Quindi
\[
\frac{1}{\theta} \gtrsim 1.5\ \text{rad/s}.
\]

Negli appunti, la banda passante desiderata del sistema chiuso viene quindi scelta attorno a
\[
\omega_{B,-3}\approx 2\ \text{rad/s}.
\]

\section{Specifica di reiezione del rumore}

La richiesta di attenuare il rumore di misura di almeno \(40\) dB per
\[
\omega>60\ \text{rad/s}
\]
si traduce in una barriera sul modulo della funzione d'anello:
alle alte frequenze \(L(s)\) deve essere sufficientemente piccolo,
cos\`i che il rumore venga attenuato.

Questo vincolo impone che, oltre una certa pulsazione, il modulo di \(L(s)\)
scenda rapidamente.

\section{Regolatori di tentativo e loop shaping}

Negli appunti il regolatore viene costruito per tentativi, cio\`e plasmando il Bode di \(L(s)\).

\subsection{Primo tentativo}

Una prima scelta naturale \`e
\[
C_1(s)=\frac{k}{s},
\]
in modo da:
\begin{itemize}
    \item introdurre il polo nell'origine;
    \item avere elevato guadagno a bassa frequenza;
    \item iniziare con una pendenza di \(-20\) dB/dec nel diagramma del modulo.
\end{itemize}

\subsection{Correzione con una rete anticipatrice/ritardatrice}

Per soddisfare contemporaneamente:
\begin{itemize}
    \item crossover desiderato;
    \item margine di fase elevato;
    \item attenuazione ad alta frequenza;
\end{itemize}
si modifica il regolatore introducendo uno zero e un polo.

Dalle formule annotate in foto emerge la forma finale del regolatore continuo:
\[
C(s)=\frac{3}{10}\,\frac{1}{s}\,
\frac{1+\dfrac{s}{0.7}}{1+\dfrac{s}{1.5}}.
\]

Questa struttura:
\begin{itemize}
    \item mantiene l'integratore;
    \item aggiunge un aumento di fase nell'intorno della banda d'interesse;
    \item consente di posizionare il taglio a \(0\) dB circa in
    \[
    \omega_c \approx 2\ \text{rad/s};
    \]
    \item lascia poi al polo dell'impianto e agli altri poli il compito di far scendere il modulo alle alte frequenze.
\end{itemize}

\section{Interpretazione del Bode del progetto}

La logica del progetto riportata in lavagna \`e la seguente:
\begin{itemize}
    \item a bassissima frequenza serve l'integratore;
    \item vicino alla zona del crossover si introducono zeri per aumentare la fase;
    \item si lascia poi che i poli del sistema riportino il modulo verso il basso;
    \item ad alta frequenza si aggiungono poli per rispettare la barriera di attenuazione del rumore.
\end{itemize}

Questo \`e il tipico schema di loop shaping:
\[
\text{guadagno alto in bassa frequenza, fase sufficiente al crossover, modulo piccolo in alta frequenza.}
\]

\section{Scelta del periodo di campionamento}

Dopo aver fissato le prestazioni in ciclo chiuso, si passa alla scelta del campionamento.

Negli appunti si usa la regola empirica:
\[
\omega_c^\text{camp} \approx 10\,\omega_{B,-3},
\]
dove \(\omega_c^\text{camp}\) \`e la pulsazione di campionamento.

Poich\'e nel progetto si \`e scelto
\[
\omega_{B,-3}\approx 2\ \text{rad/s},
\]
una prima scelta naturale sarebbe
\[
\omega_c^\text{camp}=20\ \text{rad/s}.
\]

La pulsazione di Nyquist corrispondente sarebbe
\[
\omega_N=\frac{\omega_c^\text{camp}}{2}=10\ \text{rad/s}.
\]

\section{Scelta del filtro anti-alias}

Il filtro anti-alias analogico deve avere una banda passante non superiore alla Nyquist,
per evitare aliasing del rumore e delle componenti veloci.

Quindi, con la prima scelta
\[
\omega_N=10\ \text{rad/s},
\]
bisognerebbe scegliere un anti-alias con frequenza di taglio non superiore a tale valore.

Negli appunti si osserva per\`o che questa scelta pu\`o essere troppo penalizzante,
e si preferisce aumentare la frequenza di campionamento.

Un esempio riportato \`e scegliere
\[
\omega_c^\text{camp}=60\ \text{rad/s},
\qquad
\omega_N=30\ \text{rad/s},
\]
cos\`i da poter usare un filtro anti-alias con pulsazione caratteristica attorno a \(30\) rad/s.

Una possibile scelta di anti-alias richiamata a lezione \`e un passa-basso del terzo ordine del tipo
\[
Q_1(s)=\left(\frac{30}{s+30}\right)^3.
\]

\subsection*{Periodo di campionamento}

Con
\[
\omega_c^\text{camp}=60\ \text{rad/s},
\]
si ottiene
\[
T=\frac{2\pi}{\omega_c^\text{camp}}=\frac{2\pi}{60}\approx 0.1\ \text{s}.
\]

\section{Effetto della tenuta di ordine zero sul margine di fase}

Una volta fissato \(T\), bisogna tenere conto del convertitore D/A con tenuta di ordine zero.

Alle frequenze basse rispetto alla Nyquist, la tenuta si comporta come
\[
F(s)\approx T e^{-sT/2},
\]
quindi introduce un ritardo puro approssimato pari a \(T/2\).

Questo significa che alla pulsazione di crossover \(\omega\) si ha un contributo di fase circa
\[
-\omega \frac{T}{2}.
\]

Nel caso dell'esempio, con
\[
\omega=2\ \text{rad/s},
\qquad
T\approx 0.1\ \text{s},
\]
si ottiene
\[
-\omega\frac{T}{2}\approx -0.1\ \text{rad}\approx -6^\circ.
\]

Negli appunti questa perdita di fase viene ritenuta non trascurabile:
se il margine di fase \`e gi\`a al limite, \(6^\circ\) possono compromettere la specifica.

\section{Campionare pi\`u velocemente}

Per ridurre la perdita di fase dovuta alla tenuta, conviene campionare pi\`u rapidamente,
cio\`e diminuire \(T\).

L'idea annotata a lezione \`e quindi:
\begin{itemize}
    \item mantenere l'anti-alias fissato sulla banda scelta;
    \item aumentare ulteriormente la frequenza di campionamento;
    \item in questo modo il ritardo equivalente \(T/2\) diventa pi\`u piccolo.
\end{itemize}

Cos\`i la perdita di fase della tenuta si riduce e il margine di fase resta compatibile con la specifica.

\section{Da \(C(s)\) a \(R(z)\)}

Una volta definito il regolatore continuo \(C(s)\), si passa a costruire il regolatore digitale \(R(z)\).

La trasformazione scelta a lezione \`e la trasformata di Tustin:
\[
s=\frac{2}{T}\frac{z-1}{z+1}.
\]

\subsection{Pre-warping: serve o no?}

Negli appunti compare la domanda:
\begin{quote}
\emph{Da \(C(s)\) a \(R(z)\) faccio la trasformata di Tustin. Ma faccio o non faccio il pre-warping?}
\end{quote}

La risposta data nel caso in esame \`e:
\begin{itemize}
    \item s\`i, conviene farlo;
    \item perch\'e il progetto del margine di fase \`e gi\`a piuttosto tirato;
    \item senza pre-warping, la deformazione in frequenza della Tustin potrebbe far perdere il rispetto della specifica.
\end{itemize}

Quindi si usa la Tustin con pre-warping sulla frequenza di maggiore interesse,
tipicamente la frequenza di crossover o la frequenza di taglio desiderata.

\section{Riserva di fase nel progetto continuo}

Un messaggio importante della lezione \`e che,
nel progetto continuo del regolatore,
non conviene soddisfare la specifica sul margine di fase in modo troppo poco conservativo.

Infatti, una parte del margine viene poi ``mangiata'' da:
\begin{itemize}
    \item filtro anti-alias;
    \item tenuta di ordine zero;
    \item eventuali altre dinamiche implementative.
\end{itemize}

Per questo \`e buona norma:
\begin{itemize}
    \item o campionare abbastanza velocemente da rendere trascurabili questi effetti;
    \item oppure progettare \(C(s)\) con una riserva di fase aggiuntiva.
\end{itemize}

\section{Procedura complessiva di progetto}

La procedura emersa in questa lezione pu\`o essere riassunta cos\`i:

\begin{enumerate}
    \item Si modellano tutte le dinamiche continue presenti:
    \[
    \tilde G(s)=Q_1(s)Q_2(s)Q_3(s)G(s).
    \]

    \item Si fissano le specifiche:
    \begin{itemize}
        \item errore a regime;
        \item tempo di assestamento;
        \item overshoot;
        \item reiezione del rumore.
    \end{itemize}

    \item Si progetta un regolatore continuo \(C(s)\) mediante loop shaping,
    imponendo che
    \[
    H(s)=\frac{C(s)\tilde G(s)}{1+C(s)\tilde G(s)}
    \]
    soddisfi le specifiche.

    \item Si sceglie la banda passante desiderata del sistema chiuso.

    \item Si sceglie una frequenza di campionamento sufficientemente maggiore della banda chiusa.

    \item Si sceglie un filtro anti-alias compatibile con la Nyquist.

    \item Si verifica la perdita di fase introdotta dalla tenuta di ordine zero.

    \item Si discretizza il regolatore con Tustin, eventualmente con pre-warping.

    \item Si verifica infine il comportamento del sistema discreto tramite simulazione,
    ad esempio con risposta al gradino, rampa e analisi del margine.
\end{enumerate}

\section{Sintesi della lezione}

I punti chiave della lezione sono:
\begin{enumerate}
    \item nel progetto digitale conviene inglobare le dinamiche implementative nel modello continuo
    \[
    \tilde G(s)=Q_1(s)Q_2(s)Q_3(s)G(s);
    \]
    \item il regolatore si progetta inizialmente come \(C(s)\), usando gli strumenti del controllo continuo;
    \item nell'esempio con
    \[
    G(s)=\frac{10}{s+3},
    \]
    si richiedono errore nullo al gradino, transitorio rapido, assenza di overshoot e attenuazione del rumore;
    \item per avere errore nullo al gradino serve un polo nell'origine nel regolatore;
    \item per evitare overshoot si vuole un margine di fase elevato, dell'ordine di \(75^\circ\);
    \item un regolatore continuo compatibile con gli appunti \`e
    \[
    C(s)=\frac{3}{10}\,\frac{1}{s}\,
    \frac{1+\dfrac{s}{0.7}}{1+\dfrac{s}{1.5}};
    \]
    \item la banda passante chiusa viene scelta circa a
    \[
    \omega_{B,-3}\approx 2\ \text{rad/s};
    \]
    \item il periodo di campionamento va scelto a partire da questa banda;
    \item il filtro anti-alias e la tenuta di ordine zero introducono effetti dinamici non trascurabili;
    \item la tenuta introduce una perdita di fase circa pari a
    \[
    -\omega T/2;
    \]
    \item se questa perdita \`e rilevante, bisogna campionare pi\`u velocemente oppure aumentare la riserva di fase del progetto continuo;
    \item il passaggio finale da \(C(s)\) a \(R(z)\) si fa con Tustin, in questo caso preferibilmente con pre-warping.
\end{enumerate}
\chapter{Lezione 12 --- Discretizzazione dell'impianto, Eulero in avanti e all'indietro, modello esatto con tenuta di ordine zero}

\section{Idea della lezione}

Nelle lezioni precedenti abbiamo seguito soprattutto questa filosofia:
\begin{itemize}
    \item progetto un regolatore continuo \(C(s)\);
    \item poi lo discretizzo ottenendo \(R(z)\), ad esempio con Tustin.
\end{itemize}

In questa lezione compare invece un'altra possibilit\`a:
\begin{itemize}
    \item invece di trasformare il controllore discreto in continuo,
    \item discretizzo direttamente l'impianto continuo,
    \item e ottengo un modello esatto o approssimato dell'impianto nel dominio discreto.
\end{itemize}

Questo approccio \`e particolarmente naturale quando l'impianto \`e descritto in \textbf{forma di stato}.

\section{Segnale in uscita dal convertitore D/A}

Il convertitore digitale-analogico con tenuta di ordine zero produce un segnale
\[
u(t)
\]
\textbf{costante a tratti}.

Quindi:
\begin{itemize}
    \item il segnale \`e continuo a tratti ma non liscio;
    \item presenta discontinuit\`a agli istanti di aggiornamento;
    \item dal punto di vista dell'attuatore, non sempre un ingresso cos\`i \`e desiderabile.
\end{itemize}

In pratica, a seconda dell'attuatore e della struttura fisica del sistema,
un segnale a gradini pu\`o essere accettabile oppure no.

\subsection*{Filtro addizionale di smoothing}

Se il segnale costante a tratti non \`e adatto all'attuatore,
si pu\`o inserire dopo il D/A un ulteriore filtro passa-basso analogico,
in modo da ottenere un comando pi\`u regolare.

L'idea \`e:
\begin{itemize}
    \item il D/A genera un segnale costante a tratti;
    \item il filtro lo \emph{addolcisce};
    \item se il filtro ha banda sufficientemente alta, aggiunge poli ad alta frequenza
    ma non modifica in modo sostanziale le prestazioni in ciclo chiuso nella banda di interesse.
\end{itemize}

Quindi, dal punto di vista progettuale, non serve rifare tutto da capo
se questo filtro \`e collocato abbastanza in alto in frequenza.

\section{Derivata e approssimazioni numeriche}

Negli appunti si richiama il fatto che l'operatore \(s\) nel continuo
pu\`o essere visto come operatore di derivazione.

Se ho una funzione continua \(y(t)\), fissato un passo \(h\),
la derivata si pu\`o approssimare in modo differenziale.

\subsection{Eulero in avanti}

L'approssimazione in avanti \`e
\[
\dot y(t)\approx \frac{y(t+h)-y(t)}{h}.
\]

Nel campionamento con periodo \(T\), questo diventa
\[
\dot y(kT)\approx \frac{y((k+1)T)-y(kT)}{T}.
\]

Applicando la trasformata z:
\[
\mathcal{Z}\left\{\dot y(kT)\right\}
\approx
\frac{zY(z)-Y(z)}{T}
=
\frac{z-1}{T}Y(z).
\]

Da qui si ottiene l'approssimazione
\[
s \approx \frac{z-1}{T},
\]
equivalentemente
\[
z=1+Ts.
\]

\subsection{Osservazione}

Questa \textbf{non} \`e la trasformata di Tustin.
\`E una discretizzazione molto pi\`u rozza, detta \textbf{Eulero in avanti}.

\section{Mappatura del piano \(s\) con Eulero in avanti}

Supponiamo di avere
\[
s=\sigma+j\omega.
\]
Con Eulero in avanti, il corrispondente punto nel piano \(z\) \`e
\[
z=1+Ts = 1+T\sigma + jT\omega.
\]

Questa \`e una mappa affine:
\begin{itemize}
    \item le rette verticali del piano \(s\) restano rette verticali nel piano \(z\);
    \item l'asse immaginario \(\sigma=0\) viene mappato nella retta
    \[
    \Re(z)=1;
    \]
    \item il semipiano sinistro non viene mappato interamente dentro il cerchio unitario.
\end{itemize}

Quindi l'Eulero in avanti \textbf{non preserva automaticamente la stabilit\`a}.

\subsection{Esempio}

Se un polo continuo \`e in
\[
s=-10,
\]
allora il polo discreto approssimato diventa
\[
z=1-10T.
\]

Per avere stabilit\`a discreta occorre
\[
|z|<1
\quad\Longrightarrow\quad
|1-10T|<1.
\]

Da cui segue
\[
0<T<\frac{1}{5}.
\]

Quindi un polo stabile nel continuo pu\`o diventare instabile nel discreto
se \(T\) non \`e sufficientemente piccolo.

\section{Eulero all'indietro}

Si considera ora l'approssimazione all'indietro:
\[
\dot y(kT)\approx \frac{y(kT)-y((k-1)T)}{T}.
\]

Applicando la trasformata z:
\[
\mathcal{Z}\left\{\dot y(kT)\right\}
\approx
\frac{Y(z)-z^{-1}Y(z)}{T}
=
\frac{1-z^{-1}}{T}Y(z).
\]

Quindi
\[
s \approx \frac{1-z^{-1}}{T}
=
\frac{z-1}{zT}.
\]

Risolvendo rispetto a \(z\):
\[
z=\frac{1}{1-sT}.
\]

Questa \`e la discretizzazione di \textbf{Eulero a ritroso} o \textbf{backward Euler}.

\section{Mappatura del piano \(s\) con Eulero all'indietro}

Se
\[
s=\sigma+j\omega,
\]
allora
\[
z=\frac{1}{1-sT}
=
\frac{1}{1-(\sigma+j\omega)T}.
\]

Se \(\sigma<0\), cio\`e il polo continuo sta nel semipiano sinistro,
si verifica che
\[
|z|<1.
\]

Quindi, con Eulero a ritroso:
\begin{itemize}
    \item tutto il semipiano sinistro viene mappato all'interno del cerchio unitario;
    \item i poli stabili nel continuo diventano poli stabili nel discreto.
\end{itemize}

Questo rende l'Eulero a ritroso numericamente molto pi\`u robusto di quello in avanti.

\subsection*{Osservazione pratica}

Negli appunti il messaggio \`e chiaro:
\begin{itemize}
    \item Eulero in avanti pu\`o essere usato come approssimazione veloce e rozza;
    \item Eulero a ritroso \`e pi\`u robusto dal punto di vista della stabilit\`a;
    \item ma se si deve \emph{progettare} un controllore, in generale non conviene usare Eulero:
    \`e meglio usare metodi pi\`u accurati, ad esempio Tustin oppure la discretizzazione esatta dell'impianto.
\end{itemize}

\section{Discretizzare direttamente l'impianto}

La seconda parte della lezione introduce il metodo corretto per ottenere
una rappresentazione discreta dell'impianto a partire da un modello continuo in forma di stato.

Supponiamo che il sistema continuo sia descritto da
\[
\begin{cases}
\dot x(t)=Ax(t)+Bu(t),\\[4pt]
y(t)=Cx(t).
\end{cases}
\]

Sappiamo che la soluzione dell'equazione di stato \`e
\[
x(t)=e^{At}x(0)+\int_0^t e^{A(t-\tau)}Bu(\tau)\,d\tau.
\]

Qui
\[
e^{At}
\]
\`e l'esponenziale di matrice, definito tramite serie:
\[
e^{At}=I+At+\frac{A^2t^2}{2!}+\dots+\frac{(At)^k}{k!}+\dots
\]

Nel caso matriciale, questa serie \`e una \emph{definizione}, non solo una conseguenza
della funzione esponenziale scalare.

\section{Ingresso costante a tratti}

Nel nostro schema di controllo digitale, l'ingresso che entra nell'impianto
\`e costante a tratti a causa della tenuta di ordine zero.

Quindi
\[
u(t)=u_k
\qquad
\forall t\in[kT,(k+1)T).
\]

Questa ipotesi \`e fondamentale:
stiamo discretizzando il sistema continuo tenendo conto del fatto che
l'ingresso viene mantenuto costante tra due campioni successivi.

\section{Passaggio da \(t=kT\) a \(t=(k+1)T\)}

Applichiamo la formula della soluzione tra gli istanti \(kT\) e \((k+1)T\):
\[
x((k+1)T)=e^{AT}x(kT)+\int_{kT}^{(k+1)T} e^{A((k+1)T-\tau)}B\,u(\tau)\,d\tau.
\]

Poich\'e nell'intervallo considerato \(u(\tau)=u_k\), si pu\`o portare fuori dall'integrale:
\[
x((k+1)T)=e^{AT}x(kT)+\left(\int_{kT}^{(k+1)T} e^{A((k+1)T-\tau)}\,d\tau\right)B\,u_k.
\]

Facendo il cambio di variabile
\[
\alpha=(k+1)T-\tau,
\]
si ottiene
\[
x((k+1)T)=e^{AT}x(kT)+\left(\int_0^T e^{A\alpha}\,d\alpha\right)B\,u_k.
\]

Definendo
\[
F=e^{AT},
\qquad
G=\left(\int_0^T e^{A\alpha}\,d\alpha\right)B,
\]
si arriva al modello discreto esatto
\[
x(k+1)=Fx(k)+Gu(k).
\]

Questa \`e l'equazione di stato discreta dell'impianto.

\section{Formula alternativa per \(G\)}

Se \(A\) \`e invertibile, allora si pu\`o anche scrivere
\[
G=A^{-1}(e^{AT}-I)B.
\]

Questa formula \`e utile dal punto di vista computazionale,
ma la forma integrale
\[
G=\int_0^T e^{A\alpha}\,d\alpha\,B
\]
resta la pi\`u generale.

\section{Significato del risultato}

Il modello discreto
\[
x(k+1)=Fx(k)+Gu(k)
\]
non \`e un'approssimazione come l'Eulero in avanti o all'indietro:
\`e il \textbf{modello esatto campionato} del sistema continuo
sotto l'ipotesi di ingresso costante a tratti.

Quindi:
\begin{itemize}
    \item \(F=e^{AT}\) descrive l'evoluzione libera del sistema tra due campioni;
    \item \(G\) descrive l'effetto dell'ingresso costante nell'intervallo di campionamento.
\end{itemize}

Questo \`e il modo corretto di ottenere un impianto discreto equivalente al continuo
quando si ha un D/A con tenuta di ordine zero.

\section{Perch\'e questo metodo \`e importante}

Questa discretizzazione esatta \`e molto importante perch\'e permette:
\begin{itemize}
    \item di trattare l'impianto direttamente in tempo discreto;
    \item di evitare le distorsioni introdotte da approssimazioni rozze della derivata;
    \item di progettare controllori discreti direttamente sul modello campionato dell'impianto.
\end{itemize}

Quindi, invece di fare
\[
C(s)\ \longrightarrow\ R(z),
\]
si pu\`o anche fare
\[
G(s)\ \longrightarrow\ \tilde G(z),
\]
e poi progettare direttamente il regolatore discreto su \(\tilde G(z)\).

\section{Sintesi della lezione}

I punti chiave della lezione sono:
\begin{enumerate}
    \item l'uscita del convertitore D/A \`e un segnale costante a tratti;
    \item se necessario, si pu\`o aggiungere un filtro passa-basso analogico per rendere il comando pi\`u regolare;
    \item l'approssimazione della derivata con Eulero in avanti porta a
    \[
    s\approx \frac{z-1}{T},
    \qquad
    z=1+Ts,
    \]
    ma non preserva sempre la stabilit\`a;
    \item l'approssimazione di Eulero a ritroso porta a
    \[
    s\approx \frac{1-z^{-1}}{T}
    =
    \frac{z-1}{zT},
    \qquad
    z=\frac{1}{1-sT},
    \]
    e mappa il semipiano sinistro dentro il cerchio unitario;
    \item per il progetto del controllo, gli schemi di Eulero non sono in generale la scelta migliore;
    \item se l'impianto \`e in forma di stato
    \[
    \dot x=Ax+Bu,\qquad y=Cx,
    \]
    allora, sotto l'ipotesi di ingresso costante a tratti, il modello discreto esatto \`e
    \[
    x(k+1)=Fx(k)+Gu(k),
    \]
    con
    \[
    F=e^{AT},
    \qquad
    G=\left(\int_0^T e^{A\alpha}\,d\alpha\right)B;
    \]
    \item se \(A\) \`e invertibile, si pu\`o scrivere anche
    \[
    G=A^{-1}(e^{AT}-I)B;
    \]
    \item questa discretizzazione tiene conto in modo esatto della tenuta di ordine zero.
\end{enumerate}
\chapter{Lezione 13 --- Assegnamento degli autovalori mediante retroazione completa dello stato}

\section{Impostazione del problema}

Consideriamo un sistema lineare tempo-invariante in forma di stato
\[
\begin{cases}
\dot x(t)=Ax(t)+Bu(t),\\[4pt]
y(t)=Cx(t),
\end{cases}
\]
dove, per evitare ambiguità di notazione,
\[
x(t)\in\mathbb{R}^n,\qquad
u(t)\in\mathbb{R}^r,\qquad
y(t)\in\mathbb{R}^\ell,
\]
e
\[
A\in\mathbb{R}^{n\times n},\qquad
B\in\mathbb{R}^{n\times r},\qquad
C\in\mathbb{R}^{\ell\times n}.
\]

Vogliamo controllare il sistema attraverso una \textbf{retroazione dello stato}.

\subsection{Problema di regolazione e problema di tracking}

Negli appunti si distinguono due casi:
\begin{itemize}
    \item se il riferimento \(y^\star\) è costante, si ha un \textbf{problema di regolazione};
    \item se il riferimento è tempo-variabile,
    \[
    y^\star=y^\star(t),
    \]
    si ha un \textbf{problema di tracking}.
\end{itemize}

Nel seguito ci concentriamo soprattutto sulla parte di stabilizzazione e assegnamento degli autovalori,
che riguarda la dinamica interna del sistema retroazionato.

\section{Retroazione completa dello stato}

Supponiamo che tutto lo stato sia misurabile e disponibile.
In tal caso possiamo usare una legge di controllo del tipo
\[
u(t)=-Kx(t)+n(t),
\]
dove:
\begin{itemize}
    \item \(K\) è la matrice di guadagno di retroazione;
    \item \(n(t)\) è un ingresso esterno, usato ad esempio per il riferimento.
\end{itemize}

Nel caso SISO, \(K\) è un vettore riga:
\[
K\in\mathbb{R}^{1\times n}.
\]

\subsection*{Osservazione}

L'ipotesi cruciale è che lo stato sia \emph{tutto misurabile}.
Qui non stiamo ancora affrontando il problema dell'osservatore:
si assume direttamente di conoscere \(x(t)\).

\section{Dinamica del sistema retroazionato}

Sostituendo la legge di controllo
\[
u(t)=-Kx(t)+n(t)
\]
nell'equazione di stato,
si ottiene
\[
\dot x(t)=Ax(t)+B\bigl(-Kx(t)+n(t)\bigr).
\]

Quindi
\[
\dot x(t)=Ax(t)-BKx(t)+Bn(t),
\]
ossia
\[
\dot x(t)=(A-BK)x(t)+Bn(t).
\]

La matrice di stato del sistema retroazionato è dunque
\[
A_{\text{cl}}=A-BK.
\]

\section{Autovalori del sistema retroazionato}

Gli autovalori della matrice \(A\) sono le radici del polinomio caratteristico
\[
p_A(\lambda)=\det(\lambda I-A).
\]

Dopo la retroazione di stato, gli autovalori del sistema diventano quelli della matrice
\[
A-BK,
\]
cioè le radici di
\[
\det(\lambda I-A+BK)=0.
\]

La domanda fondamentale è quindi:

\begin{quote}
\emph{Posso progettare \(K\) in modo da assegnare arbitrariamente gli autovalori di \(A-BK\)?}
\end{quote}

\section{Teorema di assegnamento degli autovalori}

\subsection*{Caso continuo}

Negli appunti il teorema è formulato nel caso tempo continuo.

\begin{theorem}
Se la coppia \((A,B)\) è completamente controllabile, allora esiste una matrice di retroazione \(K\)
tale da assegnare arbitrariamente gli \(n\) autovalori della matrice
\[
A-BK.
\]
\end{theorem}

Nel caso continuo, controllabilità e raggiungibilità sono equivalenti,
quindi l'ipotesi può essere formulata in entrambi i modi.

\subsection*{Osservazione}

Il termine di riferimento \(n(t)\) non entra nel teorema:
esso non modifica la matrice di stato del sistema retroazionato.
Per lo studio dell'assegnamento degli autovalori si può quindi porre
\[
n(t)=0.
\]

\section{Caso SISO e forma canonica di controllo}

La dimostrazione negli appunti viene svolta nel caso SISO, cioè per
\[
\dot x=Ax+bu,
\qquad b\in\mathbb{R}^n.
\]

Se la coppia \((A,b)\) è completamente controllabile,
allora esiste una trasformazione di coordinate invertibile
\[
z=P^{-1}x
\]
che porta il sistema in \textbf{forma canonica di controllo}:
\[
\dot z=A_{cc}z+b_{cc}u.
\]

La matrice \(A_{cc}\) ha la forma
\[
A_{cc}=
\begin{bmatrix}
0 & 1 & 0 & \cdots & 0\\
0 & 0 & 1 & \cdots & 0\\
\vdots & \vdots & \vdots & \ddots & \vdots\\
0 & 0 & 0 & \cdots & 1\\
-a_0 & -a_1 & -a_2 & \cdots & -a_{n-1}
\end{bmatrix},
\]
mentre
\[
b_{cc}=
\begin{bmatrix}
0\\
0\\
\vdots\\
0\\
1
\end{bmatrix}.
\]

I coefficienti \(a_0,\dots,a_{n-1}\) sono quelli del polinomio caratteristico di \(A\):
\[
\det(\lambda I-A)
=
\lambda^n+a_{n-1}\lambda^{n-1}+\cdots+a_1\lambda+a_0.
\]

\subsection*{Osservazione sulla trasformazione di similitudine}

Poiché
\[
A_{cc}=P^{-1}AP,
\]
la trasformazione è una trasformazione di similitudine, quindi conserva gli autovalori.

In particolare,
\[
\det(\lambda I-A)=\det(\lambda I-A_{cc}).
\]

\section{Retroazione in forma canonica}

Nel nuovo sistema di coordinate si applica la retroazione
\[
u=-K_{cc}z,
\qquad
K_{cc}=
\begin{bmatrix}
k_1 & k_2 & \cdots & k_n
\end{bmatrix}.
\]

La dinamica chiusa diventa
\[
\dot z=(A_{cc}-b_{cc}K_{cc})z.
\]

Poiché
\[
b_{cc}K_{cc}
=
\begin{bmatrix}
0\\
0\\
\vdots\\
0\\
1
\end{bmatrix}
\begin{bmatrix}
k_1 & k_2 & \cdots & k_n
\end{bmatrix}
=
\begin{bmatrix}
0 & 0 & \cdots & 0\\
0 & 0 & \cdots & 0\\
\vdots & \vdots & \ddots & \vdots\\
k_1 & k_2 & \cdots & k_n
\end{bmatrix},
\]
si ottiene
\[
A_{cc}-b_{cc}K_{cc}
=
\begin{bmatrix}
0 & 1 & 0 & \cdots & 0\\
0 & 0 & 1 & \cdots & 0\\
\vdots & \vdots & \vdots & \ddots & \vdots\\
0 & 0 & 0 & \cdots & 1\\
-(a_0+k_1) & -(a_1+k_2) & \cdots & -(a_{n-2}+k_{n-1}) & -(a_{n-1}+k_n)
\end{bmatrix}.
\]

Questa è ancora una matrice in forma canonica di controllo.

\section{Polinomio caratteristico del sistema retroazionato}

Il polinomio caratteristico della matrice retroazionata è
\[
\det(\lambda I-A_{cc}+b_{cc}K_{cc})
=
\lambda^n
+(a_{n-1}+k_n)\lambda^{n-1}
+\cdots
+(a_1+k_2)\lambda
+(a_0+k_1).
\]

Supponiamo di voler assegnare gli autovalori desiderati
\[
\lambda_1,\lambda_2,\dots,\lambda_n.
\]

Allora imponiamo come polinomio desiderato
\[
(\lambda-\lambda_1)(\lambda-\lambda_2)\cdots(\lambda-\lambda_n)
=
\lambda^n+\alpha_{n-1}\lambda^{n-1}+\cdots+\alpha_1\lambda+\alpha_0.
\]

Uguagliando i coefficienti dei due polinomi, si ricava:
\[
\alpha_{n-1}=a_{n-1}+k_n,
\]
\[
\alpha_{n-2}=a_{n-2}+k_{n-1},
\]
\[
\vdots
\]
\[
\alpha_0=a_0+k_1.
\]

Quindi i guadagni si scelgono come
\[
k_1=\alpha_0-a_0,
\qquad
k_2=\alpha_1-a_1,
\qquad
\dots,
\qquad
k_n=\alpha_{n-1}-a_{n-1}.
\]

Questo mostra che, in forma canonica di controllo, i guadagni di retroazione
permettono di assegnare arbitrariamente il polinomio caratteristico.

\section{Ritorno alle coordinate originali}

Poiché
\[
z=P^{-1}x
\qquad\Longrightarrow\qquad
x=Pz,
\]
si ha
\[
u=-Kx=-KPz.
\]

Confrontando con
\[
u=-K_{cc}z,
\]
segue
\[
K_{cc}=KP.
\]

Dunque
\[
K=K_{cc}P^{-1}.
\]

Questa formula permette di tornare dal sistema in forma canonica
al sistema nelle coordinate originarie.

\section{Sufficienza dell'ipotesi di controllabilità}

La costruzione appena fatta dimostra la \textbf{sufficienza} dell'ipotesi:
se \((A,b)\) è completamente controllabile,
allora è possibile assegnare arbitrariamente tutti gli autovalori del sistema retroazionato.

\section{Necessità dell'ipotesi di controllabilità}

Negli appunti viene anche mostrato perché l'ipotesi di controllabilità è necessaria.

Supponiamo che la coppia \((A,b)\) non sia completamente controllabile
e che il sottospazio di controllabilità abbia dimensione
\[
p<n.
\]

Allora esiste un cambiamento di coordinate
\[
z=Q^{-1}x
\]
che separa la parte controllabile e quella non controllabile:
\[
z=
\begin{bmatrix}
z_c\\
z_{nc}
\end{bmatrix},
\qquad
z_c\in\mathbb{R}^p,
\qquad
z_{nc}\in\mathbb{R}^{n-p}.
\]

Nel nuovo sistema di coordinate la dinamica assume la forma
\[
\dot z
=
\begin{bmatrix}
\tilde A_c & \tilde A_{12}\\
0 & \tilde A_{nc}
\end{bmatrix}
z
+
\begin{bmatrix}
\tilde b_c\\
0
\end{bmatrix}u.
\]

Scrivendo anche la retroazione in tali coordinate,
\[
\tilde K=KQ=
\begin{bmatrix}
\tilde K_c & \tilde K_{nc}
\end{bmatrix},
\]
la matrice chiusa diventa
\[
\tilde A-\tilde b\tilde K
=
\begin{bmatrix}
\tilde A_c-\tilde b_c\tilde K_c & \tilde A_{12}-\tilde b_c\tilde K_{nc}\\
0 & \tilde A_{nc}
\end{bmatrix}.
\]

Essendo una matrice triangolare a blocchi, il polinomio caratteristico fattorizza:
\[
\det(\lambda I-\tilde A+\tilde b\tilde K)
=
\det(\lambda I_p-\tilde A_c+\tilde b_c\tilde K_c)\,
\det(\lambda I_{n-p}-\tilde A_{nc}).
\]

Quindi:
\begin{itemize}
    \item i primi \(p\) autovalori, associati alla parte controllabile, possono essere spostati;
    \item i restanti \(n-p\) autovalori, associati alla parte non controllabile, \textbf{non vengono modificati} dalla retroazione di stato.
\end{itemize}

\subsection*{Conclusione}

Se la coppia non è completamente controllabile, non posso assegnare arbitrariamente tutti gli autovalori:
posso al più assegnarne \(p\), dove \(p\) è la dimensione del sottospazio controllabile.

Questo dimostra la necessità dell'ipotesi del teorema.

\section{Interpretazione fisica}

Gli autovalori della matrice \(A\) sono determinati dalla fisica del sistema.
La matrice \(B\) descrive come l'ingresso agisce sul sistema.

La retroazione di stato modifica la dinamica attraverso il termine
\[
-BKx,
\]
quindi cambia la matrice di stato da \(A\) a
\[
A-BK.
\]

Ma questa modifica può influenzare soltanto i modi dinamici che sono effettivamente raggiungibili
dall'ingresso.

Per questo:
\begin{itemize}
    \item i modi controllabili possono essere spostati;
    \item i modi non controllabili restano dov'erano.
\end{itemize}

\section{Forma finale del problema di controllo}

Il problema di assegnamento degli autovalori si può quindi formulare così:

dato il sistema
\[
\begin{cases}
\dot x=Ax+Bu,\\
y=Cx,
\end{cases}
\]
si vuole trovare una legge di controllo
\[
u(t)=-Kx(t)+n(t)
\]
tale che la matrice chiusa
\[
A-BK
\]
abbia autovalori desiderati.

Se \((A,B)\) è completamente controllabile, questo è sempre possibile.

\section{Sintesi della lezione}

I punti chiave della lezione sono:
\begin{enumerate}
    \item si considera una retroazione completa dello stato:
    \[
    u(t)=-Kx(t)+n(t);
    \]
    \item la matrice di stato del sistema retroazionato è
    \[
    A-BK;
    \]
    \item gli autovalori del sistema chiuso sono le radici di
    \[
    \det(\lambda I-A+BK)=0;
    \]
    \item se la coppia \((A,B)\) è completamente controllabile, allora è possibile assegnare arbitrariamente tutti gli autovalori del sistema chiuso;
    \item nel caso SISO, la dimostrazione si fa portando il sistema in forma canonica di controllo;
    \item in forma canonica, i coefficienti del polinomio caratteristico dipendono linearmente dai guadagni \(k_i\);
    \item uguagliando i coefficienti del polinomio desiderato si ricavano direttamente i guadagni;
    \item tornando alle coordinate originali si ottiene il guadagno \(K\) del sistema reale;
    \item se il sistema non è completamente controllabile, soltanto gli autovalori della parte controllabile possono essere assegnati;
    \item gli autovalori non controllabili non vengono modificati dalla retroazione di stato.
\end{enumerate}
\chapter{Lezione 14 --- Formula di Ackermann e assegnamento dell'autostruttura}

\section{Impostazione del problema}

Consideriamo un sistema lineare tempo-invariante in forma di stato
\[
\begin{cases}
\dot x(t)=Ax(t)+Bu(t),\\[4pt]
y(t)=Cx(t),
\end{cases}
\]
con
\[
x(t)\in\mathbb{R}^n,\qquad
u(t)\in\mathbb{R}^m,\qquad
y(t)\in\mathbb{R}^\ell,
\]
e
\[
A\in\mathbb{R}^{n\times n},\qquad
B\in\mathbb{R}^{n\times m},\qquad
C\in\mathbb{R}^{\ell\times n}.
\]

Vogliamo assegnare la dinamica del sistema in ciclo chiuso mediante una
\textbf{retroazione completa dello stato}
\[
u(t)=-Kx(t),
\]
dove
\[
K\in\mathbb{R}^{m\times n}.
\]

Sostituendo nella dinamica di stato si ottiene
\[
\dot x(t)=Ax(t)+B(-Kx(t))=(A-BK)x(t).
\]

Quindi la matrice di stato del sistema retroazionato è
\[
A_{\text{cl}}=A-BK.
\]

Il problema consiste nello scegliere \(K\) in modo che \(A-BK\) abbia gli autovalori,
e possibilmente anche gli autovettori, desiderati.

\section{Richiamo: assegnamento degli autovalori}

Dalla lezione precedente sappiamo che:
\begin{itemize}
    \item nel caso SISO, se \((A,b)\) è completamente controllabile,
    si possono assegnare arbitrariamente tutti gli autovalori della matrice \(A-bk^T\);
    \item nel caso MIMO, se \((A,B)\) è completamente controllabile,
    si possono assegnare arbitrariamente gli autovalori della matrice \(A-BK\).
\end{itemize}

Nel caso MIMO, tuttavia, la soluzione non è unica, perché \(K\) ha più gradi di libertà.
Questo apre la possibilità di assegnare non solo gli autovalori, ma anche, entro certi limiti,
gli autovettori associati: cioè l'\textbf{autostruttura} del sistema in ciclo chiuso.

\section{Caso SISO: formula di Ackermann}

Consideriamo dapprima il caso SISO:
\[
\dot x=Ax+bu,\qquad y=c^T x,
\]
con
\[
b\in\mathbb{R}^n,\qquad c\in\mathbb{R}^n,
\]
e legge di retroazione
\[
u=-k^T x,
\qquad
k\in\mathbb{R}^n.
\]

La matrice chiusa è
\[
A-bk^T.
\]

\subsection{Polinomio caratteristico desiderato}

Supponiamo di voler assegnare gli autovalori desiderati
\[
\lambda_1,\lambda_2,\dots,\lambda_n.
\]

Il polinomio caratteristico desiderato è allora
\[
p_{\text{cl}}(\lambda)=
(\lambda-\lambda_1)(\lambda-\lambda_2)\cdots(\lambda-\lambda_n).
\]

Scrivendolo in forma espansa:
\[
p_{\text{cl}}(\lambda)=
\lambda^n+\alpha_{n-1}\lambda^{n-1}+\cdots+\alpha_1\lambda+\alpha_0.
\]

\subsection{Matrice di controllabilità}

Nel caso SISO la matrice di controllabilità è
\[
\mathcal{R}=
\begin{bmatrix}
b & Ab & A^2b & \cdots & A^{n-1}b
\end{bmatrix}.
\]

Se la coppia \((A,b)\) è completamente controllabile, allora \(\mathcal{R}\) è quadrata e invertibile:
\[
\mathcal{R}\in\mathbb{R}^{n\times n},
\qquad
\rank(\mathcal{R})=n.
\]

\subsection{Formula di Ackermann}

La \textbf{formula di Ackermann} permette di scrivere direttamente il guadagno di retroazione
senza passare esplicitamente per la forma canonica di controllo:
\[
k^T=
\begin{bmatrix}
0 & \cdots & 0 & 1
\end{bmatrix}
\mathcal{R}^{-1}\,p_{\text{cl}}(A),
\]
dove
\[
p_{\text{cl}}(A)=A^n+\alpha_{n-1}A^{n-1}+\cdots+\alpha_1A+\alpha_0 I.
\]

Questa formula è molto utile nel caso SISO, perché fornisce immediatamente il guadagno \(k\)
che assegna gli autovalori scelti.

\subsection*{Osservazione}

La formula di Ackermann è concettualmente equivalente al passaggio in forma canonica di controllo,
ma è più rapida dal punto di vista del calcolo simbolico o numerico.

\section{Caso MIMO: più gradi di libertà}

Passiamo ora al caso MIMO:
\[
\dot x=Ax+Bu,
\qquad
u=-Kx,
\qquad
K\in\mathbb{R}^{m\times n}.
\]

Anche qui la matrice chiusa è
\[
A-BK.
\]

Se la coppia \((A,B)\) è completamente controllabile,
allora è possibile assegnare arbitrariamente gli autovalori del sistema in ciclo chiuso.

Tuttavia, a differenza del caso SISO, qui:
\begin{itemize}
    \item la matrice \(K\) ha \(mn\) coefficienti;
    \item il numero di vincoli imposti dal solo polinomio caratteristico è \(n\).
\end{itemize}

Quindi, in generale, rimangono dei gradi di libertà.
Questi possono essere sfruttati per cercare di assegnare anche gli autovettori,
cioè l'autostruttura del sistema chiuso.

\section{Assegnamento dell'autostruttura}

Supponiamo di voler assegnare a \(A-BK\) gli autovalori
\[
\lambda_1,\lambda_2,\dots,\lambda_n
\]
e, per quanto possibile, anche gli autovettori
\[
v_1,v_2,\dots,v_n.
\]

Se \(\lambda_i\) è autovalore di \(A-BK\) e \(v_i\) è il corrispondente autovettore,
allora deve valere
\[
(A-BK)v_i=\lambda_i v_i.
\]

Portando tutto a sinistra:
\[
(\lambda_i I-A+B K)v_i=0.
\]

Scrivendo
\[
q_i:=Kv_i\in\mathbb{R}^m,
\]
la relazione precedente diventa
\[
(\lambda_i I-A)v_i+Bq_i=0.
\]

Questa equazione si può riscrivere in forma matriciale come
\[
\begin{bmatrix}
\lambda_i I-A & B
\end{bmatrix}
\begin{bmatrix}
v_i\\
q_i
\end{bmatrix}
=0.
\]

Definiamo quindi
\[
T_i:=
\begin{bmatrix}
\lambda_i I-A & B
\end{bmatrix}
\in\mathbb{R}^{n\times (n+m)}.
\]

Dobbiamo allora cercare vettori
\[
p_i=
\begin{bmatrix}
v_i\\
q_i
\end{bmatrix}
\in \ker(T_i).
\]

\section{Dimensione del nucleo di \(T_i\)}

Poiché \((A,B)\) è completamente controllabile, per il criterio di Hautus si ha
\[
\rank\!\begin{bmatrix}
\lambda I-A & B
\end{bmatrix}=n
\qquad
\forall \lambda\in\mathbb{C}.
\]

Quindi, per ogni \(\lambda_i\),
\[
\rank(T_i)=n.
\]

Essendo \(T_i\in\mathbb{R}^{n\times (n+m)}\), segue
\[
\dim\ker(T_i)=(n+m)-n=m.
\]

Dunque, per ciascun autovalore scelto \(\lambda_i\),
esistono \(m\) vettori linearmente indipendenti nel nucleo di \(T_i\).

\section{Costruzione dei vettori \(p_i\)}

Sia
\[
\{m_{i1},m_{i2},\dots,m_{im}\}
\]
una base di \(\ker(T_i)\).

Allora ogni vettore del nucleo si può scrivere come combinazione lineare
\[
p_i=\sum_{j=1}^{m}\alpha_{ij}\,m_{ij}.
\]

Poiché vogliamo che le prime \(n\) componenti di \(p_i\) coincidano con l'autovettore desiderato \(v_i\),
scriviamo
\[
p_i=
\begin{bmatrix}
v_i\\
q_i
\end{bmatrix}.
\]

La scelta dei coefficienti \(\alpha_{ij}\) viene fatta proprio in modo che
la parte alta del vettore \(p_i\) sia il vettore \(v_i\) desiderato.
La parte bassa fornisce automaticamente il vettore
\[
q_i=Kv_i.
\]

Ripetendo questa procedura per tutti gli autovalori \(\lambda_i\),
si ottengono i vettori
\[
p_1,p_2,\dots,p_n.
\]

\section{Costruzione di \(K\)}

A questo punto costruiamo le matrici
\[
V=
\begin{bmatrix}
v_1 & v_2 & \cdots & v_n
\end{bmatrix},
\qquad
Q=
\begin{bmatrix}
q_1 & q_2 & \cdots & q_n
\end{bmatrix}.
\]

Poiché per ciascun \(i\) vale
\[
q_i=Kv_i,
\]
mettendo insieme tutte le colonne si ottiene
\[
KV=Q.
\]

Se \(V\) è invertibile, allora
\[
K=QV^{-1}.
\]

Questa è la formula finale che permette di ricavare la matrice di guadagno
a partire dagli autovettori desiderati e dai corrispondenti vettori \(q_i\).

\section{Algoritmo operativo}

Riassumendo, l'algoritmo per l'assegnamento dell'autostruttura nel caso MIMO è:

\begin{enumerate}
    \item fissare gli autovalori desiderati
    \[
    \lambda_1,\lambda_2,\dots,\lambda_n;
    \]

    \item per ogni \(i=1,\dots,n\), costruire la matrice
    \[
    T_i=
    \begin{bmatrix}
    \lambda_i I-A & B
    \end{bmatrix};
    \]

    \item calcolare una base del nucleo \(\ker(T_i)\);

    \item scegliere una combinazione lineare dei vettori della base in modo che
    la parte alta del vettore risultante sia l'autovettore desiderato \(v_i\);

    \item leggere la parte bassa come \(q_i\), cioè
    \[
    p_i=
    \begin{bmatrix}
    v_i\\
    q_i
    \end{bmatrix};
    \]

    \item costruire le matrici
    \[
    V=
    \begin{bmatrix}
    v_1 & \cdots & v_n
    \end{bmatrix},
    \qquad
    Q=
    \begin{bmatrix}
    q_1 & \cdots & q_n
    \end{bmatrix};
    \]

    \item se \(V\) è invertibile, porre
    \[
    K=QV^{-1}.
    \]
\end{enumerate}

\section{Esempio}

Consideriamo il sistema
\[
\dot x=
\begin{bmatrix}
-1 & 1\\
0 & -1
\end{bmatrix}x
+
\begin{bmatrix}
1 & 0\\
0 & 1
\end{bmatrix}u.
\]

Qui:
\[
n=2,\qquad m=2.
\]

Vogliamo assegnare in ciclo chiuso gli autovalori
\[
\lambda_1=-5,\qquad \lambda_2=-10.
\]

\subsection{Primo autovalore}

Per \(\lambda_1=-5\):
\[
T_1=
\begin{bmatrix}
-4 & -1 & 1 & 0\\
0 & -4 & 0 & 1
\end{bmatrix}.
\]

Una base del nucleo è, ad esempio,
\[
m_{11}=
\begin{bmatrix}
1\\0\\4\\0
\end{bmatrix},
\qquad
m_{12}=
\begin{bmatrix}
0\\1\\1\\4
\end{bmatrix}.
\]

Scegliamo come autovettore desiderato
\[
v_1=
\begin{bmatrix}
1\\0
\end{bmatrix}.
\]

Allora basta prendere
\[
p_1=m_{11}=
\begin{bmatrix}
1\\0\\4\\0
\end{bmatrix},
\]
da cui
\[
q_1=
\begin{bmatrix}
4\\0
\end{bmatrix}.
\]

\subsection{Secondo autovalore}

Per \(\lambda_2=-10\):
\[
T_2=
\begin{bmatrix}
-9 & -1 & 1 & 0\\
0 & -9 & 0 & 1
\end{bmatrix}.
\]

Una base del nucleo è
\[
m_{21}=
\begin{bmatrix}
1\\0\\9\\0
\end{bmatrix},
\qquad
m_{22}=
\begin{bmatrix}
0\\1\\1\\9
\end{bmatrix}.
\]

Scegliamo come autovettore desiderato
\[
v_2=
\begin{bmatrix}
0\\1
\end{bmatrix}.
\]

Allora possiamo prendere
\[
p_2=m_{22}=
\begin{bmatrix}
0\\1\\1\\9
\end{bmatrix},
\]
da cui
\[
q_2=
\begin{bmatrix}
1\\9
\end{bmatrix}.
\]

\subsection{Costruzione finale di \(K\)}

Si ottiene quindi
\[
V=
\begin{bmatrix}
1 & 0\\
0 & 1
\end{bmatrix},
\qquad
Q=
\begin{bmatrix}
4 & 1\\
0 & 9
\end{bmatrix}.
\]

Poiché \(V=I\), segue immediatamente
\[
K=QV^{-1}=Q=
\begin{bmatrix}
4 & 1\\
0 & 9
\end{bmatrix}.
\]

Verifica:
\[
A-BK=
\begin{bmatrix}
-1 & 1\\
0 & -1
\end{bmatrix}
-
\begin{bmatrix}
4 & 1\\
0 & 9
\end{bmatrix}
=
\begin{bmatrix}
-5 & 0\\
0 & -10
\end{bmatrix}.
\]

Quindi gli autovalori assegnati sono effettivamente
\[
-5,\qquad -10,
\]
e in questo caso si è assegnata completamente anche l'autostruttura.

\section{Osservazioni numeriche e robustezza}

Nell'esempio precedente l'inversione di \(V\) è banale.
In generale, però, la costruzione
\[
K=QV^{-1}
\]
può essere numericamente delicata.

Se \(V\) è invertibile ma mal condizionata, cioè con determinante molto piccolo o colonne quasi dipendenti,
l'inversione numerica può introdurre errori rilevanti.
In tal caso il guadagno calcolato potrebbe non collocare gli autovalori esattamente dove previsto.

Per migliorare la robustezza numerica conviene, quando possibile, scegliere autovettori
il più possibile ortogonali o comunque ben condizionati.

\section{Quando si può assegnare tutta l'autostruttura?}

Nel caso dell'esempio avevamo
\[
n=m.
\]

Quindi il numero di ingressi indipendenti era uguale al numero di stati,
e questo ha permesso di assegnare completamente l'autostruttura:
non solo gli autovalori, ma anche tutti gli autovettori scelti.

In generale:
\begin{itemize}
    \item gli autovalori si possono assegnare se \((A,B)\) è completamente controllabile;
    \item l'assegnamento completo degli autovettori dipende dal numero di ingressi e dalla struttura del sistema;
    \item in MIMO ci sono più gradi di libertà, ma non sempre è possibile assegnare arbitrariamente \emph{ogni} autovettore.
\end{itemize}

\section{Caso tempo discreto}

Negli appunti si osserva che nei sistemi a tempo discreto tutto funziona in modo del tutto analogo.

Consideriamo
\[
x(k+1)=Ax(k)+Bu(k),
\]
con retroazione di stato
\[
u(k)=-Kx(k).
\]

Il sistema chiuso diventa
\[
x(k+1)=(A-BK)x(k).
\]

La teoria, le formule e la procedura sono le stesse del caso continuo,
con l'unica differenza concettuale che nel tempo discreto l'ipotesi corretta è che
la coppia \((A,B)\) sia \textbf{completamente raggiungibile}.

Nel caso tempo discreto:
\begin{itemize}
    \item si assegnano autovalori desiderati nel piano \(z\);
    \item per la stabilità, tali autovalori devono stare dentro il cerchio unitario;
    \item la formula di Ackermann e l'algoritmo MIMO si applicano nello stesso modo,
    sostituendo al polinomio in \(\lambda\) il polinomio desiderato in \(z\).
\end{itemize}

\section{Sintesi della lezione}

I punti principali della lezione sono:
\begin{enumerate}
    \item nel caso SISO, il guadagno di retroazione può essere calcolato direttamente con la formula di Ackermann:
    \[
    k^T=
    \begin{bmatrix}
    0 & \cdots & 0 & 1
    \end{bmatrix}
    \mathcal{R}^{-1}p_{\text{cl}}(A);
    \]

    \item nel caso MIMO, la soluzione non è unica perché \(K\) ha più gradi di libertà;

    \item questi gradi di libertà possono essere sfruttati per assegnare non solo gli autovalori,
    ma anche gli autovettori, cioè l'autostruttura;

    \item per ogni autovalore desiderato \(\lambda_i\), si costruisce
    \[
    T_i=
    \begin{bmatrix}
    \lambda_i I-A & B
    \end{bmatrix};
    \]

    \item il vettore
    \[
    p_i=
    \begin{bmatrix}
    v_i\\
    q_i
    \end{bmatrix}
    \]
    deve appartenere al nucleo di \(T_i\);

    \item ripetendo la costruzione per tutti gli autovalori si ottengono
    \[
    KV=Q,
    \qquad
    K=QV^{-1};
    \]

    \item in presenza di matrici mal condizionate, l'inversione di \(V\) può creare problemi numerici;

    \item nel caso tempo discreto la teoria è la stessa del continuo, purché la coppia \((A,B)\) sia completamente raggiungibile.
\end{enumerate}
\chapter{Lezione 15 --- Dead-beat control, sistemi stabilizzabili, stato aumentato e introduzione all'osservatore}

\section{Dead-beat controller}

La prima osservazione della lezione riguarda un sistema a tempo discreto
\[
x(k+1)=Ax(k)+Bu(k),
\]
controllato con una retroazione completa dello stato
\[
u(k)=-Kx(k).
\]

La dinamica in ciclo chiuso è quindi
\[
x(k+1)=(A-BK)x(k).
\]

Nel tempo discreto è possibile scegliere gli autovalori della matrice \(A-BK\) \emph{tutti in zero}.
Questo porta al cosiddetto \textbf{dead-beat controller}.

\subsection{Idea fondamentale}

Se riesco ad assegnare
\[
\lambda_1=\lambda_2=\cdots=\lambda_n=0,
\]
allora il polinomio caratteristico del sistema chiuso è
\[
p_{cl}(z)=z^n.
\]

In questo caso la matrice \(A-BK\) è nilpotente e, per un sistema di ordine \(n\),
lo stato va esattamente a zero in un numero finito di passi (al più \(n\), nel caso genericamente controllabile).

Questa è una differenza importante rispetto al tempo continuo:
nel continuo ``mandare un autovalore a \(-\infty\)'' non ha significato fisico,
mentre nel discreto
\[
z=0
\]
è un valore perfettamente lecito e corrisponde a un modo che si annulla in tempo finito.

\subsection{Interpretazione nel caso SISO}

Nel caso SISO, un sistema con tutti i poli in zero ha funzione di trasferimento del tipo
\[
G(z)=\frac{N(z)}{z^n}.
\]

Questa è la struttura tipica di un \textbf{filtro FIR} (\emph{Finite Impulse Response}),
cioè un sistema la cui risposta all'impulso si esaurisce in un tempo finito.

Per questo motivo un controllo dead-beat è molto rapido, ma presenta un problema pratico importante:
può richiedere azioni di controllo molto intense, quindi
\begin{itemize}
    \item si rischia di superare i limiti dell'attuatore;
    \item si può andare in saturazione;
    \item si perde la validità del modello lineare.
\end{itemize}

Quindi il dead-beat controller è concettualmente molto elegante,
ma spesso è troppo aggressivo per l'implementazione fisica.

\section{Sistemi stabilizzabili}

Si passa poi al caso in cui la coppia \((A,B)\) non sia completamente controllabile.

Consideriamo il sistema continuo
\[
\dot x=Ax+Bu.
\]

Se il sistema non è completamente controllabile, esiste un cambio di coordinate che separa
la parte controllabile da quella non controllabile:
\[
\dot z=
\begin{bmatrix}
A_c & A_{12}\\
0 & A_{nc}
\end{bmatrix}z
+
\begin{bmatrix}
B_c\\
0
\end{bmatrix}u,
\]
dove:
\begin{itemize}
    \item \(A_c\) descrive la parte controllabile;
    \item \(A_{nc}\) descrive la parte non controllabile.
\end{itemize}

Applicando una retroazione di stato
\[
u=-Kz,
\]
si può agire solo sulla parte controllabile, cioè sugli autovalori di
\[
A_c-B_cK_c.
\]

Gli autovalori della parte non controllabile, cioè quelli di \(A_{nc}\), \textbf{non possono essere spostati}.

\subsection{Definizione di stabilizzabilità}

Un sistema non completamente controllabile si dice \textbf{stabilizzabile} se tutti i modi non controllabili sono già stabili.

Nel caso tempo continuo:
\[
(A,B)\ \text{stabilizzabile}
\quad\Longleftrightarrow\quad
\Re(\lambda)<0\ \text{per tutti gli autovalori non controllabili}.
\]

Nel caso tempo discreto:
\[
(A,B)\ \text{stabilizzabile}
\quad\Longleftrightarrow\quad
|\lambda|<1\ \text{per tutti gli autovalori non raggiungibili/non controllabili}.
\]

Quindi:
\begin{itemize}
    \item la controllabilità completa permette di assegnare \emph{tutti} gli autovalori;
    \item la stabilizzabilità richiede solo che quelli non assegnabili siano già stabili.
\end{itemize}

\section{Perché la sola retroazione di stato non basta per il tracking perfetto}

Riprendiamo ora un sistema continuo
\[
\begin{cases}
\dot x=Ax+Bu,\\
y=Cx,
\end{cases}
\]
e consideriamo una legge di controllo del tipo
\[
u=-Kx+\bar r,
\]
dove \(\bar r\) rappresenta un riferimento costante.

Allora la dinamica diventa
\[
\dot x=(A-BK)x+B\bar r.
\]

All'equilibrio, imponendo \(\dot x=0\), si ottiene
\[
(A-BK)\bar x+B\bar r=0.
\]

Se la matrice
\[
N:=A-BK
\]
è invertibile, allora
\[
\bar x=-N^{-1}B\bar r.
\]

L'uscita a regime vale quindi
\[
\bar y=C\bar x=-CN^{-1}B\bar r.
\]

In generale \(\bar y\neq \bar r\), quindi la sola retroazione di stato \(-Kx\) non garantisce,
da sola, errore nullo a regime per un riferimento costante.

\section{Introduzione dell'integratore dell'errore}

Per eliminare l'errore a regime si introduce uno \textbf{stato integratore}.

Definiamo l'errore
\[
e(t)=r(t)-y(t)=r(t)-Cx(t),
\]
e introduciamo una nuova variabile di stato \(x_i(t)\) tale che
\[
\dot x_i(t)=e(t)=r(t)-Cx(t).
\]

Negli appunti compare prima la struttura in cui l'uscita dell'integratore entra come nuovo segnale di comando:
\[
u(t)=-Kx(t)+x_i(t).
\]

Sostituendo nella dinamica del sistema,
\[
\dot x=(A-BK)x+Bx_i.
\]

Quindi il sistema aumentato, con stato
\[
w=
\begin{bmatrix}
x\\
x_i
\end{bmatrix},
\]
si scrive
\[
\dot w
=
\begin{bmatrix}
A-BK & B\\
-C & 0
\end{bmatrix}
w
+
\begin{bmatrix}
0\\
I
\end{bmatrix}r.
\]

\subsection{Versione con ulteriore feedback sullo stato integratore}

Negli appunti compare poi una generalizzazione in cui si aggiunge un termine di retroazione anche sullo stato integratore:
\[
\dot x_i=r-Cx-K_i x_i.
\]

In questo caso il sistema aumentato diventa
\[
\dot w
=
\begin{bmatrix}
A-BK & B\\
-C & -K_i
\end{bmatrix}
w
+
\begin{bmatrix}
0\\
I
\end{bmatrix}r.
\]

L'idea è che anche la dinamica dell'integratore può essere modellata e ``chiusa''
attraverso un opportuno guadagno \(K_i\), così da permettere l'assegnamento
degli autovalori della matrice aumentata.

\section{Interpretazione dello stato aumentato}

L'introduzione di \(x_i\) ha due scopi:
\begin{itemize}
    \item memorizzare l'errore accumulato;
    \item forzare il sistema a correggere eventuali offset a regime.
\end{itemize}

Dal punto di vista del progetto, questo significa passare da un sistema di ordine \(n\)
a un sistema di ordine \(n+\ell\) (oppure \(n+1\) nel caso SISO),
su cui si può poi tentare una nuova retroazione di stato.

\section{Problema di osservazione dello stato}

La seconda grande parte della lezione introduce il problema dell'osservazione.

Finora si è sempre ipotizzato di conoscere tutto lo stato \(x(t)\),
ma nella pratica spesso sono misurabili solo:
\begin{itemize}
    \item l'ingresso \(u(t)\);
    \item l'uscita \(y(t)\).
\end{itemize}

Il sistema è
\[
\begin{cases}
\dot x=Ax+Bu,\\
y=Cx.
\end{cases}
\]

La domanda è:

\begin{quote}
\emph{Se lo stato non è direttamente accessibile, cosa posso ricostruire di \(x(t)\)
a partire da \(u(t)\), \(y(t)\) e dal modello \((A,B,C)\)?}
\end{quote}

\section{Primo tentativo: copia del modello}

Un primo tentativo consiste nel costruire una copia del sistema:
\[
\dot{\hat x}=A\hat x+Bu,
\]
dove \(\hat x(t)\) è una stima dello stato vero.

Se conoscessi esattamente:
\begin{itemize}
    \item il modello \(A,B\);
    \item la condizione iniziale \(x(0)\);
\end{itemize}
e scegliessi
\[
\hat x(0)=x(0),
\]
allora la stima coinciderebbe sempre con lo stato reale.

Infatti, definendo l'errore di stima
\[
e(t)=x(t)-\hat x(t),
\]
si ottiene
\[
\dot e=Ae.
\]

Se inizialmente \(e(0)=0\), allora \(e(t)\equiv 0\).

\subsection*{Problema pratico}

Questa idea però non è sufficiente nella realtà, perché:
\begin{itemize}
    \item non conosco esattamente la condizione iniziale;
    \item il modello può non essere perfetto;
    \item rumore e incertezze possono rendere l'errore non nullo.
\end{itemize}

Se ad esempio si usano matrici stimate \(\hat A,\hat B\), l'errore non soddisfa più una dinamica autonoma semplice,
e in generale non si può garantire che tenda a zero.

\section{Osservatore di Luenberger: idea}

Per correggere la stima si introduce una retroazione aggiuntiva basata sull'errore di uscita.

Poiché
\[
y=Cx,
\qquad
\hat y=C\hat x,
\]
la differenza
\[
y-\hat y = Cx-C\hat x = C(x-\hat x)=Ce
\]
contiene informazione sull'errore di stima.

Si costruisce allora un osservatore del tipo
\[
\dot{\hat x}=A\hat x+Bu+L\bigl(y-\hat y\bigr),
\]
cioè
\[
\dot{\hat x}=A\hat x+Bu+L\bigl(y-C\hat x\bigr),
\]
dove \(L\) è il guadagno dell'osservatore.

Negli appunti la stessa idea compare anche nella forma equivalente
\[
\dot{\hat x}=A\hat x+Bu-\mu(C\hat x-y),
\]
dove il simbolo del guadagno può cambiare, ma il significato resta lo stesso.

\section{Dinamica dell'errore di osservazione}

Definendo ancora
\[
e=x-\hat x,
\]
si ha
\[
\dot e
=
(Ax+Bu)-\bigl(A\hat x+Bu+L(y-C\hat x)\bigr).
\]

Poiché \(y=Cx\),
\[
\dot e=A(x-\hat x)-LC(x-\hat x),
\]
quindi
\[
\dot e=(A-LC)e.
\]

Questo è il punto chiave:
scegliendo \(L\), posso modificare gli autovalori della matrice
\[
A-LC.
\]

Quindi il problema dell'osservazione è duale a quello della retroazione di stato:
\begin{itemize}
    \item nella retroazione di stato si progettava \(K\) per agire su \(A-BK\);
    \item nell'osservatore si progetta \(L\) per agire su \(A-LC\).
\end{itemize}

\section{Commento finale sulla lezione}

La lezione mette insieme tre idee importanti:

\begin{enumerate}
    \item nel discreto è possibile progettare un \emph{dead-beat controller},
    assegnando tutti gli autovalori in zero;

    \item anche se un sistema non è completamente controllabile,
    può comunque essere \emph{stabilizzabile}, purché i modi non controllabili siano già stabili;

    \item quando la sola retroazione di stato non basta a garantire tracking perfetto,
    si introduce uno stato integratore e si passa a un sistema aumentato;

    \item quando lo stato non è interamente misurabile,
    nasce il problema dell'osservazione, che porta all'introduzione dell'osservatore dinamico.
\end{enumerate}

\section{Sintesi della lezione}

I punti principali sono:
\begin{itemize}
    \item per un sistema discreto
    \[
    x(k+1)=Ax(k)+Bu(k),
    \qquad
    u(k)=-Kx(k),
    \]
    è possibile assegnare gli autovalori del sistema chiuso in
    \[
    z=0,
    \]
    ottenendo un \textbf{dead-beat controller};

    \item il dead-beat controllo è molto rapido, ma può richiedere azioni di controllo troppo elevate;

    \item un sistema non completamente controllabile è \textbf{stabilizzabile}
    se i modi non controllabili sono già stabili;

    \item per eliminare l'errore a regime si introduce un integratore dell'errore,
    ottenendo un sistema aumentato;

    \item se lo stato non è direttamente misurabile, si introduce una stima \(\hat x\);

    \item una semplice copia del modello non basta in presenza di incertezze;

    \item usando l'errore di uscita \(y-\hat y\) si costruisce un osservatore dinamico:
    \[
    \dot{\hat x}=A\hat x+Bu+L(y-C\hat x);
    \]

    \item l'errore di osservazione evolve secondo
    \[
    \dot e=(A-LC)e,
    \]
    e quindi il progetto dell'osservatore consiste nello scegliere \(L\)
    in modo che questa dinamica sia stabile e sufficientemente veloce.
\end{itemize}
\chapter{Lezione 16 --- Osservatore asintotico, dualità e principio di separazione}

\section{Richiamo: osservatore dello stato}

Consideriamo il sistema lineare tempo-invariante
\[
\begin{cases}
\dot x = Ax + Bu,\\[4pt]
y = Cx,
\end{cases}
\]
dove:
\[
x(t)\in\mathbb{R}^n,\qquad
u(t)\in\mathbb{R}^m,\qquad
y(t)\in\mathbb{R}^\ell,
\]
e
\[
A\in\mathbb{R}^{n\times n},\qquad
B\in\mathbb{R}^{n\times m},\qquad
C\in\mathbb{R}^{\ell\times n}.
\]

Lo stato \(x(t)\) non è direttamente accessibile. Per questo si introduce uno
\textbf{stato osservato} o \textbf{stimato}
\[
\hat x(t),
\]
costruito mediante un osservatore dinamico.

\section{Equazioni dell'osservatore}

L'osservatore considerato a lezione è del tipo di Luenberger:
\[
\begin{cases}
\dot{\hat x} = A\hat x + Bu + L(y-\hat y),\\[4pt]
\hat y = C\hat x,
\end{cases}
\]
dove
\[
L\in\mathbb{R}^{n\times \ell}
\]
è la matrice di guadagno dell'osservatore.

L'idea è semplice:
\begin{itemize}
    \item copio la dinamica del sistema;
    \item confronto l'uscita vera \(y\) con l'uscita stimata \(\hat y\);
    \item correggo la stima mediante il termine \(L(y-\hat y)\).
\end{itemize}

\section{Errore di osservazione}

Definiamo l'errore di osservazione:
\[
e := x-\hat x.
\]

Allora
\[
\dot e = \dot x - \dot{\hat x}.
\]

Sostituendo le equazioni del sistema e dell'osservatore:
\[
\dot e
=
(Ax+Bu)-\bigl(A\hat x+Bu+L(y-\hat y)\bigr).
\]

Poiché
\[
y-\hat y = Cx-C\hat x = C(x-\hat x)=Ce,
\]
si ottiene
\[
\dot e = A(x-\hat x)-LC(x-\hat x),
\]
cioè
\[
\dot e = (A-LC)e.
\]

Questa è la \textbf{dinamica dell'errore di osservazione}.

\section{Osservatore asintotico}

Se tutti gli autovalori della matrice
\[
A-LC
\]
hanno parte reale strettamente negativa, allora
\[
e(t)\to 0 \qquad \text{per } t\to+\infty.
\]

Quindi
\[
\hat x(t)\to x(t).
\]

Per questo motivo l'equazione dell'osservatore viene chiamata
\textbf{osservatore asintotico}.

\subsection*{Interpretazione}

Non imponiamo che \(\hat x(t)\) coincida subito con \(x(t)\),
ma facciamo in modo che converga asintoticamente allo stato vero.

\section{Scelta di \(L\) e assegnamento degli autovalori dell'osservatore}

La matrice \(L\) è scelta da noi.  
Quindi il problema è del tutto analogo a quello già visto per la retroazione di stato:
\begin{itemize}
    \item lì si progettava \(K\) per modificare la matrice \(A-BK\);
    \item qui si progetta \(L\) per modificare la matrice \(A-LC\).
\end{itemize}

Vale allora il seguente risultato.

\begin{theorem}
Se la coppia \((A,C)\) è completamente osservabile, allora esiste una matrice
\[
L\in\mathbb{R}^{n\times \ell}
\]
tale da assegnare arbitrariamente gli autovalori della matrice
\[
A-LC.
\]
\end{theorem}

Quindi, sotto ipotesi di osservabilità completa, possiamo scegliere la velocità
di convergenza dell'osservatore.

\section{Dualità tra osservabilità e controllabilità}

Negli appunti si sottolinea che \textbf{osservabilità e controllabilità sono concetti duali}.

Dato il sistema
\[
\begin{cases}
\dot x = Ax + Bu,\\
y = Cx,
\end{cases}
\]
il suo sistema duale è
\[
\begin{cases}
\dot z = A^T z + C^T w,\\[4pt]
\rho = B^T z.
\end{cases}
\]

La proprietà chiave è:
\[
(A,C)\ \text{completamente osservabile}
\quad\Longleftrightarrow\quad
(A^T,C^T)\ \text{completamente controllabile}.
\]

Allora, applicando al sistema duale il teorema di assegnamento degli autovalori
mediante retroazione completa dello stato, sappiamo che esiste una matrice \(K_d\)
tale che gli autovalori di
\[
A^T-C^T K_d
\]
siano arbitrari.

Ma
\[
A^T-C^T K_d = (A-K_d^T C)^T,
\]
e una matrice ha gli stessi autovalori della sua trasposta. Dunque
\[
\lambda(A^T-C^T K_d)=\lambda(A-K_d^T C).
\]

Questo mostra che possiamo porre
\[
L = K_d^T,
\]
e assegnare arbitrariamente gli autovalori di
\[
A-LC.
\]

Questa è la giustificazione teorica del progetto dell'osservatore tramite dualità.

\section{Schema a blocchi dell'osservatore}

Dal punto di vista strutturale:
\begin{itemize}
    \item il sistema vero evolve con stato \(x\) e uscita \(y\);
    \item l'osservatore evolve con stato \(\hat x\) e uscita \(\hat y=C\hat x\);
    \item la differenza \(y-\hat y\) viene amplificata da \(L\) e reiniettata nell'osservatore.
\end{itemize}

Un punto concettuale importante è che lo \textbf{stato osservato} \(\hat x\)
è sempre accessibile, perché l'osservatore è un algoritmo che costruiamo noi.
Quindi, anche se lo stato vero \(x\) non è misurabile, lo stato \(\hat x\) lo è.

\section{Retroazione sullo stato osservato}

Poiché \(\hat x\) è disponibile, possiamo usarlo nella legge di controllo:
\[
u=-K\hat x.
\]

Questa è la soluzione pratica al problema della non accessibilità dello stato:
invece di retroazionare lo stato vero, retroazioniamo il suo stimatore.

\section{Equazioni del sistema accoppiato}

Consideriamo quindi il sistema vero
\[
\begin{cases}
\dot x = Ax + Bu,\\
y = Cx,
\end{cases}
\]
e l'osservatore
\[
\begin{cases}
\dot{\hat x} = A\hat x + Bu + L(y-\hat y),\\
\hat y = C\hat x,
\end{cases}
\]
con legge di controllo
\[
u=-K\hat x.
\]

\subsection{Dinamica del sistema vero}

Sostituendo \(u=-K\hat x\) nella dinamica del sistema:
\[
\dot x = Ax - BK\hat x.
\]

Poiché
\[
\hat x = x-e,
\]
segue
\[
\dot x = A x - BK(x-e)
= (A-BK)x + BK\,e.
\]

\subsection{Dinamica dell'osservatore}

L'equazione dell'osservatore diventa
\[
\dot{\hat x}
=
A\hat x - BK\hat x + L(Cx-C\hat x),
\]
cioè
\[
\dot{\hat x}
=
(A-BK-LC)\hat x + LC\,x.
\]

\subsection{Dinamica dell'errore}

Calcoliamo ora di nuovo l'errore:
\[
e=x-\hat x.
\]

Derivando:
\[
\dot e = \dot x - \dot{\hat x}.
\]

Sostituendo le espressioni precedenti:
\[
\dot e
=
\bigl((A-BK)x + BK e\bigr)
-
\bigl((A-BK)\hat x + LC(x-\hat x)\bigr).
\]

Poiché \(x-\hat x=e\), si ottiene
\[
\dot e = (A-BK)e + BK e - LC e,
\]
e quindi
\[
\dot e = (A-LC)e.
\]

Questo è un risultato fondamentale:
la dinamica dell'errore di osservazione \textbf{non dipende da \(K\)},
ma solo da \(L\).

\section{Dinamica congiunta di stato ed errore}

Se come nuovo vettore di stato prendiamo
\[
\begin{bmatrix}
x\\
e
\end{bmatrix},
\]
la dinamica complessiva diventa
\[
\begin{bmatrix}
\dot x\\
\dot e
\end{bmatrix}
=
\underbrace{
\begin{bmatrix}
A-BK & BK\\
0 & A-LC
\end{bmatrix}
}_{D}
\begin{bmatrix}
x\\
e
\end{bmatrix}.
\]

La matrice \(D\) è \textbf{triangolare a blocchi}.
Quindi il suo polinomio caratteristico è
\[
p_D(\lambda)
=
\det(\lambda I-D)
=
\det\!\bigl(\lambda I-(A-BK)\bigr)\,
\det\!\bigl(\lambda I-(A-LC)\bigr).
\]

Ne segue che gli autovalori di \(D\) sono l'unione di:
\begin{itemize}
    \item autovalori di \(A-BK\), legati alla dinamica del controllo;
    \item autovalori di \(A-LC\), legati alla dinamica dell'osservatore.
\end{itemize}

\section{Principio di separazione}

Il risultato precedente è il \textbf{principio di separazione}.

\begin{quote}
Gli autovalori del controllore (\(A-BK\)) e quelli dell'osservatore (\(A-LC\))
sono assegnabili \emph{indipendentemente}.
\end{quote}

Questo significa che:
\begin{itemize}
    \item posso progettare \(K\) come se avessi accesso diretto allo stato;
    \item posso progettare \(L\) come se stessi costruendo soltanto l'osservatore;
    \item il sistema complessivo ha come autovalori l'unione dei due insiemi.
\end{itemize}

\subsection*{Attenzione}

Il principio di separazione \textbf{non} significa che la dinamica effettiva di \(x\)
sia identica al caso in cui si retroaziona direttamente lo stato vero.

Infatti nella dinamica di \(x\) compare il termine
\[
BK\,e,
\]
che agisce come perturbazione finché l'errore di osservazione non è ancora nullo.

Quindi:
\begin{itemize}
    \item i \emph{modi} del sistema sono assegnabili separatamente;
    \item ma lo stato vero \(x\) è influenzato dall'errore \(e\) finché questo non si è estinto.
\end{itemize}

Se però
\[
e(t)\to 0,
\]
allora asintoticamente la retroazione sullo stato osservato si comporta come
la retroazione sullo stato vero.

\section{Scelta degli autovalori dell'osservatore}

Poiché vogliamo che l'errore di osservazione si annulli rapidamente,
conviene scegliere gli autovalori dell'osservatore più veloci di quelli del sistema controllato.

Nel caso continuo questo significa tipicamente scegliere gli autovalori di
\[
A-LC
\]
più a sinistra nel piano complesso rispetto a quelli di
\[
A-BK.
\]

In altre parole:
\begin{itemize}
    \item il sistema vero deve essere stabilizzato e avere la dinamica desiderata;
    \item l'osservatore deve convergere ancora più rapidamente,
    così l'errore \(e\) si annulla prima che domini la dinamica del sistema.
\end{itemize}

\subsection*{Osservazione pratica}

Negli appunti compare una nota importante:
\begin{itemize}
    \item sul sistema fisico non posso rendere arbitrariamente veloci gli autovalori, perché rischierei
    di richiedere segnali di controllo troppo grandi e saturare gli attuatori;
    \item per l'osservatore questo problema è meno critico, perché è un algoritmo interno.
\end{itemize}

Quindi in genere si scelgono gli autovalori dell'osservatore significativamente più veloci
di quelli del controllore.

\section{Interpretazione strutturale}

Dal punto di vista concettuale, il sistema completo è costituito da due sottosistemi:
\begin{enumerate}
    \item il \textbf{sistema vero}, con stato \(x\);
    \item il \textbf{sistema osservato}, con stato \(\hat x\).
\end{enumerate}

L'osservatore utilizza:
\begin{itemize}
    \item l'ingresso \(u\), che è noto;
    \item l'uscita misurata \(y\), che è nota;
    \item il proprio stato interno \(\hat x\), che è accessibile.
\end{itemize}

Una volta che \(\hat x\) converge a \(x\), posso trattare la retroazione
\[
u=-K\hat x
\]
come un sostituto pratico di
\[
u=-Kx.
\]

\section{Sintesi della lezione}

I punti principali della lezione sono:
\begin{enumerate}
    \item l'osservatore asintotico è descritto da
    \[
    \dot{\hat x}=A\hat x+Bu+L(y-\hat y),
    \qquad
    \hat y=C\hat x;
    \]

    \item definendo
    \[
    e=x-\hat x,
    \]
    la dinamica dell'errore è
    \[
    \dot e=(A-LC)e;
    \]

    \item se tutti gli autovalori di \(A-LC\) hanno parte reale negativa,
    allora
    \[
    e(t)\to 0
    \quad\text{e quindi}\quad
    \hat x(t)\to x(t);
    \]

    \item se la coppia \((A,C)\) è completamente osservabile, gli autovalori di \(A-LC\)
    si possono assegnare arbitrariamente;

    \item questo risultato deriva per dualità dalla teoria della retroazione di stato;

    \item usando la retroazione sullo stato osservato
    \[
    u=-K\hat x,
    \]
    il sistema vero soddisfa
    \[
    \dot x=(A-BK)x + BK\,e;
    \]

    \item la dinamica complessiva di \((x,e)\) è
    \[
    \begin{bmatrix}
    \dot x\\
    \dot e
    \end{bmatrix}
    =
    \begin{bmatrix}
    A-BK & BK\\
    0 & A-LC
    \end{bmatrix}
    \begin{bmatrix}
    x\\
    e
    \end{bmatrix};
    \]

    \item gli autovalori del sistema complessivo sono l'unione di quelli di \(A-BK\)
    e di quelli di \(A-LC\);

    \item controllore e osservatore si progettano indipendentemente:
    questo è il principio di separazione;

    \item in pratica conviene scegliere gli autovalori dell'osservatore più veloci
    di quelli del sistema controllato.
\end{enumerate}
\chapter{Lezione 17 --- Esercizio su risposta al gradino, realizzazione minima e assegnamento dei poli}

\section{Esercizio: studio di una funzione di trasferimento discreta}

Consideriamo la funzione di trasferimento
\[
G(z)=\frac{z^2}{\left(z-\frac14\right)\left[\left(z-\frac12\right)^2+\frac14\right]}.
\]

Osserviamo subito che
\[
\left(z-\frac12\right)^2+\frac14
=
z^2-z+\frac12,
\]
quindi
\[
G(z)=\frac{z^2}{\left(z-\frac14\right)\left(z^2-z+\frac12\right)}.
\]

\subsection{Poli e modi del sistema}

I poli sono:
\[
z_1=\frac14,
\qquad
z_{2,3}=\frac12\pm j\,\frac12.
\]

Per i poli complessi coniugati conviene passare alla forma polare:
\[
r=\sqrt{\left(\frac12\right)^2+\left(\frac12\right)^2}
=\frac{\sqrt{2}}{2},
\qquad
\theta=\frac{\pi}{4}.
\]

Infatti
\[
\frac12\pm j\frac12
=
\frac{\sqrt2}{2}e^{\pm j\pi/4}.
\]

Quindi i modi naturali associati sono:
\[
\left(\frac14\right)^k,
\qquad
\left(\frac{\sqrt2}{2}\right)^k \cos\!\left(\frac{k\pi}{4}\right),
\qquad
\left(\frac{\sqrt2}{2}\right)^k \sin\!\left(\frac{k\pi}{4}\right).
\]

\section{Risposta al gradino}

Vogliamo calcolare la risposta al gradino unitario.

Nel dominio z:
\[
U(z)=\mathcal{Z}\{1(k)\}=\frac{z}{z-1}.
\]

Dunque
\[
Y(z)=G(z)U(z)
=
\frac{z^3}{(z-1)\left(z-\frac14\right)\left(z^2-z+\frac12\right)}.
\]

Per l'antitrasformata conviene scomporre in fratti semplici
\[
\frac{Y(z)}{z}
=
\frac{A}{z-1}
+
\frac{B}{z-\frac14}
+
C\,\frac{z-r\cos\theta}{z^2-2r\cos\theta\,z+r^2}
+
D\,\frac{r\sin\theta}{z^2-2r\cos\theta\,z+r^2}.
\]

Nel nostro caso
\[
r=\frac{\sqrt2}{2},
\qquad
\theta=\frac{\pi}{4},
\qquad
r\cos\theta=\frac12,
\qquad
r\sin\theta=\frac12,
\]
quindi
\[
\frac{Y(z)}{z}
=
\frac{A}{z-1}
+
\frac{B}{z-\frac14}
+
C\,\frac{z-\frac12}{z^2-z+\frac12}
+
D\,\frac{\frac12}{z^2-z+\frac12}.
\]

\subsection{Calcolo dei coefficienti}

Per i termini del primo ordine si può usare il teorema di Heaviside.

\paragraph{Coefficiente \(A\)}
\[
A=\lim_{z\to 1}\frac{Y(z)}{z}(z-1)
=
\frac{1}{\left(1-\frac14\right)\left(1-1+\frac12\right)}
=
\frac{1}{\frac34\cdot \frac12}
=
\frac{8}{3}.
\]

\paragraph{Coefficiente \(B\)}
\[
B=\lim_{z\to \frac14}\frac{Y(z)}{z}\left(z-\frac14\right)
=
\frac{1}{\left(\frac14-1\right)\left(\frac1{16}-\frac14+\frac12\right)}.
\]

Poiché
\[
\frac1{16}-\frac14+\frac12=\frac{5}{16},
\]
segue
\[
B
=
\frac{1}{-\frac34\cdot \frac5{16}}
=
-\frac{64}{15}
\cdot \frac{1}{16}
=
-\frac{4}{15}.
\]

\paragraph{Coefficienti \(C\) e \(D\)}

Scriviamo
\[
\frac{Y(z)}{z}
=
\frac{A}{z-1}
+
\frac{B}{z-\frac14}
+
C\,\frac{z-\frac12}{z^2-z+\frac12}
+
D\,\frac{\frac12}{z^2-z+\frac12}.
\]

Moltiplicando per il denominatore comune e uguagliando i coefficienti si ottiene:
\[
C=-\frac{12}{5},
\qquad
D=-\frac{4}{5}.
\]

\subsection{Antitrasformata z}

Usiamo le coppie note:
\[
\mathcal{Z}^{-1}\left\{\frac{z}{z-a}\right\}=a^k\,1(k),
\]
\[
\mathcal{Z}^{-1}\left\{
\frac{z(z-r\cos\theta)}{z^2-2r\cos\theta\,z+r^2}
\right\}
=
r^k\cos(k\theta)\,1(k),
\]
\[
\mathcal{Z}^{-1}\left\{
\frac{zr\sin\theta}{z^2-2r\cos\theta\,z+r^2}
\right\}
=
r^k\sin(k\theta)\,1(k).
\]

Moltiplicando di nuovo per \(z\), si ha
\[
Y(z)
=
\frac{8}{3}\frac{z}{z-1}
-
\frac{4}{15}\frac{z}{z-\frac14}
-
\frac{12}{5}\frac{z\left(z-\frac12\right)}{z^2-z+\frac12}
-
\frac{4}{5}\frac{z\cdot \frac12}{z^2-z+\frac12}.
\]

Quindi
\[
y(k)
=
\frac{8}{3}\,1(k)
-
\frac{4}{15}\left(\frac14\right)^k 1(k)
-
\frac{12}{5}\left(\frac{\sqrt2}{2}\right)^k
\cos\!\left(\frac{k\pi}{4}\right)1(k)
-
\frac{4}{5}\left(\frac{\sqrt2}{2}\right)^k
\sin\!\left(\frac{k\pi}{4}\right)1(k).
\]

Questa è la risposta al gradino.

\section{Realizzazione minima in equazioni di stato}

Vogliamo ora costruire una realizzazione minima del sistema, cioè del tipo
\[
\begin{cases}
x(k+1)=Ax(k)+Bu(k),\\
y(k)=Cx(k).
\end{cases}
\]

\subsection{Denominatore come polinomio}

Espandiamo il denominatore:
\[
\left(z-\frac14\right)\left(z^2-z+\frac12\right)
=
z^3-\frac54 z^2+\frac34 z-\frac18.
\]

Allora
\[
G(z)=\frac{z^2}{z^3-\frac54 z^2+\frac34 z-\frac18}.
\]

Poiché
\[
Y(z)=G(z)U(z),
\]
segue
\[
\left(z^3-\frac54 z^2+\frac34 z-\frac18\right)Y(z)=z^2U(z).
\]

\subsection{Equazione alle differenze}

Antitrasformando si può scrivere, ad esempio,
\[
y(k)-\frac54 y(k-1)+\frac34 y(k-2)-\frac18 y(k-3)=u(k-1).
\]

Questa equazione suggerisce una realizzazione canonica di controllo di ordine \(3\).

\subsection{Forma canonica di controllo}

Poiché il denominatore ha grado \(3\), la dimensione minima dello stato è \(3\).  
Una realizzazione minima è quindi:
\[
x(k+1)=
\begin{bmatrix}
0 & 1 & 0\\
0 & 0 & 1\\
\frac18 & -\frac34 & \frac54
\end{bmatrix}x(k)
+
\begin{bmatrix}
0\\
0\\
1
\end{bmatrix}u(k),
\]
\[
y(k)=
\begin{bmatrix}
0 & 0 & 1
\end{bmatrix}x(k).
\]

\section{Osservazione sulle realizzazioni minime}

Dato un sistema in forma di stato, esiste un'unica funzione di trasferimento associata.

Viceversa, data una funzione di trasferimento, esistono in generale infiniti sistemi in forma di stato che la realizzano.

Tra tutte queste realizzazioni, si chiamano \textbf{realizzazioni minime} quelle con dimensione dello stato minima possibile.

Una realizzazione minima è, per costruzione, completamente:
\begin{itemize}
    \item raggiungibile (o controllabile);
    \item osservabile.
\end{itemize}

In particolare, i poli della funzione di trasferimento corrispondono agli autovalori raggiungibili e osservabili del sistema.

Quindi, una volta ottenuta una realizzazione minima, sono soddisfatte le ipotesi necessarie
per applicare le tecniche di assegnamento dei poli mediante retroazione dello stato.

\section{Progetto in ciclo chiuso: scelta dei poli desiderati}

Il sistema ha un polo reale in
\[
\frac14
\]
e una coppia di poli complessi coniugati in
\[
\frac12\pm j\frac12.
\]

Non essendoci specifiche particolari, scegliamo poli desiderati:
\begin{itemize}
    \item reali;
    \item più vicini a zero;
    \item quindi più veloci dei poli originari.
\end{itemize}

Una scelta semplice è:
\[
z_1^\star=\frac15,
\qquad
z_2^\star=\frac1{10},
\qquad
z_3^\star=\frac1{20}.
\]

\subsection{Polinomio caratteristico desiderato}

Il polinomio caratteristico desiderato è
\[
p_c^\star(\lambda)
=
\left(\lambda-\frac15\right)
\left(\lambda-\frac1{10}\right)
\left(\lambda-\frac1{20}\right).
\]

Sviluppando:
\[
p_c^\star(\lambda)
=
\lambda^3
-\frac{7}{20}\lambda^2
+\frac{7}{200}\lambda
-\frac{1}{1000}.
\]

\section{Retroazione completa dello stato}

Introduciamo una retroazione completa dello stato del tipo
\[
u(k)=-H\,x(k),
\qquad
H=
\begin{bmatrix}
h_1 & h_2 & h_3
\end{bmatrix}.
\]

La dinamica in ciclo chiuso diventa
\[
x(k+1)=(A-BH)x(k).
\]

Poiché
\[
B=
\begin{bmatrix}
0\\0\\1
\end{bmatrix},
\]
si ha
\[
BH=
\begin{bmatrix}
0\\0\\1
\end{bmatrix}
\begin{bmatrix}
h_1 & h_2 & h_3
\end{bmatrix}
=
\begin{bmatrix}
0 & 0 & 0\\
0 & 0 & 0\\
h_1 & h_2 & h_3
\end{bmatrix}.
\]

Quindi
\[
A-BH=
\begin{bmatrix}
0 & 1 & 0\\
0 & 0 & 1\\
\frac18-h_1 & -\frac34-h_2 & \frac54-h_3
\end{bmatrix}.
\]

Essendo ancora in forma canonica di controllo, il suo polinomio caratteristico è
\[
\lambda^3
+\left(-\frac54+h_3\right)\lambda^2
+\left(\frac34+h_2\right)\lambda
+\left(-\frac18+h_1\right).
\]

Imponiamo l'uguaglianza con il polinomio desiderato:
\[
\lambda^3
+\left(-\frac54+h_3\right)\lambda^2
+\left(\frac34+h_2\right)\lambda
+\left(-\frac18+h_1\right)
=
\lambda^3-\frac{7}{20}\lambda^2+\frac{7}{200}\lambda-\frac{1}{1000}.
\]

Da cui:
\[
-\frac54+h_3=-\frac{7}{20},
\qquad
\frac34+h_2=\frac{7}{200},
\qquad
-\frac18+h_1=-\frac{1}{1000}.
\]

Si ottiene quindi
\[
h_3=\frac{9}{10},
\qquad
h_2=-\frac{143}{200},
\qquad
h_1=\frac{31}{250}.
\]

Dunque la matrice di guadagno è
\[
H=
\begin{bmatrix}
\frac{31}{250} & -\frac{143}{200} & \frac{9}{10}
\end{bmatrix}.
\]

\section{Commento finale}

Questa lezione mette insieme tre aspetti importanti.

\begin{enumerate}
    \item Data una funzione di trasferimento discreta, si possono ricavare:
    \begin{itemize}
        \item poli;
        \item modi naturali;
        \item risposta al gradino;
        \item una realizzazione minima in forma di stato.
    \end{itemize}

    \item Una realizzazione minima è la più conveniente per il progetto, perché è:
    \begin{itemize}
        \item completamente raggiungibile;
        \item completamente osservabile.
    \end{itemize}

    \item Se la realizzazione è in forma canonica di controllo, l'assegnamento dei poli in ciclo chiuso
    è particolarmente semplice: basta confrontare i coefficienti del polinomio caratteristico.
\end{enumerate}

\section{Sintesi della lezione}

I risultati principali sono:
\begin{itemize}
    \item i poli di
    \[
    G(z)=\frac{z^2}{\left(z-\frac14\right)\left(z^2-z+\frac12\right)}
    \]
    sono
    \[
    \frac14,\qquad \frac12\pm j\frac12;
    \]

    \item la risposta al gradino è
    \[
    y(k)
    =
    \frac{8}{3}\,1(k)
    -
    \frac{4}{15}\left(\frac14\right)^k 1(k)
    -
    \frac{12}{5}\left(\frac{\sqrt2}{2}\right)^k
    \cos\!\left(\frac{k\pi}{4}\right)1(k)
    -
    \frac{4}{5}\left(\frac{\sqrt2}{2}\right)^k
    \sin\!\left(\frac{k\pi}{4}\right)1(k);
    \]

    \item una realizzazione minima è
    \[
    x(k+1)=
    \begin{bmatrix}
    0 & 1 & 0\\
    0 & 0 & 1\\
    \frac18 & -\frac34 & \frac54
    \end{bmatrix}x(k)
    +
    \begin{bmatrix}
    0\\0\\1
    \end{bmatrix}u(k),
    \qquad
    y(k)=
    \begin{bmatrix}
    0 & 0 & 1
    \end{bmatrix}x(k);
    \]

    \item scegliendo i poli desiderati
    \[
    \frac15,\quad \frac1{10},\quad \frac1{20},
    \]
    il polinomio caratteristico desiderato è
    \[
    \lambda^3-\frac{7}{20}\lambda^2+\frac{7}{200}\lambda-\frac{1}{1000};
    \]

    \item la retroazione completa dello stato
    \[
    u(k)=-H x(k)
    \]
    con
    \[
    H=
    \begin{bmatrix}
    \frac{31}{250} & -\frac{143}{200} & \frac{9}{10}
    \end{bmatrix}
    \]
    assegna esattamente tali poli al sistema in ciclo chiuso.
\end{itemize}
\chapter{Lezione 18 --- Controllo ottimo e programmazione dinamica}

\section{Introduzione al controllo ottimo}

Vogliamo impostare un problema di controllo in retroazione che non si limiti a
stabilizzare il sistema o a imporre certi poli, ma che soddisfi un
\textbf{criterio di ottimalità}.

L'idea è minimizzare una funzione di costo che tenga conto contemporaneamente di:
\begin{itemize}
    \item andamento della traiettoria di stato;
    \item sforzo di controllo;
    \item eventuali vincoli o penalità sullo stato finale.
\end{itemize}

Il controllo ottimo è particolarmente importante per sistemi MIMO e per problemi
in cui non basta ``chiudere bene il sistema'', ma bisogna anche farlo nel modo
più conveniente possibile.

Un'applicazione moderna di questa idea è il \textbf{Model Predictive Control} (MPC),
che può essere visto come uno sviluppo del controllo ottimo.

\section{Due approcci al controllo ottimo}

Negli appunti vengono richiamati due approcci classici:
\begin{enumerate}
    \item \textbf{Calcolo delle variazioni / principio di Pontryagin};
    \item \textbf{Programmazione dinamica / principio di Bellman}.
\end{enumerate}

Entrambi descrivono lo stesso problema di ottimo, ma con strumenti diversi.
In questa lezione ci concentriamo sulla \textbf{programmazione dinamica} per sistemi
a tempo discreto.

\section{Principio di Bellman e idea del cost-to-go}

Per introdurre il metodo si considera il classico esempio del percorso ottimo:
si deve andare da un punto iniziale \(x_A\) a un punto finale \(x_B\), attraversando
eventuali punti di transito organizzati per \emph{stadi}.

Ogni arco del grafo ha un costo, che può rappresentare:
\begin{itemize}
    \item tempo di percorrenza;
    \item consumo di carburante;
    \item rischio;
    \item una qualsiasi grandezza da minimizzare.
\end{itemize}

Il \textbf{principio di Bellman} dice che:

\begin{quote}
Se un percorso è ottimo, allora ogni suo tratto finale è a sua volta ottimo
rispetto al punto da cui parte.
\end{quote}

Questo porta al concetto di \textbf{cost-to-go}:

\begin{quote}
il cost-to-go è il costo minimo necessario per andare dallo stato attuale fino
alla destinazione finale.
\end{quote}

La strategia consiste quindi nel lavorare \emph{all'indietro}, dallo stadio finale
verso lo stadio iniziale.

\section{Problema dinamico discreto generale}

Consideriamo un sistema discreto generico
\[
x(i+1)=f(x(i),u(i)),
\]
e una funzione di costo additiva del tipo
\[
J=\sum_{i=0}^{N-1} h_i(x(i),u(i)).
\]

Il problema è:

\begin{quote}
determinare la sequenza di controlli
\[
u(0),u(1),\dots,u(N-1)
\]
che minimizza il costo \(J\), compatibilmente con la dinamica del sistema.
\end{quote}

\subsection{Definizione di cost-to-go ottimo}

Si definisce il costo ottimo a partire dallo stadio \(i\) come
\[
J_i^o(x(i))
=
\min_{u(i),u(i+1),\dots,u(N-1)}
\sum_{k=i}^{N-1} h_k(x(k),u(k)).
\]

In particolare:
\[
J_0^o(x(0))=\min J.
\]

\subsection{Equazione di Bellman}

Poiché il costo è additivo, vale la ricorsione
\[
J_i^o(x(i))
=
\min_{u(i)}
\left[
h_i(x(i),u(i))
+
J_{i+1}^o\bigl(f(x(i),u(i))\bigr)
\right].
\]

Questa è la \textbf{equazione di Bellman}.

L'inizializzazione è data dalla condizione terminale. Nel caso più semplice:
\[
J_N^o(x(N))=0,
\]
oppure, più in generale, si può assegnare una penalità terminale.

\subsection{Osservazione}

In generale non esiste una soluzione analitica esplicita del problema.
La strategia ottima dipende dallo stato in cui il sistema si trova, quindi il controllo
ottimo è in generale una legge di retroazione
\[
u(i)=\mu_i(x(i)).
\]

\section{Caso particolare: controllo ottimo LQ}

Un caso molto importante, e in effetti uno dei pochi con soluzione analitica,
è il problema \textbf{LQ}:
\begin{itemize}
    \item \textbf{L} = sistema lineare;
    \item \textbf{Q} = costo quadratico.
\end{itemize}

\subsection{Sistema}

Consideriamo il sistema lineare discreto
\[
x_{k+1}=Ax_k+Bu_k,
\qquad
x_0=\hat x,
\]
dove
\[
x_k\in\mathbb{R}^n,\qquad
u_k\in\mathbb{R}^m.
\]

\subsection{Indice di costo}

Il costo è scelto nella forma
\[
J
=
\sum_{k=0}^{N-1}
\left(
x_k^T V x_k + u_k^T P u_k
\right)
+
x_N^T V_N x_N,
\]
dove:
\begin{itemize}
    \item \(V=V^T\succeq 0\): penalità sullo stato;
    \item \(P=P^T\succ 0\): penalità sul controllo;
    \item \(V_N=V_N^T\succeq 0\): penalità sullo stato finale.
\end{itemize}

L'ultimo termine penalizza lo stato di arrivo al tempo \(N\).

\section{Forma attesa della soluzione}

Cerchiamo una legge ottima del tipo
\[
u_k^o=-H_k x_k,
\]
cioè una retroazione lineare sullo stato.

La cosa non è banale a priori, ma la programmazione dinamica mostrerà che,
nel caso LQ, questa forma è effettivamente ottima.

\section{Richiamo utile: minimo di una forma quadratica}

Negli appunti compare il seguente fatto elementare.

Se
\[
Q\succ 0,
\]
allora la funzione
\[
\phi(z)=z^TQz+2c^Tz
\]
ha minimo unico nel punto
\[
z^\star=-Q^{-1}c.
\]

Questo risultato verrà usato ripetutamente nella minimizzazione rispetto a \(u_k\).

\section{Passo finale della ricorsione: stadio \(N-1\)}

Partiamo dallo stadio finale utile, cioè \(k=N-1\).

Poiché
\[
x_N=Ax_{N-1}+Bu_{N-1},
\]
si ha
\[
J_{N-1}^o(x_{N-1})
=
\min_{u_{N-1}}
\left[
x_{N-1}^T V x_{N-1}
+
u_{N-1}^T P u_{N-1}
+
x_N^T V_N x_N
\right].
\]

Sostituendo \(x_N=Ax_{N-1}+Bu_{N-1}\):
\[
\begin{aligned}
J_{N-1}^o(x_{N-1})
&=
\min_{u_{N-1}}
\Big[
x_{N-1}^T V x_{N-1}
+
u_{N-1}^T P u_{N-1}\\
&\qquad\qquad
+
(Ax_{N-1}+Bu_{N-1})^T
V_N
(Ax_{N-1}+Bu_{N-1})
\Big].
\end{aligned}
\]

Sviluppando:
\[
\begin{aligned}
J_{N-1}^o(x_{N-1})
&=
x_{N-1}^T(V+A^TV_NA)x_{N-1}\\
&\quad+
\min_{u_{N-1}}
\Big[
u_{N-1}^T(P+B^TV_NB)u_{N-1}
+
2x_{N-1}^T A^TV_NBu_{N-1}
\Big].
\end{aligned}
\]

Ponendo
\[
Q=P+B^TV_NB,
\qquad
c=B^TV_NA\,x_{N-1},
\]
e usando il risultato precedente, si ottiene
\[
u_{N-1}^o(x_{N-1})
=
-(P+B^TV_NB)^{-1}B^TV_NA\,x_{N-1}.
\]

Definendo
\[
H_{N-1}
=
(P+B^TV_NB)^{-1}B^TV_NA,
\]
si ha
\[
u_{N-1}^o=-H_{N-1}x_{N-1}.
\]

Sostituendo il valore ottimo, il costo ottimo allo stadio \(N-1\) assume forma quadratica:
\[
J_{N-1}^o(x_{N-1})=x_{N-1}^T T_{N-1}x_{N-1},
\]
con
\[
T_{N-1}
=
V+A^TV_NA
-
A^TV_NB(P+B^TV_NB)^{-1}B^TV_NA.
\]

\section{Passo ricorsivo generale}

Supponiamo ora che, a uno stadio generico \(k+1\), il cost-to-go ottimo sia già noto
nella forma quadratica
\[
J_{k+1}^o(x_{k+1})=x_{k+1}^T T_{k+1}x_{k+1}.
\]

Allora allo stadio \(k\):
\[
J_k^o(x_k)
=
\min_{u_k}
\left[
x_k^TVx_k+u_k^TPu_k+J_{k+1}^o(x_{k+1})
\right],
\]
con
\[
x_{k+1}=Ax_k+Bu_k.
\]

Sostituendo:
\[
\begin{aligned}
J_k^o(x_k)
&=
\min_{u_k}
\Big[
x_k^TVx_k
+
u_k^TPu_k
+
(Ax_k+Bu_k)^TT_{k+1}(Ax_k+Bu_k)
\Big].
\end{aligned}
\]

Sviluppando:
\[
\begin{aligned}
J_k^o(x_k)
&=
x_k^T(V+A^TT_{k+1}A)x_k\\
&\quad+
\min_{u_k}
\Big[
u_k^T(P+B^TT_{k+1}B)u_k
+
2x_k^TA^TT_{k+1}Bu_k
\Big].
\end{aligned}
\]

Usando di nuovo il minimo della forma quadratica:
\[
u_k^o(x_k)
=
-(P+B^TT_{k+1}B)^{-1}B^TT_{k+1}A\,x_k.
\]

Definendo
\[
H_k
=
(P+B^TT_{k+1}B)^{-1}B^TT_{k+1}A,
\]
si ottiene la legge ottima
\[
u_k^o=-H_kx_k.
\]

Sostituendo nel costo:
\[
J_k^o(x_k)=x_k^TT_kx_k,
\]
dove
\[
T_k
=
V
+
A^TT_{k+1}A
-
A^TT_{k+1}B
(P+B^TT_{k+1}B)^{-1}
B^TT_{k+1}A.
\]

\section{Ricorsione di Riccati discreta}

La ricorsione
\[
T_k
=
V
+
A^TT_{k+1}A
-
A^TT_{k+1}B
(P+B^TT_{k+1}B)^{-1}
B^TT_{k+1}A
\]
è l'\textbf{equazione di Riccati a tempo discreto} (all'indietro).

L'inizializzazione è
\[
T_N=V_N.
\]

Quindi l'algoritmo di soluzione del problema LQ su orizzonte finito è:

\begin{enumerate}
    \item fissare la condizione terminale
    \[
    T_N=V_N;
    \]
    \item calcolare all'indietro
    \[
    T_{N-1},T_{N-2},\dots,T_0;
    \]
    \item ad ogni passo ricavare
    \[
    H_k=(P+B^TT_{k+1}B)^{-1}B^TT_{k+1}A;
    \]
    \item applicare il controllo ottimo
    \[
    u_k^o=-H_kx_k.
    \]
\end{enumerate}

\section{Forma finale della soluzione}

Riassumendo, la soluzione del problema LQ discreto a orizzonte finito è:
\[
u_k^o=-H_kx_k,
\]
con
\[
H_k=(P+B^TT_{k+1}B)^{-1}B^TT_{k+1}A,
\]
e
\[
J_k^o(x_k)=x_k^TT_kx_k,
\]
dove \(T_k\) soddisfa la ricorsione di Riccati
\[
T_k
=
V
+
A^TT_{k+1}A
-
A^TT_{k+1}B
(P+B^TT_{k+1}B)^{-1}
B^TT_{k+1}A,
\qquad
T_N=V_N.
\]

\section{Osservazioni importanti}

\subsection{Controllo in retroazione}

La soluzione ottima non è una sequenza aperta di controlli, ma una \textbf{retroazione}
lineare sullo stato:
\[
u_k^o=-H_kx_k.
\]

Questa è una delle ragioni del successo del metodo:
il controllo ottimo è robusto rispetto a perturbazioni e incertezze sullo stato,
perché viene ricalcolato a ogni istante a partire da \(x_k\).

\subsection{Guadagno tempo-variante}

Anche se il sistema di partenza è tempo-invariante, il guadagno ottimo
\[
H_k
\]
dipende in generale da \(k\), perché dipende da \(T_{k+1}\) e quindi
dalla distanza dal tempo finale \(N\).

Dunque, su orizzonte finito, il controllo ottimo è in generale \textbf{tempo-variante}.

\subsection{Stabilità}

Negli appunti compare giustamente la domanda:

\begin{quote}
\emph{Ma siamo sicuri che il sistema sia stabile?}
\end{quote}

A questo stadio della trattazione la risposta non è ancora automatica:
la ricorsione fornisce il controllo ottimo su orizzonte finito, ma la stabilità
asintotica del sistema chiuso richiede un'analisi ulteriore.

Questo porterà naturalmente, nelle lezioni successive, al caso a orizzonte infinito,
in cui il guadagno tende a una matrice costante e compare l'equazione algebrica di Riccati.

\section{Interpretazione dei pesi \(V\), \(P\), \(V_N\)}

I tre pesi del problema hanno significato fisico chiaro:
\begin{itemize}
    \item \(V\): penalizza deviazioni dello stato lungo tutta la traiettoria;
    \item \(P\): penalizza lo sforzo di controllo;
    \item \(V_N\): penalizza lo stato finale.
\end{itemize}

Se aumento \(P\), il controllo ottimo diventa più prudente:
si cerca di usare meno attuazione.

Se aumento \(V\) o \(V_N\), si privilegia il rispetto della traiettoria e/o del target finale,
accettando anche uno sforzo di controllo maggiore.

\section{Sintesi della lezione}

I risultati principali della lezione sono:

\begin{itemize}
    \item il controllo ottimo si può affrontare con il principio di Bellman;
    \item si definisce il cost-to-go ottimo
    \[
    J_i^o(x(i));
    \]
    \item nel caso lineare-quadratico discreto:
    \[
    x_{k+1}=Ax_k+Bu_k,
    \qquad
    J=\sum_{k=0}^{N-1}(x_k^TVx_k+u_k^TPu_k)+x_N^TV_Nx_N;
    \]
    \item il cost-to-go ottimo è quadratico:
    \[
    J_k^o(x_k)=x_k^TT_kx_k;
    \]
    \item il controllo ottimo è lineare nello stato:
    \[
    u_k^o=-H_kx_k;
    \]
    \item il guadagno ottimo è
    \[
    H_k=(P+B^TT_{k+1}B)^{-1}B^TT_{k+1}A;
    \]
    \item la matrice \(T_k\) soddisfa la ricorsione di Riccati all'indietro:
    \[
    T_k
    =
    V
    +
    A^TT_{k+1}A
    -
    A^TT_{k+1}B
    (P+B^TT_{k+1}B)^{-1}
    B^TT_{k+1}A,
    \qquad
    T_N=V_N.
    \]
\end{itemize}
\chapter{Lezione 19 --- Riccati a regime, orizzonte infinito, MPC e programmazione dinamica in tempo continuo}

\section{Richiamo: controllo ottimo LQ a tempo discreto}

Consideriamo il sistema lineare discreto
\[
x_{k+1}=Ax_k+Bu_k,
\qquad x_0=\hat x,
\]
con funzionale di costo
\[
J=\sum_{k=0}^{N-1}\left(x_k^T V x_k+u_k^T P u_k\right)+x_N^T V_N x_N.
\]

Le matrici di costo soddisfano le ipotesi:
\[
V=V^T\succeq 0,
\qquad
P=P^T\succ 0,
\qquad
V_N=V_N^T\succeq 0.
\]

Dalle lezioni precedenti abbiamo dimostrato che il controllo ottimo ha la forma
\[
u_k^\circ=-H_k x_k,
\]
dove
\[
H_k=\left(P+B^T T_{k+1}B\right)^{-1}B^T T_{k+1}A,
\]
e il costo ottimo a partire dal passo \(k\) è
\[
J_k^\circ=x_k^T T_k x_k.
\]

La matrice \(T_k\) soddisfa la ricorsione all'indietro di Riccati:
\[
T_k
=
V
+
A^T
\left[
T_{k+1}
-
T_{k+1}B\left(P+B^T T_{k+1}B\right)^{-1}B^T T_{k+1}
\right]
A,
\]
con condizione terminale
\[
T_N=V_N.
\]

\section{Osservazioni sulla soluzione del problema LQ}

La soluzione del problema di controllo ottimo dipende dalla conoscenza del modello,
cioè delle matrici \(A\) e \(B\), e dalle matrici di peso \(V\), \(P\), \(V_N\).

Quindi, per poter implementare il controllo ottimo, è necessario:
\begin{itemize}
    \item disporre di un modello in equazioni di stato del sistema;
    \item scegliere le matrici di costo in base alle prestazioni desiderate;
    \item risolvere \emph{prima} dell'evoluzione del processo la ricorsione di Riccati.
\end{itemize}

In pratica, la ricorsione di Riccati viene risolta offline, partendo dall'istante finale
e andando all'indietro fino all'istante iniziale; una volta calcolati i guadagni
\(
H_k
\),
si può implementare online la legge di controllo
\[
u_k^\circ=-H_k x_k.
\]

\section{La ricorsione di Riccati e il concetto di regime}

Una particolarità della ricorsione di Riccati è che essa si risolve \emph{all'indietro}:
si parte da \(T_N\) e si ricavano \(T_{N-1},T_{N-2},\dots,T_0\).

Questo è diverso da ciò che accade per le normali equazioni di stato di un sistema stabile,
le cui soluzioni vengono tipicamente studiate in avanti nel tempo:
\[
x_0 \to x_1 \to x_2 \to \cdots
\]
e, nei sistemi asintoticamente stabili, si tende al regime per \(k\to+\infty\).

Nel caso di Riccati, essendo la ricorsione all'indietro, il concetto di regime va inteso
come \emph{punto fisso} della ricorsione:
\[
T_k=T_{k+1}=T.
\]

Sostituendo questa condizione nella ricorsione di Riccati si ottiene
l'\textbf{equazione algebrica di Riccati a tempo discreto}:
\[
T
=
V
+
A^T
\left[
T
-
TB\left(P+B^T T B\right)^{-1}B^T T
\right]A.
\]

Questa equazione non è più una ricorsione alle differenze, ma una
\textbf{equazione matriciale algebrica}.

\section{Orizzonte finito e orizzonte infinito}

L'orizzonte \(N\) ha un ruolo fondamentale.

\subsection{Orizzonte finito}

Se \(N\) è finito e fissato, si parla di \textbf{controllo ottimo su orizzonte finito}.
In questo caso:
\begin{itemize}
    \item bisogna calcolare una successione finita di matrici \(T_k\);
    \item si ottiene una successione finita di guadagni \(H_k\);
    \item il controllo ottimo è in generale \emph{tempo-variante}.
\end{itemize}

\subsection{Orizzonte infinito}

Se invece \(N\to+\infty\), si parla di \textbf{controllo ottimo su orizzonte infinito}.
In questo caso, sotto opportune ipotesi, la successione \(T_k\) converge a una matrice
costante \(T\), e quindi anche il guadagno converge a una matrice costante:
\[
H=\left(P+B^T T B\right)^{-1}B^T T A.
\]

Il controllo ottimo diventa allora
\[
u_k^\circ=-Hx_k,
\]
cioè una retroazione lineare a guadagno costante.

\subsection{Interpretazione}

Con orizzonte finito si ottimizza tenendo conto anche di come il sistema deve arrivare
al tempo finale fissato \(N\).

Con orizzonte infinito, invece, si rinuncia a un vincolo finale preciso:
si privilegia il comportamento complessivo su tempi lunghi.
In questo senso, l'ottimo a orizzonte infinito può essere meno raffinato sul transitorio
iniziale, ma produce una legge di controllo stazionaria molto più semplice da implementare.

\section{Implementazione pratica del controllo LQ}

Dal punto di vista implementativo:
\begin{enumerate}
    \item si costruisce il modello discreto del sistema;
    \item si scelgono le matrici \(V\), \(P\), \(V_N\);
    \item si risolve la ricorsione di Riccati all'indietro;
    \item si memorizzano i guadagni \(H_k\) oppure, nel caso a regime, il solo guadagno \(H\);
    \item durante il funzionamento si applica
    \[
    u_k^\circ=-H_kx_k
    \quad\text{oppure}\quad
    u_k^\circ=-Hx_k.
    \]
\end{enumerate}

\section{Dal controllo ottimo all'MPC}

L'idea dell'\textbf{MPC (Model Predictive Control)} nasce proprio dal controllo ottimo
su orizzonte finito.

In un regolatore MPC si fa, concettualmente, così:
\begin{enumerate}
    \item al tempo corrente si considera un orizzonte finito \(N\);
    \item si risolve il problema di controllo ottimo su tale orizzonte;
    \item si applica solo il primo tratto del controllo ottimo;
    \item si misura il nuovo stato;
    \item si sposta in avanti l'orizzonte e si ripete tutto da capo.
\end{enumerate}

Questa tecnica è detta anche \textbf{receding horizon control}.

\subsection{Idea operativa}

Negli appunti la logica è:
\begin{description}
    \item[Passo 1:] si linearizza il sistema attorno alle condizioni iniziali, si risolve
    Riccati su orizzonte \(N\), si applica il controllo LQ per un tempo \(H\ll N\);
    \item[Passo 2:] si osserva dove il sistema è arrivato, si rilinearizza attorno al nuovo
    punto di lavoro, si risolve di nuovo il problema LQ e si riapplica il controllo.
\end{description}

\section{Perché nasce l'MPC: impianti tempo-varianti}

Nella pratica gli impianti reali sono spesso:
\begin{itemize}
    \item non lineari;
    \item linearizzabili solo localmente;
    \item descrivibili da modelli che cambiano lungo la traiettoria.
\end{itemize}

Se il sistema evolve lungo una traiettoria, la linearizzazione locale può produrre matrici
tempo-varianti:
\[
A \to A_k,
\qquad
B \to B_k.
\]

In questo caso cambia anche la Riccati:
\[
T_k \to T_k \ \text{calcolata usando } A_k,B_k.
\]

Non esiste più, in generale, un'unica equazione algebrica di regime.
Bisogna invece risolvere una Riccati alle differenze con coefficienti variabili nel tempo.

Questo richiede di conoscere, a ogni istante, il modello locale del sistema.
Nella realtà questo è complesso, sia perché il modello è sempre approssimato,
sia perché possono agire disturbi e incertezze.

L'MPC è allora una soluzione pratica:
si ricalcola periodicamente il problema ottimo usando il modello locale più aggiornato.

\subsection{Condizione pratica}

Per fare funzionare questa strategia occorre che il processo non sia troppo veloce:
tra un aggiornamento e l'altro bisogna avere il tempo di
\begin{itemize}
    \item stimare o misurare lo stato;
    \item linearizzare il modello;
    \item risolvere il problema di Riccati;
    \item applicare il nuovo controllo.
\end{itemize}

\section{Programmazione dinamica nel caso a tempo continuo}

Passiamo ora al caso continuo, in forma più generale.

Consideriamo un impianto non lineare
\[
\dot x = F(x(t),u(t)),
\]
e supponiamo, per semplicità, che l'origine sia un equilibrio:
\[
F(0,0)=0.
\]

Vogliamo minimizzare, su orizzonte infinito, il funzionale di costo
\[
J=\int_{t_0}^{+\infty} \ell(x(\tau),u(\tau))\,d\tau.
\]

Siamo quindi interessati a trovare un controllo ottimo
\[
u^\circ(t)=\arg\min_{u(t)} J.
\]

\section{Richiamo del concetto di cost-to-go}

Nel tempo discreto avevamo introdotto il \textbf{cost-to-go}:
se un sistema si trova lungo una traiettoria ottima e arriva in un certo stato,
allora anche il tratto successivo della traiettoria deve restare ottimo.

Lo stesso concetto si usa qui.

Definiamo
\[
\psi(x_0)=J^\circ,
\]
cioè il costo ottimo associato alla condizione iniziale \(x_0\).

Più in generale, se al tempo \(t\) il sistema si trova nello stato \(x(t)\),
definiamo
\[
\psi(x(t))
=
\min_u \int_t^{+\infty}\ell(x(\tau),u(\tau))\,d\tau.
\]

\section{Scomposizione del costo su un intervallo breve}

Prendiamo un piccolo intervallo \(h>0\). Allora, per il principio di Bellman,
si può scrivere
\[
\psi(x(t))
=
\min_u
\left[
\int_t^{t+h}\ell(x(\tau),u(\tau))\,d\tau
+
\psi(x(t+h))
\right].
\]

Sottraendo \(\psi(x(t))\) da entrambi i membri:
\[
0
=
\min_u
\left[
\psi(x(t+h))-\psi(x(t))
+
\int_t^{t+h}\ell(x(\tau),u(\tau))\,d\tau
\right].
\]

Dividendo per \(h\):
\[
0
=
\min_u
\left[
\frac{\psi(x(t+h))-\psi(x(t))}{h}
+
\frac{1}{h}\int_t^{t+h}\ell(x(\tau),u(\tau))\,d\tau
\right].
\]

Passando al limite per \(h\to 0^+\), il secondo termine tende a
\[
\ell(x(t),u(t)),
\]
mentre il primo tende a
\[
\dot\psi(x(t)).
\]

Quindi otteniamo
\[
\min_u\left[\dot\psi(x(t))+\ell(x(t),u(t))\right]=0.
\]

\section{Regola della catena}

Poiché \(\psi\) dipende dallo stato e lo stato dipende dal tempo,
per la regola della catena:
\[
\dot\psi(x(t))
=
\nabla_x\psi(x(t))^T \dot x(t).
\]

Essendo
\[
\dot x(t)=F(x(t),u(t)),
\]
si ha
\[
\dot\psi(x(t))
=
\nabla_x\psi(x(t))^T F(x(t),u(t)).
\]

Sostituendo nella relazione precedente si ottiene la
\textbf{equazione di Bellman} (o, più precisamente, l'equazione di Hamilton--Jacobi--Bellman):
\[
\min_u
\left[
\nabla_x\psi(x)^T F(x,u)+\ell(x,u)
\right]
=0.
\]

Questa è una condizione necessaria che il costo ottimo deve soddisfare.

\section{Interpretazione dell'equazione di Bellman}

L'equazione
\[
\min_u
\left[
\nabla_x\psi(x)^T F(x,u)+\ell(x,u)
\right]
=0
\]
esprime il fatto che, in ogni stato \(x\), il controllo ottimo è quello che minimizza
la somma di:
\begin{itemize}
    \item costo istantaneo \(\ell(x,u)\);
    \item variazione istantanea del costo residuo \(\psi\) dovuta alla dinamica del sistema.
\end{itemize}

In altre parole, il controllo ottimo sceglie in ogni istante il miglior compromesso
tra costo immediato e costo futuro.

\section{Collegamento con il caso discreto}

Nel caso discreto avevamo:
\[
J_k^\circ(x_k)=\min_{u_k}\Bigl[h(x_k,u_k)+J_{k+1}^\circ(x_{k+1})\Bigr].
\]

Nel caso continuo, facendo tendere a zero il passo temporale,
questa relazione diventa la Bellman continua:
\[
\min_u
\left[
\nabla_x\psi(x)^T F(x,u)+\ell(x,u)
\right]=0.
\]

Quindi la programmazione dinamica continua è l'analogo infinitesimo
della ricorsione di Bellman discreta.

\section{Sintesi della lezione}

In questa lezione abbiamo visto che:

\begin{itemize}
    \item nel problema LQ discreto a orizzonte finito la soluzione si ottiene
    tramite la ricorsione all'indietro di Riccati;
    \item se la ricorsione raggiunge un punto fisso \(T\), si ottiene
    l'equazione algebrica di Riccati:
    \[
    T
    =
    V
    +
    A^T
    \left[
    T
    -
    TB(P+B^TTB)^{-1}B^TT
    \right]A;
    \]
    \item l'orizzonte finito produce guadagni tempo-varianti \(H_k\),
    mentre l'orizzonte infinito porta a un guadagno costante;
    \item l'MPC nasce dall'idea di risolvere ripetutamente problemi LQ su orizzonte finito,
    applicando ogni volta solo una parte del controllo ottimo;
    \item questa strategia è particolarmente utile per impianti non lineari o
    tempo-varianti, che vengono linearizzati localmente;
    \item nel caso continuo e non lineare, introducendo il cost-to-go
    \(\psi(x)\), si ottiene l'equazione di Bellman:
    \[
    \min_u\left[\nabla_x\psi(x)^T F(x,u)+\ell(x,u)\right]=0.
    \]
\end{itemize}

Questa equazione è il punto di partenza teorico della programmazione dinamica
nel continuo.
\chapter{Lezione 20 --- Controllo ottimo LQ continuo, equazioni di Riccati e valori singolari}

\section{Dal cost-to-go all'equazione di Bellman nel continuo}

Consideriamo un sistema dinamico non lineare in tempo continuo
\[
\dot x = f(x,u),
\]
e un indice di ottimalità su orizzonte infinito
\[
J=\int_{t_0}^{+\infty}\ell(x(\tau),u(\tau))\,d\tau.
\]

Definiamo il \emph{cost-to-go} ottimo
\[
\psi(x_0)=\min_{u(\cdot)} J
\qquad\text{con } x(t_0)=x_0.
\]

Ragionando come nella programmazione dinamica, il costo ottimo a partire da \(x(t)\) deve soddisfare
\[
\psi(x(t))
=
\min_{u(\cdot)}
\left[
\int_t^{t+h}\ell(x(\tau),u(\tau))\,d\tau
+
\psi(x(t+h))
\right].
\]

Portando \(\psi(x(t))\) a destra, dividendo per \(h\) e facendo tendere \(h\to 0^+\), si ottiene
\[
\min_u\left[\dot\psi(x(t))+\ell(x(t),u(t))\right]=0.
\]

Usando la regola della catena,
\[
\dot\psi(x(t))
=
\nabla_x \psi(x(t))^T \dot x(t)
=
\nabla_x \psi(x(t))^T f(x(t),u(t)),
\]
segue l'\textbf{equazione di Bellman}:
\[
\boxed{
\min_u\left[\nabla_x\psi(x)^T f(x,u)+\ell(x,u)\right]=0.
}
\]

\section{Caso LQ continuo su orizzonte infinito}

Consideriamo ora il caso lineare-quadratico:
\[
\dot x = Ax+Bu,
\]
con funzionale di costo
\[
J=\frac12\int_0^{+\infty}\left(x^TQx+u^TRu\right)\,d\tau,
\]
dove
\[
Q=Q^T\succeq 0,
\qquad
R=R^T\succ 0.
\]

Cerchiamo una soluzione del cost-to-go nella forma quadratica
\[
\psi(x)=\frac12 x^T P x,
\qquad
P=P^T.
\]

Poiché \(P\) è simmetrica,
\[
\nabla_x\psi(x)=Px.
\]

Sostituendo nell'equazione di Bellman si ottiene
\[
\min_u
\left[
x^TP(Ax+Bu)+\frac12 x^TQx+\frac12 u^TRu
\right]=0.
\]

\subsection{Calcolo del controllo ottimo}

Per trovare il minimo rispetto a \(u\), imponiamo la condizione del primo ordine:
\[
\nabla_u\!\left(
x^TPBu+\frac12 u^TRu
\right)=0.
\]

Essendo
\[
\nabla_u(x^TPBu)=B^TPx,
\qquad
\nabla_u\!\left(\frac12 u^TRu\right)=Ru,
\]
segue
\[
B^TPx+Ru=0,
\]
cioè
\[
\boxed{
u^\star=-R^{-1}B^TPx.
}
\]

Quindi il controllo ottimo è una retroazione lineare dello stato, con guadagno
\[
K=R^{-1}B^TP.
\]

\section{Equazione algebrica di Riccati nel continuo}

Sostituiamo \(u^\star\) nell'equazione di Bellman:
\[
x^TPAx
-\,
x^TPBR^{-1}B^TPx
+
\frac12 x^TQx
+
\frac12 x^TPBR^{-1}B^TPx
=0.
\]

Riordinando:
\[
x^TPAx
+\frac12 x^TQx
-\frac12 x^TPBR^{-1}B^TPx
=0.
\]

Ora osserviamo che lo scalare \(x^TPAx\) può essere simmetrizzato:
\[
x^TPAx
=
\frac12 x^T(PA+A^TP)x.
\]

Otteniamo dunque
\[
\frac12 x^T
\left(
A^TP+PA-PBR^{-1}B^TP+Q
\right)x=0
\qquad \forall x.
\]

Perciò la matrice tra parentesi deve annullarsi:
\[
\boxed{
A^TP+PA-PBR^{-1}B^TP+Q=0.
}
\]

Questa è l'\textbf{equazione algebrica di Riccati a tempo continuo}.

\subsection{Sistema in ciclo chiuso}

Ponendo
\[
u^\star=-Kx,
\qquad
K=R^{-1}B^TP,
\]
la dinamica in ciclo chiuso diventa
\[
\dot x=(A-BK)x.
\]

Quindi la matrice di stato del sistema controllato ottimamente è
\[
A_{cc}=A-BR^{-1}B^TP.
\]

\section{Stabilità e legame con Lyapunov}

Partiamo dall'equazione di Riccati
\[
A^TP+PA-PBR^{-1}B^TP+Q=0.
\]

Poiché
\[
K=R^{-1}B^TP,
\]
si ha
\[
A_{cc}=A-BK=A-BR^{-1}B^TP.
\]

Sostituendo, l'equazione di Riccati può essere riscritta come
\[
(A-BK)^TP + P(A-BK) + \widetilde Q = 0,
\]
dove
\[
\widetilde Q = Q + PBR^{-1}B^TP \succeq 0.
\]

Cioè
\[
\boxed{
A_{cc}^TP + PA_{cc} = -\widetilde Q.
}
\]

Questa è una \textbf{equazione di Lyapunov} per il sistema in ciclo chiuso.

Se \(P\succ 0\) e \(\widetilde Q \succ 0\), allora la funzione
\[
V(x)=\frac12 x^TPx
\]
è una funzione di Lyapunov e il sistema in ciclo chiuso è asintoticamente stabile.

\paragraph{Interpretazione}
La soluzione \(P\) dell'equazione di Riccati non serve solo a scrivere il controllo ottimo:
essa codifica anche una funzione energia/costo che certifica la stabilità del sistema controllato.

\section{Ipotesi standard: stabilizzabilità e rilevabilità}

Negli appunti compare la seguente idea fondamentale.

\begin{itemize}
    \item \textbf{Stabilizzabile}: gli autovalori non controllabili di \(A\) hanno parte reale strettamente negativa.
    \item \textbf{Rilevabile} (detectable): gli autovalori non osservabili hanno parte reale strettamente negativa.
\end{itemize}

Per il problema LQ continuo, sotto le ipotesi standard
\[
(A,B)\ \text{stabilizzabile},
\qquad
(A,\sqrt Q)\ \text{rilevabile},
\]
l'equazione algebrica di Riccati ammette una \emph{soluzione stabilizzante} \(P=P^T\succeq 0\), e il feedback
\[
u^\star=-R^{-1}B^TPx
\]
rende stabile il sistema in ciclo chiuso.

\medskip

\noindent
\textbf{Nota.}
Negli appunti la conclusione è espressa in forma forte come esistenza di una soluzione \(P>0\).
Per essere più precisi, la formulazione standard parla di soluzione simmetrica positiva semidefinita stabilizzante; nei casi ``ben posti'' trattati a lezione essa è spesso positiva definita.

\section{Caso LQ continuo su orizzonte finito}

Se invece considero un orizzonte finito \([0,t_f]\), il funzionale di costo diventa
\[
J=
\int_0^{t_f}
\left(
\frac12 x^TQx+\frac12 u^TRu
\right)d\tau
+
\frac12 x^T(t_f)Sx(t_f),
\]
con
\[
S=S^T\succ 0.
\]

In questo caso il cost-to-go dipende anche dal tempo:
\[
\psi(x,t)=\frac12 x^T P(t)x.
\]

La legge di controllo ottimo conserva la stessa forma
\[
\boxed{
u^\star(t)=-R^{-1}B^TP(t)x(t),
\qquad 0<t<t_f,
}
\]
ma il guadagno non è più costante: dipende dal tempo tramite \(P(t)\).

La matrice \(P(t)\) soddisfa l'\textbf{equazione differenziale di Riccati}:
\[
\boxed{
-\dot P(t)=A^TP(t)+P(t)A-P(t)BR^{-1}B^TP(t)+Q,
}
\]
con condizione terminale
\[
\boxed{
P(t_f)=S.
}
\]

\paragraph{Osservazione sul segno}
L'equazione è scritta con il segno ``\(-\dot P\)'' perché la si integra all'indietro, partendo dal tempo finale \(t_f\) e tornando verso \(0\).  
Negli appunti questa dipendenza all'indietro è il punto concettuale importante.

\section{Richiamo di algebra lineare: decomposizione ai valori singolari}

Nella seconda parte della lezione compare un richiamo sui \textbf{valori singolari}, utile per ragionare sul condizionamento numerico delle matrici.

Sia
\[
M\in\mathbb R^{m\times k},
\qquad m\ge k.
\]
Allora esistono matrici ortogonali
\[
U\in\mathbb R^{m\times m},
\qquad
V\in\mathbb R^{k\times k},
\]
tali che
\[
U^TU=I,
\qquad
V^TV=I,
\]
e una matrice \(\Sigma\in\mathbb R^{m\times k}\) della forma
\[
\Sigma=
\begin{bmatrix}
\Sigma'\\
0
\end{bmatrix},
\qquad
\Sigma'=
\operatorname{diag}(\sigma_1,\sigma_2,\dots,\sigma_k),
\]
con
\[
\sigma_1\ge \sigma_2\ge \cdots \ge \sigma_k\ge 0,
\]
tale che
\[
\boxed{
M=U\Sigma V^T.
}
\]

Questa è la \textbf{SVD} (Singular Value Decomposition).

\subsection{Interpretazione dei valori singolari}

I numeri \(\sigma_i\) sono i \textbf{valori singolari} di \(M\).
Possono essere visti come una generalizzazione degli autovalori al caso di matrici non necessariamente quadrate.

Essi misurano anche quanto la matrice sia vicina alla singolarità.

\begin{itemize}
    \item Se \(M\) è quadrata e invertibile, allora il più piccolo valore singolare è positivo:
    \[
    \sigma_k>0.
    \]
    \item Se \(\sigma_k=0\), la matrice perde rango.
\end{itemize}

\subsection{Numero di condizionamento}

La distanza dalla non invertibilità si misura tramite il rapporto
\[
\kappa(M)=\frac{\sigma_1}{\sigma_k},
\]
detto \textbf{numero di condizionamento}.

\begin{itemize}
    \item Se
    \[
    \frac{\sigma_1}{\sigma_k}\approx 1,
    \]
    la matrice è ben condizionata e lontana dalla singolarità.
    \item Se
    \[
    \frac{\sigma_1}{\sigma_k}\gg 1,
    \]
    la matrice è male condizionata, cioè quasi non invertibile.
\end{itemize}

\paragraph{Perché compare qui?}
Nel controllo ottimo e, più in generale, nei problemi di calcolo numerico del controllo, si invertono spesso matrici come
\[
R,\qquad P+B^TT_{k+1}B,\qquad \text{ecc.}
\]
Perciò il condizionamento numerico è importante: una matrice formalmente invertibile può essere problematica dal punto di vista computazionale se è molto vicina alla singolarità.

\section{Sintesi finale della lezione}

In questa lezione abbiamo visto che:

\begin{enumerate}
    \item dal cost-to-go nel continuo si ottiene l'equazione di Bellman
    \[
    \min_u\left[\nabla_x\psi(x)^Tf(x,u)+\ell(x,u)\right]=0;
    \]
    \item nel caso LQ continuo, cercando \(\psi(x)=\tfrac12 x^TPx\), si ricava il controllo ottimo
    \[
    u^\star=-R^{-1}B^TPx;
    \]
    \item la matrice \(P\) deve soddisfare l'equazione algebrica di Riccati
    \[
    A^TP+PA-PBR^{-1}B^TP+Q=0;
    \]
    \item tale equazione si collega a una Lyapunov del sistema in ciclo chiuso
    \[
    A_{cc}^TP+PA_{cc}=-\widetilde Q;
    \]
    \item su orizzonte finito compare invece l'equazione differenziale di Riccati
    \[
    -\dot P=A^TP+PA-PBR^{-1}B^TP+Q,
    \qquad
    P(t_f)=S;
    \]
    \item i valori singolari permettono di capire se una matrice è ben o mal condizionata, tramite
    \[
    \kappa(M)=\frac{\sigma_1}{\sigma_k}.
    \]
\end{enumerate}
\chapter{Lezione 21 --- Bande di Gershgorin, criterio di Nyquist per sistemi MIMO, SVD e minimi quadrati}

\section{Sistemi MIMO e dominanza diagonale}

Consideriamo un sistema dinamico MIMO descritto in frequenza da
\[
X(s)=G(s)U(s),
\qquad
X(s)\in\mathbb{C}^m,\quad U(s)\in\mathbb{C}^m,
\]
dove
\[
G(s)\in\mathbb{C}^{m\times m}.
\]

Se il sistema fosse completamente disaccoppiato, cioè con matrice di trasferimento diagonale,
\[
G(s)=
\begin{bmatrix}
g_{11}(s) & 0 & \cdots & 0\\
0 & g_{22}(s) & \cdots & 0\\
\vdots & \vdots & \ddots & \vdots\\
0 & 0 & \cdots & g_{mm}(s)
\end{bmatrix},
\]
potremmo studiare ogni canale separatamente, analizzando ciascun elemento diagonale.

Se invece il sistema non è diagonale ma \emph{a dominanza diagonale}, allora gli elementi fuori diagonale possono essere interpretati come perturbazioni rispetto ai termini diagonali dominanti.

\medskip

\noindent
\textbf{Dominanza diagonale per righe.}
Una matrice \(A\in\mathbb{R}^{m\times m}\) si dice diagonalmente dominante per righe se, per ogni \(i=1,\dots,m\),
\[
|a_{ii}| \gg \sum_{\substack{j=1\\j\neq i}}^{m}|a_{ij}|.
\]

\noindent
\textbf{Dominanza diagonale per colonne.}
Analogamente, si può richiedere che, per ogni \(i=1,\dots,m\),
\[
|a_{ii}| \gg \sum_{\substack{j=1\\j\neq i}}^{m}|a_{ji}|.
\]

Nel seguito usiamo la convenzione degli appunti, cioè quella associata alle \emph{colonne}.

\section{Teorema di Gershgorin per matrici di trasferimento}

Per ogni frequenza \(\omega\), la matrice
\[
G(j\omega)
\]
è una matrice complessa \(m\times m\). Possiamo quindi applicare il teorema di Gershgorin.

\begin{theorem}[Gershgorin]
Sia \(G(j\omega)\in\mathbb{C}^{m\times m}\). Allora tutti gli autovalori di \(G(j\omega)\) appartengono all'unione dei cerchi del piano complesso centrati nei termini diagonali \(g_{ii}(j\omega)\) e di raggio
\[
r_i(j\omega)=\sum_{\substack{j=1\\j\neq i}}^{m}|g_{ji}(j\omega)|,
\qquad i=1,\dots,m.
\]
\end{theorem}

I cerchi di raggio \(r_i(j\omega)\) sono detti \textbf{cerchi di Gershgorin}.  
Al variare di \(\omega\), il centro \(g_{ii}(j\omega)\) si muove lungo il diagramma polare del termine diagonale, mentre il raggio cambia con la frequenza.

L'inviluppo di tali cerchi al variare di \(\omega\) genera le cosiddette \textbf{bande di Gershgorin}.

\section{Criterio di Nyquist formulato sulle bande di Gershgorin}

Per un sistema MIMO diagonalmente dominante, si può usare il criterio di Nyquist non direttamente sugli autovalori della matrice, ma sulle bande di Gershgorin.

L'idea è la seguente:

\begin{itemize}
    \item per ogni \(\omega\), gli autovalori di \(G(j\omega)\) stanno dentro i cerchi di Gershgorin;
    \item quindi, al variare della frequenza, le traiettorie degli autovalori sono contenute dentro le bande di Gershgorin;
    \item se tali bande \emph{non avvolgono} il punto critico \(-1\), allora anche il sistema accoppiato risulta stabile in ciclo chiuso.
\end{itemize}

Questa è una \textbf{condizione sufficiente} di stabilità del sistema in ciclo chiuso.

\begin{figure}[h!]
\centering
\begin{tikzpicture}[scale=1.05,>=Latex]
    % axes
    \draw[->] (-3.8,0)--(3.8,0) node[right] {$\Re$};
    \draw[->] (0,-2.8)--(0,2.8) node[above] {$\Im$};

    % point -1
    \fill (-2,0) circle (1.4pt);
    \node[below] at (-2,0) {$-1$};

    % Nyquist-like curve
    \draw[very thick,green!50!black]
        (2.0,0.1)
        .. controls (1.0,-0.2) and (0.6,-1.7) ..
        (-0.8,-1.9)
        .. controls (-1.6,-1.8) and (-1.8,-0.5) ..
        (-1.2,-0.2)
        .. controls (-0.7,0.0) and (0.4,0.15) ..
        (1.8,0.1);

    % representative centers
    \coordinate (c1) at (1.3,0.05);
    \coordinate (c2) at (0.4,-0.7);
    \coordinate (c3) at (-0.5,-1.5);
    \coordinate (c4) at (-1.1,-1.2);
    \coordinate (c5) at (-1.25,-0.45);
    \coordinate (c6) at (-0.4,-0.15);

    % Gershgorin circles
    \draw[red!75!black] (c1) circle (0.32);
    \draw[red!75!black] (c2) circle (0.42);
    \draw[red!75!black] (c3) circle (0.35);
    \draw[red!75!black] (c4) circle (0.28);
    \draw[red!75!black] (c5) circle (0.22);
    \draw[red!75!black] (c6) circle (0.18);

    \node[green!50!black,right] at (1.7,-0.5) {$g_{ii}(j\omega)$};
    \node[red!75!black,left] at (-2.9,-1.7) {cerchi di Gershgorin};
\end{tikzpicture}
\caption{Idea qualitativa: il centro di ciascun cerchio si muove sul diagramma polare di \(g_{ii}(j\omega)\), mentre il raggio è dato dalla somma dei moduli dei termini fuori diagonale associati.}
\end{figure}

\section{Esempio qualitativo}

Negli appunti compare un esempio in cui si prende un termine diagonale del tipo
\[
g_{11}(s)=\frac{10}{1+\frac{\sqrt{2}}{4}s+\frac{1}{16}s^2}
\]
e un termine fuori diagonale del tipo
\[
g_{21}(s)=\frac{5}{s+2}.
\]

In tal caso il centro dei cerchi di Gershgorin è il diagramma polare di \(g_{11}(j\omega)\), mentre il raggio vale
\[
r_1(j\omega)=|g_{21}(j\omega)|=\left|\frac{5}{2+j\omega}\right|
=\frac{5}{\sqrt{4+\omega^2}}.
\]

Quindi, ad esempio,
\[
r_1(0)=\frac{5}{2}=2.5,
\qquad
r_1(10)=\frac{5}{\sqrt{104}}\approx 0.49.
\]

Per \(\omega=0\), se \(g_{11}(0)=10\), il cerchio iniziale taglia l'asse reale nel punto
\[
10-\frac{5}{2}=7.5.
\]

Questo spiega il disegno negli appunti: a bassa frequenza i cerchi sono più grandi, poi diventano progressivamente più piccoli all'aumentare della frequenza.

\section{SVD: decomposizione ai valori singolari}

Sia ora
\[
A\in\mathbb{R}^{n\times m},
\qquad n\ge m.
\]
Allora esistono matrici ortogonali
\[
U\in\mathbb{R}^{n\times n},
\qquad
V\in\mathbb{R}^{m\times m},
\]
e una matrice diagonale ``rettangolare'' \(\Sigma\in\mathbb{R}^{n\times m}\) tali che
\[
A=U\Sigma V^T,
\]
con
\[
\Sigma=
\begin{bmatrix}
\widetilde\Sigma\\
0
\end{bmatrix},
\qquad
\widetilde\Sigma=
\operatorname{diag}(\sigma_1,\dots,\sigma_m),
\qquad
\sigma_1\ge \sigma_2\ge \cdots \ge \sigma_m\ge 0.
\]

I numeri \(\sigma_1,\dots,\sigma_m\) sono i \textbf{valori singolari} di \(A\).

Se
\[
\sigma_m>0,
\]
allora \(A\) ha rango pieno.

\subsection{Interpretazione dei valori singolari}

I valori singolari forniscono una misura di quanto una matrice sia vicina alla singolarità:

\begin{itemize}
    \item se tutti i valori singolari sono ``lontani da zero'', la matrice è ben condizionata;
    \item se almeno uno è molto piccolo, la matrice è vicina alla non invertibilità.
\end{itemize}

In particolare, il rapporto
\[
\frac{\sigma_1}{\sigma_m}
\]
misura il condizionamento della matrice.

\section{Problema dei minimi quadrati}

Consideriamo il problema
\[
Ax \approx y,
\qquad
A\in\mathbb{R}^{n\times m},
\quad
x\in\mathbb{R}^m,
\quad
y\in\mathbb{R}^n,
\]
con \(n\ge m\).

In generale il sistema può essere sovradeterminato e non avere soluzione esatta. Si cerca allora \(x\) che minimizza lo scostamento tra \(Ax\) e \(y\), cioè
\[
\min_x \frac12\|Ax-y\|^2.
\]

Nel caso di errore gaussiano, questa scelta coincide con la stima ottima nel senso dei minimi quadrati.

\begin{figure}[h!]
\centering
\begin{tikzpicture}[scale=0.95,>=Latex]
    \draw[->] (0,0)--(6.5,0) node[right] {$x$};
    \draw[->] (0,0)--(0,4.8) node[above] {$y$};

    % fitted line
    \draw[thick,teal!70!black] (0.5,0.9)--(6.0,4.2);
    \node[teal!70!black,right] at (5.2,4.0) {$y^\star=Kx$};

    % points
    \foreach \p/\q in {
        0.8/1.4,
        1.3/1.1,
        2.0/1.8,
        2.7/2.4,
        3.1/2.2,
        3.8/3.1,
        4.3/2.9,
        5.0/3.8,
        5.6/3.6
    }{
        \fill (\p,\q) circle (1.5pt);
    }

    % residuals
    \draw[densely dashed] (0.8,1.4)--($(0.8,0.9)!0.12!(6.0,4.2)$);
    \draw[densely dashed] (2.7,2.4)--($(2.7,0.9)!0.42!(6.0,4.2)$);
    \draw[densely dashed] (4.3,2.9)--($(4.3,0.9)!0.70!(6.0,4.2)$);

    \node[below right] at (3.0,1.45) {residui};
\end{tikzpicture}
\caption{Interpretazione geometrica del problema dei minimi quadrati: si cerca la retta/modello che minimizza gli scarti.}
\end{figure}

\subsection{Equazioni normali}

Scriviamo il funzionale:
\[
J(x)=\frac12(Ax-y)^T(Ax-y)
=\frac12\left(x^TA^TAx+y^Ty-2x^TA^Ty\right).
\]

La condizione di ottimo del primo ordine è
\[
\nabla_x J(x)=0
\quad\Longrightarrow\quad
A^TAx-A^Ty=0.
\]

Quindi
\[
\boxed{
A^TAx=A^Ty.
}
\]

Se \(A^TA\) è invertibile, la soluzione dei minimi quadrati è
\[
\boxed{
x=(A^TA)^{-1}A^Ty.
}
\]

La matrice
\[
\boxed{
A^+=(A^TA)^{-1}A^T
}
\]
è la \textbf{pseudo-inversa} di \(A\) (nel caso di rango pieno sulle colonne).

\section{Pseudo-inversa tramite SVD}

Se
\[
A=U\Sigma V^T,
\]
allora
\[
A^TA = V\Sigma^T\Sigma V^T
      = V\widetilde\Sigma^2V^T.
\]

Perciò gli autovalori di \(A^TA\) sono
\[
\sigma_1^2,\dots,\sigma_m^2,
\]
cioè \textbf{i quadrati dei valori singolari} di \(A\).

La pseudo-inversa si può allora scrivere come
\[
\boxed{
A^+=V\widetilde\Sigma^{-1}U^T.
}
\]

Quindi la soluzione dei minimi quadrati diventa
\[
\boxed{
x=A^+y=V\widetilde\Sigma^{-1}U^Ty.
}
\]

\subsection{Interpretazione in coordinate ortonormali}

Se definiamo
\[
\theta = U^Ty,
\qquad
z = V^Tx,
\]
allora il problema si diagonalizza e diventa, componente per componente,
\[
z_i = \frac{\theta_i}{\sigma_i},
\qquad i=1,\dots,m.
\]

Questa lettura è molto utile:
\begin{itemize}
    \item \(U^Ty\) sono le coordinate di \(y\) nella base delle colonne di \(U\);
    \item \(V^Tx\) sono le coordinate di \(x\) nella base delle colonne di \(V\);
    \item il passaggio da \(\theta_i\) a \(z_i\) consiste nel dividere per il valore singolare \(\sigma_i\).
\end{itemize}

Se un \(\sigma_i\) è molto piccolo, il corrispondente coefficiente viene fortemente amplificato: questo spiega la sensibilità numerica del problema.

\section{Condizionamento numerico}

Supponiamo ora di avere una matrice quadrata invertibile
\[
Fx=y.
\]

Se il dato \(y\) è perturbato,
\[
y=y^\star+\delta y,
\]
allora anche la soluzione sarà perturbata:
\[
x=x^\star+\delta x,
\qquad
\delta x = F^{-1}\delta y.
\]

L'errore sulla soluzione dipende quindi da \(F^{-1}\).  
Una misura della sensibilità del problema è il \textbf{numero di condizionamento}
\[
\kappa(F)=\|F\|\,\|F^{-1}\|.
\]

Nel caso della norma \(2\),
\[
\boxed{
\kappa_2(F)=\frac{\sigma_{\max}(F)}{\sigma_{\min}(F)}.
}
\]

Se \(\kappa_2(F)\approx 1\), il problema è ben condizionato; se invece \(\kappa_2(F)\gg 1\), piccoli errori nei dati possono produrre grandi errori nella soluzione.

Nel caso della matrice \(A^TA\),
\[
\kappa_2(A^TA)=\frac{\sigma_1^2}{\sigma_m^2}=\kappa_2(A)^2.
\]

Questa è la ragione per cui le equazioni normali peggiorano il condizionamento.

\section{Sintesi della lezione}

In questa lezione abbiamo visto due blocchi di argomenti.

\subsection*{1. Stabilità di sistemi MIMO mediante Gershgorin e Nyquist}
\begin{itemize}
    \item Se un sistema MIMO è diagonalmente dominante, si può approssimare il comportamento considerando la diagonale come parte principale e i termini fuori diagonale come perturbazioni.
    \item Per ogni frequenza \(\omega\), gli autovalori di \(G(j\omega)\) stanno nei cerchi di Gershgorin.
    \item L'inviluppo dei cerchi al variare di \(\omega\) costituisce le bande di Gershgorin.
    \item Se tali bande non circondano il punto critico \(-1\), allora si ha una condizione sufficiente di stabilità in ciclo chiuso.
\end{itemize}

\subsection*{2. SVD, minimi quadrati e pseudo-inversa}
\begin{itemize}
    \item La SVD scrive una matrice come
    \[
    A=U\Sigma V^T.
    \]
    \item I valori singolari misurano quanto la matrice sia vicina alla singolarità.
    \item Il problema dei minimi quadrati
    \[
    \min_x \frac12\|Ax-y\|^2
    \]
    porta alle equazioni normali
    \[
    A^TAx=A^Ty.
    \]
    \item La soluzione è
    \[
    x=A^+y,
    \qquad
    A^+=(A^TA)^{-1}A^T = V\widetilde\Sigma^{-1}U^T.
    \]
    \item Il condizionamento dipende dal rapporto tra valore singolare massimo e minimo.
\end{itemize}
\chapter{Lezione 22 --- Controllo decentralizzato e disaccoppiamento nei sistemi MIMO}

\section{Richiamo: margine di fase e margine di guadagno}

Nel diagramma di Nyquist, la stabilità del sistema in ciclo chiuso si studia rispetto al punto critico
\[
-1.
\]
Le due quantità classiche che misurano la ``distanza'' dal punto critico sono:

\begin{itemize}
    \item \textbf{margine di fase} \(M_\varphi\),
    \item \textbf{margine di guadagno} \(M_G\).
\end{itemize}

\begin{figure}[h!]
\centering
\begin{tikzpicture}[scale=1.1,>=Latex]
    % assi
    \draw[->] (-2.3,0)--(2.8,0) node[right] {};
    \draw[->] (0,-2.4)--(0,2.3) node[above] {};

    % punto critico
    \fill (-1,0) circle (1.3pt);
    \node[below] at (-1,0) {$-1$};

    % cerchi qualitativi
    \draw[thick] (0,0) circle (1.35);
    \draw[thick] (0,0) circle (0.75);

    % curve Nyquist qualitative
    \draw[very thick,red]
        (2.4,0) .. controls (1.5,-1.8) and (0.3,-2.1) .. (-1.1,-1.2)
        .. controls (-1.8,-0.8) and (-1.2,-0.1) .. (-0.35,0.15)
        .. controls (0.4,0.25) and (0.8,-0.05) .. (0.55,-0.45)
        .. controls (0.3,-0.8) and (-0.2,-0.55) .. (-0.05,-0.1)
        .. controls (0.1,0.15) and (0.35,0.12) .. (0.48,0.0);

    \draw[very thick]
        (2.1,0) .. controls (1.3,-1.5) and (0.2,-1.8) .. (-0.95,-1.05)
        .. controls (-1.4,-0.65) and (-1.0,-0.1) .. (-0.3,0.05)
        .. controls (0.25,0.15) and (0.55,-0.05) .. (0.38,-0.33)
        .. controls (0.2,-0.58) and (-0.08,-0.38) .. (0.0,-0.08)
        .. controls (0.08,0.08) and (0.22,0.05) .. (0.3,0.0);

    % frecce direzione
    \draw[->,thick] (-1.15,-1.15)--(-1.32,-1.0);
    \draw[->,thick] (-0.55,-0.55)--(-0.7,-0.4);
    \draw[->,thick,red] (-1.3,-1.25)--(-1.5,-1.1);
    \draw[->,thick,red] (-0.68,-0.72)--(-0.82,-0.56);

    % margine di guadagno
    \draw[dashed] (-1,0)--(-1,1.5);
    \draw[dashed] (-1.45,0)--(-1.45,1.5);
    \draw[<->,blue] (-1.45,1.38)--(-1,1.38);
    \node[blue,above] at (-1.225,1.42) {\small \(1/M_G\)};

    % margine di fase
    \draw[blue,thick] (-1,0)--(-1.0,-0.85);
    \draw[blue,thick] (-1,0)--(-0.42,-0.55);
    \draw[->,blue,thick] (-0.86,-0.35) arc[start angle=-100,end angle=-52,radius=0.52];
    \node[blue] at (-0.92,-0.72) {\small \(M_\varphi\)};
\end{tikzpicture}
\caption{Interpretazione qualitativa di margine di fase \(M_\varphi\) e margine di guadagno \(M_G\) sul diagramma di Nyquist.}
\end{figure}

\section{Sistema MIMO \(2\times 2\)}

Consideriamo un impianto MIMO \(2\times 2\) con matrice di trasferimento
\[
G(s)=
\begin{bmatrix}
G_{11}(s) & G_{12}(s)\\
G_{21}(s) & G_{22}(s)
\end{bmatrix}.
\]

I singoli elementi hanno il significato seguente:

\begin{itemize}
    \item \(G_{11}(s)\): funzione di trasferimento tra ingresso \(1\) e uscita \(1\),
    \item \(G_{22}(s)\): funzione di trasferimento tra ingresso \(2\) e uscita \(2\),
    \item \(G_{21}(s)\): funzione di trasferimento tra ingresso \(1\) e uscita \(2\),
    \item \(G_{12}(s)\): funzione di trasferimento tra ingresso \(2\) e uscita \(1\).
\end{itemize}

I termini fuori diagonale \(G_{12}\) e \(G_{21}\) rappresentano le \textbf{interferenze tra canali}.

\section{Controllo decentralizzato}

Nel controllo decentralizzato si usano due regolatori separati, uno per ciascun canale, ignorando nella sintesi iniziale i termini di accoppiamento.

\[
R(s)=
\begin{bmatrix}
R_1(s) & 0\\
0 & R_2(s)
\end{bmatrix}.
\]

La struttura a blocchi è la seguente.

\begin{figure}[h!]
\centering
\begin{tikzpicture}[scale=0.95,>=Latex]
    % r1 path
    \node at (-0.2,2.8) {$r_1$};
    \draw (-0.1,2.8) circle (0.08);
    \draw[->] (0,2.8)--(0.8,2.8);
    \node at (1.0,2.8) {$\sum$};
    \draw[->] (1.2,2.8)--(2.1,2.8);
    \draw (2.1,2.45) rectangle (3.2,3.15);
    \node at (2.65,2.8) {$R_1$};

    \draw[->] (3.2,2.8)--(4.2,2.8);
    \node at (4.4,2.8) {$\sum$};
    \draw[->] (4.6,2.8)--(5.5,2.8);
    \node[above] at (4.95,2.9) {$u_1$};
    \draw (5.5,2.45) rectangle (6.8,3.15);
    \node at (6.15,2.8) {$G_{11}(s)$};
    \draw[->] (6.8,2.8)--(8.0,2.8);
    \node[right] at (8.0,2.8) {$y_1$};

    % feedback y1
    \draw (7.6,2.8)--(7.6,4.1)--(0.8,4.1)--(0.8,2.92);
    \draw[->] (0.8,2.92)--(0.8,2.88);

    % r2 path
    \node at (-0.2,0.8) {$r_2$};
    \draw (-0.1,0.8) circle (0.08);
    \draw[->] (0,0.8)--(0.8,0.8);
    \node at (1.0,0.8) {$\sum$};
    \draw[->] (1.2,0.8)--(2.1,0.8);
    \draw (2.1,0.45) rectangle (3.2,1.15);
    \node at (2.65,0.8) {$R_2$};

    \draw[->] (3.2,0.8)--(4.2,0.8);
    \node at (4.4,0.8) {$\sum$};
    \draw[->] (4.6,0.8)--(5.5,0.8);
    \node[above] at (4.95,0.9) {$u_2$};
    \draw (5.5,0.45) rectangle (6.8,1.15);
    \node at (6.15,0.8) {$G_{22}(s)$};
    \draw[->] (6.8,0.8)--(8.0,0.8);
    \node[right] at (8.0,0.8) {$y_2$};

    % feedback y2
    \draw (7.6,0.8)--(7.6,-0.5)--(0.8,-0.5)--(0.8,0.68);
    \draw[->] (0.8,0.68)--(0.8,0.72);

    % interaction blocks
    \draw (4.3,1.45) rectangle (5.2,2.05);
    \node at (4.75,1.75) {$G_{12}$};
    \draw[->] (4.75,1.15)--(4.75,1.45);
    \draw[->] (4.75,2.05)--(4.75,2.6);

    \draw (3.5,1.45) rectangle (4.4,2.05);
    \node at (3.95,1.75) {$G_{21}$};
    \draw[->] (3.95,2.45)--(3.95,2.05);
    \draw[->] (3.95,1.45)--(3.95,1.0);

    \node at (8.6,2.1) {\small controllo decentralizzato};
\end{tikzpicture}
\caption{Struttura di controllo decentralizzato per un sistema \(2\times 2\): i canali principali sono regolati separatamente, ma restano le interferenze incrociate \(G_{12}\) e \(G_{21}\).}
\end{figure}

Il controllo decentralizzato, da solo, \textbf{non elimina} le interferenze tra canali.

\section{Inserimento del disaccoppiatore}

Per ridurre o annullare le interazioni, si introduce un sistema dinamico detto \textbf{disaccoppiatore}
\[
D(s),
\]
posto tra il regolatore e l'impianto.

Lo schema complessivo diventa

\[
R(s)\;\longrightarrow\; D(s)\;\longrightarrow\; G(s).
\]

Definiamo
\[
M(s)=G(s)D(s).
\]

L'obiettivo è scegliere \(D(s)\) in modo che \(M(s)\) risulti \textbf{diagonale}, o almeno quasi diagonale. Se questo accade, il sistema diventa disaccoppiato e il controllo decentralizzato diventa naturale.

\begin{figure}[h!]
\centering
\begin{tikzpicture}[scale=1,>=Latex]
    \draw (0,0) circle (0.08);
    \draw[->] (0.08,0)--(1.0,0);
    \node at (0.5,0.25) {$e$};
    \draw (1.0,-0.35) rectangle (2.0,0.35);
    \node at (1.5,0) {$R(s)$};
    \draw[->] (2.0,0)--(3.0,0);
    \draw (3.0,-0.35) rectangle (4.0,0.35);
    \node at (3.5,0) {$D(s)$};
    \draw[->] (4.0,0)--(5.0,0);
    \draw (5.0,-0.35) rectangle (6.0,0.35);
    \node at (5.5,0) {$G(s)$};
    \draw[->] (6.0,0)--(7.0,0);
    \node[right] at (7.0,0) {$y$};

    \draw (6.6,0)--(6.6,-1.0)--(0.5,-1.0)--(0.5,-0.08);
    \draw[->] (0.5,-0.08)--(0.5,-0.02);
    \node at (6.9,0.25) {};
\end{tikzpicture}
\caption{Struttura di controllo con disaccoppiatore dinamico \(D(s)\).}
\end{figure}

\section{Richiamo sul controllo classico}

Nel caso SISO classico, con
\[
e=y^\star-y,\qquad u=Re,\qquad y=Gu,
\]
si ha
\[
y=GR(y^\star-y)=GRy^\star-GRy.
\]
Da cui
\[
(I+GR)y=GR\,y^\star
\]
e quindi
\[
y=(I+GR)^{-1}GR\,y^\star.
\]

La matrice/funzione
\[
H=(I+GR)^{-1}GR
\]
è la funzione di trasferimento in ciclo chiuso.

Nel caso MIMO con disaccoppiatore, si desidera che il blocco equivalente
\[
M(s)=G(s)D(s)
\]
sia diagonale.

\section{Disaccoppiamento ideale nel caso \(2\times 2\)}

Supponiamo di volere
\[
M(s)=
\begin{bmatrix}
M_{11}(s) & 0\\
0 & M_{22}(s)
\end{bmatrix}.
\]
Se \(G(s)\) è invertibile, allora
\[
D(s)=G^{-1}(s)M(s).
\]

Per una matrice \(2\times 2\),
\[
G^{-1}(s)=\frac{1}{\Delta(s)}
\begin{bmatrix}
G_{22}(s) & -G_{12}(s)\\
-G_{21}(s) & G_{11}(s)
\end{bmatrix},
\]
dove
\[
\Delta(s)=G_{11}(s)G_{22}(s)-G_{12}(s)G_{21}(s).
\]

Quindi
\[
G^{-1}(s)M(s)
=\frac{1}{\Delta(s)}
\begin{bmatrix}
G_{22}M_{11} & -G_{12}M_{22}\\
-G_{21}M_{11} & G_{11}M_{22}
\end{bmatrix}.
\]

Una scelta naturale è mantenere sulle diagonali le dinamiche principali originarie:
\[
M_{11}(s)=G_{11}(s),\qquad M_{22}(s)=G_{22}(s).
\]
Con questa scelta si ottiene il \textbf{disaccoppiatore ideale}
\[
\boxed{
D(s)=
\begin{bmatrix}
\dfrac{G_{22}G_{11}}{G_{11}G_{22}-G_{12}G_{21}}
&
-\dfrac{G_{12}G_{22}}{G_{11}G_{22}-G_{12}G_{21}}
\\[1.2em]
-\dfrac{G_{21}G_{11}}{G_{11}G_{22}-G_{12}G_{21}}
&
\dfrac{G_{11}G_{22}}{G_{11}G_{22}-G_{12}G_{21}}
\end{bmatrix}.
}
\]

Ed effettivamente
\[
G(s)D(s)=
\begin{bmatrix}
G_{11}(s) & 0\\
0 & G_{22}(s)
\end{bmatrix}.
\]

\section{Interpretazione del disaccoppiamento ideale}

Formalmente il risultato sembra perfetto: il sistema risultante è diagonale. Tuttavia questo disaccoppiamento avviene tramite \textbf{cancellazioni dinamiche}.

Infatti
\[
M(s)=G(s)D(s)=G(s)G^{-1}(s)\,\operatorname{diag}(G(s))
=\operatorname{diag}(G(s)).
\]

Quindi stiamo cancellando, attraverso \(D(s)\), le parti non diagonali e parte delle dinamiche del sistema.

Dal punto di vista fisico questo è delicato, perché le cancellazioni possono eliminare modi interni che risultano poi non raggiungibili o non osservabili nello schema equivalente.

\section{Quando si possono cancellare poli e zeri?}

Nel progetto del regolatore si possono cancellare soltanto poli e zeri \textbf{a parte reale strettamente negativa}, cioè associati a \textbf{fase minima}.

Quindi, poiché il disaccoppiamento ideale lavora proprio per cancellazioni, è necessario richiedere che:

\begin{itemize}
    \item le dinamiche coinvolte siano a fase minima;
    \item non ci siano ritardi di tempo significativi;
    \item le cancellazioni siano fisicamente robuste.
\end{itemize}

In sintesi:

\[
\boxed{\text{il disaccoppiatore ideale richiede condizioni di fase minima}}
\]

ed è inoltre molto sensibile a:

\begin{itemize}
    \item incertezze parametriche,
    \item errori di modello,
    \item ritardi di tempo.
\end{itemize}

\section{Caso quasi diagonale}

Se i termini fuori diagonale sono piccoli, non è necessario disaccoppiare in modo esatto.

In particolare, se
\[
G_{11}(s)G_{22}(s)\gg G_{12}(s)G_{21}(s),
\]
allora nel determinante
\[
\Delta(s)=G_{11}(s)G_{22}(s)-G_{12}(s)G_{21}(s)
\]
il termine dominante è \(G_{11}G_{22}\), mentre il prodotto \(G_{12}G_{21}\) può essere trascurato.

In questo caso il sistema è \textbf{quasi diagonale}.

\section{Disaccoppiatore semplificato}

Trascurando il prodotto \(G_{12}G_{21}\) nel denominatore, dal disaccoppiatore ideale si ottiene il \textbf{disaccoppiatore semplificato}
\[
\boxed{
D(s)\approx
\begin{bmatrix}
1 & -\dfrac{G_{12}(s)}{G_{11}(s)}\\[0.8em]
-\dfrac{G_{21}(s)}{G_{22}(s)} & 1
\end{bmatrix}.
}
\]

Questo non disaccoppia perfettamente il sistema, ma riduce sensibilmente l'interazione cross-channel.

\subsection{Caso triangolare}

Se il sistema è triangolare, ad esempio
\[
G(s)=
\begin{bmatrix}
G_{11}(s) & 0\\
G_{21}(s) & G_{22}(s)
\end{bmatrix},
\]
si può scegliere
\[
D(s)=
\begin{bmatrix}
1 & 0\\[0.4em]
-\dfrac{G_{21}(s)}{G_{22}(s)} & 1
\end{bmatrix}.
\]

In questo caso la cancellazione coinvolge solo il termine \(G_{21}\) rispetto a \(G_{22}\), e la struttura è più favorevole del caso generale.

\section{Schema a blocchi con disaccoppiatore semplificato}

Il disaccoppiatore semplificato introduce due rami di compensazione incrociati che cercano di annullare, almeno in prima approssimazione, le interferenze.

\begin{figure}[h!]
\centering
\begin{tikzpicture}[scale=0.9,>=Latex]
    % upper channel
    \node at (-0.3,3.2) {$r_1$};
    \draw (-0.15,3.2) circle (0.08);
    \draw[->] (-0.07,3.2)--(0.7,3.2);
    \node at (0.9,3.2) {$\sum$};
    \draw[->] (1.1,3.2)--(2.0,3.2);
    \draw (2.0,2.85) rectangle (3.1,3.55);
    \node at (2.55,3.2) {$R_1(s)$};
    \draw[->] (3.1,3.2)--(4.2,3.2);

    \node at (4.4,3.2) {$\sum$};
    \draw[->] (4.6,3.2)--(5.5,3.2);
    \draw (5.5,2.85) rectangle (6.8,3.55);
    \node at (6.15,3.2) {$G_{11}(s)$};
    \draw[->] (6.8,3.2)--(8.0,3.2);
    \node[right] at (8.0,3.2) {$y_1$};

    % lower channel
    \node at (-0.3,0.8) {$r_2$};
    \draw (-0.15,0.8) circle (0.08);
    \draw[->] (-0.07,0.8)--(0.7,0.8);
    \node at (0.9,0.8) {$\sum$};
    \draw[->] (1.1,0.8)--(2.0,0.8);
    \draw (2.0,0.45) rectangle (3.1,1.15);
    \node at (2.55,0.8) {$R_2(s)$};
    \draw[->] (3.1,0.8)--(4.2,0.8);

    \node at (4.4,0.8) {$\sum$};
    \draw[->] (4.6,0.8)--(5.5,0.8);
    \draw (5.5,0.45) rectangle (6.8,1.15);
    \node at (6.15,0.8) {$G_{22}(s)$};
    \draw[->] (6.8,0.8)--(8.0,0.8);
    \node[right] at (8.0,0.8) {$y_2$};

    % feedbacks
    \draw (7.6,3.2)--(7.6,4.3)--(0.7,4.3)--(0.7,3.28);
    \draw[->] (0.7,3.28)--(0.7,3.24);

    \draw (7.6,0.8)--(7.6,-0.3)--(0.7,-0.3)--(0.7,0.72);
    \draw[->] (0.7,0.72)--(0.7,0.76);

    % decoupler blocks
    \draw (3.55,1.75) rectangle (4.8,2.35);
    \node at (4.175,2.05) {$-\dfrac{G_{12}}{G_{11}}$};
    \draw[->] (4.175,1.15)--(4.175,1.75);
    \draw[->] (4.175,2.35)--(4.175,3.0);

    \draw (3.0,1.75) rectangle (4.2,2.35);
    \node at (3.6,2.05) {$-\dfrac{G_{21}}{G_{22}}$};
    \draw[->] (3.6,2.85)--(3.6,2.35);
    \draw[->] (3.6,1.75)--(3.6,1.0);

    % cross couplings remaining
    \draw (5.0,1.75) rectangle (5.8,2.35);
    \node at (5.4,2.05) {$G_{21}$};
    \draw[->] (5.4,2.85)--(5.4,2.35);
    \draw[->] (5.4,1.75)--(5.4,1.0);

    \draw (5.9,1.75) rectangle (6.7,2.35);
    \node at (6.3,2.05) {$G_{12}$};
    \draw[->] (6.3,1.15)--(6.3,1.75);
    \draw[->] (6.3,2.35)--(6.3,3.0);

    % ellipse around decoupler
    \draw[green!60!black,thick] (3.95,2.0) ellipse (1.35 and 1.5);
    \node[green!60!black] at (2.9,2.95) {$D(s)$};
\end{tikzpicture}
\caption{Schema qualitativo con disaccoppiatore semplificato: l'interazione non viene eliminata del tutto, ma viene ridotta.}
\end{figure}

\section{Generalizzazione qualitativa al caso \(n\times n\)}

Nel caso \(n\times n\), l'idea del disaccoppiatore semplificato è costruire una matrice del tipo
\[
D(s)\approx
\begin{bmatrix}
1 & -\dfrac{G_{12}}{G_{11}} & \cdots & -\dfrac{G_{1n}}{G_{11}}\\[0.8em]
-\dfrac{G_{21}}{G_{22}} & 1 & \cdots & -\dfrac{G_{2n}}{G_{22}}\\
\vdots & \vdots & \ddots & \vdots\\
-\dfrac{G_{n1}}{G_{nn}} & -\dfrac{G_{n2}}{G_{nn}} & \cdots & 1
\end{bmatrix}.
\]

Questa espressione è utile come linea guida, ma non sempre può essere realizzata in pratica: bisogna verificare che non compaiano cancellazioni non ammissibili, ritardi, o termini instabili.

\section{Osservazioni finali}

\begin{itemize}
    \item Il \textbf{controllo decentralizzato} usa un regolatore per ciascun canale, ma lascia le interazioni.
    \item Il \textbf{disaccoppiatore} viene inserito per trasformare l'impianto in uno quasi diagonale o diagonale.
    \item Il \textbf{disaccoppiatore ideale} si ottiene imponendo
    \[
    M(s)=G(s)D(s)
    \]
    diagonale e usando
    \[
    D(s)=G^{-1}(s)M(s).
    \]
    \item Questa soluzione è però basata su \textbf{cancellazioni dinamiche}, quindi richiede condizioni di fase minima ed è molto sensibile a ritardi e incertezze.
    \item Se i canali sono già debolmente accoppiati, cioè se
    \[
    G_{11}G_{22}\gg G_{12}G_{21},
    \]
    conviene usare il \textbf{disaccoppiatore semplificato}, più robusto e più facile da realizzare.
\end{itemize}

\end{document}